% El emprendimiento y los intercambios comerciales — Espagnol 1ère A — Cours signé OnBuch+
% Programme officiel MINESEC. Compiler avec : tectonic EA14-emprendimiento-intercambios-comerciales.tex
\newcommand{\DOCMATIERE}{Espagnol}
\newcommand{\DOCNIVEAU}{1ère A}
\newcommand{\DOCMODULE}{Módulo 4 : Vie économique}
\newcommand{\DOCLECON}{EA14}
\newcommand{\DOCTITRE}{El emprendimiento y los intercambios comerciales}
% OnBuch+ — préambule « cours signés » ENRICHI (compilé avec tectonic / XeLaTeX)
% Système visuel « Pop » d'OnBuch+ : fond crème (jamais blanc), bordures encre
% épaisses, ombres « dures » décalées, titres Archivo Black en capitales.
% Version MATHÉMATIQUES : graphes (pgfplots/tikz), boîtes pédagogiques
% (définition, propriété, méthode, attention, le savais-tu, exercice résolu,
% exercice, TP, bilan), page de garde et signature OnBuch+ ancrée en bas.
%
% Macros attendues dans main.tex AVANT \input{preamble} :
%   \DOCMATIERE  \DOCNIVEAU  \DOCMODULE  \DOCLECON (numéro)  \DOCTITRE
\documentclass[11pt]{article}

\usepackage{fontspec}
\usepackage[spanish,french]{babel} % français = langue principale ; l'espagnol via \esp{..} / espagnol
\usepackage[a4paper,margin=2cm,top=3.4cm,bottom=2.8cm,headheight=1.4cm,footskip=1.3cm]{geometry}
\usepackage{amsmath,amssymb,amsfonts}
\usepackage{mathtools} % \xmapsto, \coloneqq... souvent utilisés par les modèles
\usepackage{cancel} % simplifications barrées (bilans de forces)
% \abs{z} (module, valeur absolue) et \norm{u} (norme), utilisés par les modèles
\providecommand{\abs}{}\let\abs\relax\DeclarePairedDelimiter{\abs}{\lvert}{\rvert}
\providecommand{\norm}{}\let\norm\relax\DeclarePairedDelimiter{\norm}{\lVert}{\rVert}
\usepackage[table]{xcolor} % [table] : fournit aussi \rowcolors (alternance de lignes)
\usepackage{pagecolor}
\usepackage{tikz}
\usetikzlibrary{arrows.meta,positioning,calc,shapes.geometric,shapes.misc,shapes.arrows,decorations.pathmorphing,decorations.pathreplacing,decorations.markings,patterns,shadows,mindmap,trees,fit,backgrounds,matrix,intersections,angles,quotes}
\usepackage{pgfplots}
\usepackage[version=4]{mhchem} % équations nucléaires \ce{^{235}_{92}U}
\usepackage[european,siunitx]{circuitikz} % schémas électriques
\pgfplotsset{compat=1.18}
\usepgfplotslibrary{fillbetween,groupplots} % aires entre courbes (calcul intégral)
\usepackage{enumitem}
\usepackage{fancyhdr}
% [most] charge toutes les bibliothèques tcolorbox (listings, minted, raster,
% documentation...) et gonfle inutilement la mémoire TeX sur un document long
% (~300 boîtes) : on ne charge que ce qui est réellement utilisé.
\usepackage[skins,breakable]{tcolorbox}
\usepackage{graphicx}
\usepackage{colortbl}
\usepackage{booktabs}
\usepackage{tabularx}
\usepackage{diagbox} % case d'en-tête diagonale des tableaux à double entrée
% Colonnes extensibles centrée / gauche / droite (C, L, R), souvent utilisées par les modèles
\newcolumntype{C}{>{\centering\arraybackslash}X}
\newcolumntype{L}{>{\raggedright\arraybackslash}X}
\newcolumntype{R}{>{\raggedleft\arraybackslash}X}
\usepackage{array}
\usepackage{multicol}
\usepackage{siunitx}
% parse-numbers=false : accepte aussi \SI{2,5 \times 10^{-3}}{...}
\sisetup{locale=FR,output-decimal-marker={,},group-minimum-digits=4,parse-numbers=false}
\DeclareSIUnit\ohm{\ensuremath{\symup{\Omega}}} % U+2126 absent de la police
\DeclareSIUnit\division{div} % réglages d'oscilloscope (ms/div, V/div)
\DeclareSIUnit\tr{tr} % tours (vitesse de rotation, ex. \SI{1500}{\tr\per\minute})
\DeclareSIUnit\molar{M} % concentration molaire (ex. \SI{3}{\molar}, \SI{1}{\micro\molar})
\DeclareSIUnit\an{an} % année (le modèle confond parfois avec \year, un compteur TeX) — ex. \SI{5}{\kilo\meter\per\an}
\DeclareSIUnit\FCFA{FCFA} % franc CFA (ex. \SI{500}{\FCFA\per\kilo\gram})
\DeclareSIUnit\degreeBrix{\textdegree Brix} % teneur en sucre d'un sirop/jus (réfractométrie)
\DeclareSIUnit\month{mois} % le modèle utilise parfois \month, un compteur TeX, comme unité de durée
\usepackage{qrcode}
% \degree : commande fréquemment utilisée par les modèles pour un angle ou une
% température en dehors de \SI{}{} (non fournie par les paquets chargés).
\providecommand{\degree}{^{\circ}}
% \qed (fin de démonstration), utilisé par les modèles sans amsthm
\providecommand{\qed}{\hfill$\square$}
% \barre : double barre des valeurs interdites dans un tableau de variations (tkz-tab)
\providecommand{\barre}{\ensuremath{\|}}
% environnement proof (amsthm non chargé) utilisé par les modèles
\ifdefined\proof\else\newenvironment{proof}[1][Démonstration]{\par\noindent\textit{#1.}\ }{\qed\par}\fi
\usepackage{lastpage}
\usepackage{hyperref}
\hypersetup{hidelinks}

% ============================================================
% Palette « Pop » OnBuch+ — mêmes hex que la classe P (plus_ui.dart)
% ============================================================
\definecolor{poporange}{HTML}{F59321}
\definecolor{popdark}{HTML}{E07A0C}
\definecolor{popcream}{HTML}{FFF6E8}
\definecolor{popink}{HTML}{1C1714}
\definecolor{poppurple}{HTML}{7A5AE0}
\definecolor{popgreen}{HTML}{2E9E6A}
\definecolor{poppink}{HTML}{E0457B}
\definecolor{popblue}{HTML}{2F7DE1}
\definecolor{popgold}{HTML}{C98A12}
\definecolor{popmuted}{HTML}{8A7F76}
\definecolor{popline}{HTML}{EADFD2}
\definecolor{popgray}{HTML}{8A7F76}
\colorlet{popgrey}{popgray}
% Alias tolérants (le modèle nomme parfois les couleurs différemment)
\colorlet{popwhite}{white}
\colorlet{popblack}{popink}
\colorlet{popred}{poppink}
\colorlet{popyellow}{popgold}
% Teintes claires (fonds de tableaux/encadrés) — le modèle nomme parfois
% les couleurs avec un suffixe L (ex. popblueL) sans les avoir définies.
\colorlet{poporangeL}{poporange!20}
\colorlet{popdarkL}{popdark!20}
\colorlet{poppurpleL}{poppurple!20}
\colorlet{popgreenL}{popgreen!20}
\colorlet{poppinkL}{poppink!20}
\colorlet{popblueL}{popblue!20}
\colorlet{popgoldL}{popgold!20}
\colorlet{popredL}{popred!20}
\colorlet{popgrayL}{popgray!20}
% Teintes claires pour fonds de tableaux / zones
\colorlet{popblueL}{popblue!10!white}
\colorlet{poporangeL}{poporange!14!white}
\colorlet{popgreenL}{popgreen!12!white}
\colorlet{poppurpleL}{poppurple!12!white}
\colorlet{poppinkL}{poppink!12!white}
\colorlet{popgoldL}{popgold!14!white}

\pagecolor{popcream}

% ============================================================
% Typographie
% ============================================================
\defaultfontfeatures{Ligatures=TeX}
\setmainfont{PlusJakartaSans-Medium.ttf}[Path=../../../fonts/,BoldFont=PlusJakartaSans-Bold.ttf]
% Maths en sans-serif (Fira Math) avec chiffres et lettres droites de Plus
% Jakarta, pour que les formules se fondent dans le texte.
\usepackage[math-style=ISO,bold-style=ISO]{unicode-math}
\setmathfont{FiraMath-Regular.otf}[Path=../../../fonts/]
\setmathfont{PlusJakartaSans-Medium.ttf}[Path=../../../fonts/,range={up/num,up/{Latin,latin},\mathup}]
\usepackage{icomma} % virgule décimale sans espace en mode math
% Caractères mathématiques Unicode absents des polices de texte (Plus Jakarta,
% Archivo Black) : rendus par leur équivalent mathématique, en texte comme en
% formule — sinon ils s'affichent en carré vide (ex. « Arithmétique dans ℤ »).
\usepackage{newunicodechar}
\newunicodechar{ℝ}{\ensuremath{\mathbb{R}}}
\newunicodechar{ℤ}{\ensuremath{\mathbb{Z}}}
\newunicodechar{ℂ}{\ensuremath{\mathbb{C}}}
\newunicodechar{ℕ}{\ensuremath{\mathbb{N}}}
\newunicodechar{ℚ}{\ensuremath{\mathbb{Q}}}
\newunicodechar{𝔻}{\ensuremath{\mathbb{D}}}
\newunicodechar{α}{\ensuremath{\alpha}}
\newunicodechar{β}{\ensuremath{\beta}}
\newunicodechar{γ}{\ensuremath{\gamma}}
\newunicodechar{δ}{\ensuremath{\delta}}
\newunicodechar{ε}{\ensuremath{\varepsilon}}
\newunicodechar{θ}{\ensuremath{\theta}}
\newunicodechar{λ}{\ensuremath{\lambda}}
\newunicodechar{μ}{\ensuremath{\mu}}
\newunicodechar{σ}{\ensuremath{\sigma}}
\newunicodechar{τ}{\ensuremath{\tau}}
\newunicodechar{φ}{\ensuremath{\varphi}}
\newunicodechar{ψ}{\ensuremath{\psi}}
\newunicodechar{ω}{\ensuremath{\omega}}
\newunicodechar{Σ}{\ensuremath{\Sigma}}
% majuscules produites par \MakeUppercase dans les titres (α-aminés -> Α)
\newunicodechar{Α}{\ensuremath{\alpha}}
\newunicodechar{Β}{\ensuremath{\beta}}
\newunicodechar{Γ}{\ensuremath{\Gamma}}
\newunicodechar{Δ}{\ensuremath{\Delta}}
\newunicodechar{Ε}{\ensuremath{\varepsilon}}
\newunicodechar{Θ}{\ensuremath{\Theta}}
\newunicodechar{Λ}{\ensuremath{\Lambda}}
\newunicodechar{Μ}{\ensuremath{\mu}}
\newunicodechar{Τ}{\ensuremath{\tau}}
\newunicodechar{Φ}{\ensuremath{\Phi}}
\newunicodechar{Ψ}{\ensuremath{\Psi}}
\newunicodechar{Ω}{\ensuremath{\Omega}}
\newunicodechar{↦}{\ensuremath{\mapsto}}
\newunicodechar{∈}{\ensuremath{\in}}
\newunicodechar{∉}{\ensuremath{\notin}}
\newunicodechar{∧}{\ensuremath{\wedge}}
\newunicodechar{⇔}{\ensuremath{\Leftrightarrow}}
\newunicodechar{⟺}{\ensuremath{\Longleftrightarrow}}
\newunicodechar{⇒}{\ensuremath{\Rightarrow}}
\newunicodechar{⟹}{\ensuremath{\Longrightarrow}}
\newunicodechar{∘}{\ensuremath{\circ}}
\newunicodechar{∪}{\ensuremath{\cup}}
\newunicodechar{∩}{\ensuremath{\cap}}
\newunicodechar{≡}{\ensuremath{\equiv}}
\newunicodechar{⊂}{\ensuremath{\subset}}
\newunicodechar{⊥}{\ensuremath{\perp}}
\newunicodechar{∃}{\ensuremath{\exists}}
\newunicodechar{∀}{\ensuremath{\forall}}
\newunicodechar{∛}{\ensuremath{\sqrt[3]{\,}}}
\newunicodechar{∜}{\ensuremath{\sqrt[4]{\,}}}
\newunicodechar{⃗}{\ensuremath{{}^{\to}}}
\newunicodechar{ⁿ}{\ifmmode^{n}\else\textsuperscript{\ensuremath{n}}\fi}
\newunicodechar{ˣ}{\ifmmode^{x}\else\textsuperscript{\ensuremath{x}}\fi}
\newunicodechar{ᵇ}{\ifmmode^{b}\else\textsuperscript{\ensuremath{b}}\fi}
\newunicodechar{ʳ}{\ifmmode^{r}\else\textsuperscript{\ensuremath{r}}\fi}
\newunicodechar{ᵘ}{\ifmmode^{u}\else\textsuperscript{\ensuremath{u}}\fi}
\newunicodechar{ʸ}{\ifmmode^{y}\else\textsuperscript{\ensuremath{y}}\fi}
\newunicodechar{ᵃ}{\ifmmode^{a}\else\textsuperscript{\ensuremath{a}}\fi}
\newunicodechar{ᵉ}{\ifmmode^{e}\else\textsuperscript{\ensuremath{e}}\fi}
\newunicodechar{ᵏ}{\ifmmode^{k}\else\textsuperscript{\ensuremath{k}}\fi}
\newunicodechar{ᵖ}{\ifmmode^{p}\else\textsuperscript{\ensuremath{p}}\fi}
\newunicodechar{ᴺ}{\ifmmode^{N}\else\textsuperscript{\ensuremath{N}}\fi}
\newunicodechar{ᶜ}{\ifmmode^{c}\else\textsuperscript{\ensuremath{c}}\fi}
\newunicodechar{ʰ}{\ifmmode^{h}\else\textsuperscript{\ensuremath{h}}\fi}
\newunicodechar{ᵠ}{\ifmmode^{q}\else\textsuperscript{\ensuremath{q}}\fi}
\newunicodechar{ₙ}{\ifmmode_{n}\else\textsubscript{\ensuremath{n}}\fi}
\newunicodechar{ₐ}{\ifmmode_{a}\else\textsubscript{\ensuremath{a}}\fi}
\newunicodechar{ᵢ}{\ifmmode_{i}\else\textsubscript{\ensuremath{i}}\fi}
\newunicodechar{ᵣ}{\ifmmode_{r}\else\textsubscript{\ensuremath{r}}\fi}
\newunicodechar{ₖ}{\ifmmode_{k}\else\textsubscript{\ensuremath{k}}\fi}
\newunicodechar{ᵦ}{\ifmmode_{\beta}\else\textsubscript{\ensuremath{\beta}}\fi}
\newunicodechar{ₓ}{\ifmmode_{x}\else\textsubscript{\ensuremath{x}}\fi}
\newfontfamily\popdisplay{ArchivoBlack-Regular.ttf}[Path=../../../fonts/]
\newfontfamily\poplogo{SpaceGrotesk-Bold.ttf}[Path=../../../fonts/]
\newfontfamily\popsemi{PlusJakartaSans-SemiBold.ttf}[Path=../../../fonts/]
\newfontfamily\popxbold{PlusJakartaSans-ExtraBold.ttf}[Path=../../../fonts/]
\newfontfamily\popxboldls{PlusJakartaSans-ExtraBold.ttf}[Path=../../../fonts/,LetterSpace=3.0]

\setlist[enumerate]{leftmargin=1.7em,itemsep=5pt,topsep=4pt}
\setlist[itemize]{leftmargin=1.7em,itemsep=4pt,topsep=4pt,label={\color{poporange}\textbullet}}
\setlength{\parindent}{0pt}
\setlength{\parskip}{5pt}
\renewcommand{\arraystretch}{1.25}

% pgfplots : style Pop par défaut
\pgfplotsset{
  every axis/.append style={axis line style={popink,line width=1pt},tick style={popink},
    grid style={popline,line width=0.5pt},label style={font=\small},tick label style={font=\footnotesize}},
  popcurve/.style={poporange,line width=1.8pt,no markers,smooth},
}
% Même style disponible dans un \draw TikZ simple (hors pgfplots)
\tikzset{popcurve/.style={poporange,line width=1.8pt}}
\tikzset{popgrid/.style={popline,very thin}}
\colorlet{popcurve}{poporange} % le nom du style est parfois utilisé comme couleur

% ============================================================
% Pill / squiggle
% ============================================================
\newcommand{\poppill}[3]{%
  \tikz[baseline=-0.55ex]\node[draw=popink,line width=1.2pt,rounded corners=2.6pt,%
    fill=#2,inner xsep=6.5pt,inner ysep=3pt]{%
    \popxboldls\fontsize{7}{7}\selectfont\color{#3}\MakeUppercase{#1}};%
}
\newcommand{\popsquiggle}[1]{%
  \tikz[baseline=-0.4ex]{
    \draw[#1,line width=2.6pt,line cap=round]
      plot [smooth] coordinates {(0,0.09) (0.34,0.01) (0.68,0.21) (1.02,0.01) (1.36,0.10)};
  }%
}
% Mot-clé surligné (marqueur orange sous le texte)
\newcommand{\cle}[1]{{\popxbold\color{popdark}#1}}

% ============================================================
% En-tête / pied
% ============================================================
\pagestyle{fancy}
\fancyhf{}
\renewcommand{\headrulewidth}{0pt}
\renewcommand{\footrulewidth}{0pt}
\lhead{\raisebox{-0.15ex}{\poplogo\fontsize{13}{13}\selectfont\color{popink}OnBuch\color{poporange}+}}
\rhead{\poppill{\DOCMATIERE\ \textbullet\ \DOCNIVEAU\ \textbullet\ Leçon \DOCLECON}{white}{popink}}
\lfoot{\popxbold\fontsize{7.6}{7.6}\selectfont\color{popmuted}\MakeUppercase{Cours signé OnBuch+}\ \textbullet\ {\popsemi\DOCTITRE}}
\rfoot{\tikz[baseline=-0.6ex]\node[rounded corners=8.5pt,fill=popink,minimum height=17pt,minimum width=17pt,inner xsep=5pt,inner sep=0]{\popxbold\fontsize{8}{8}\selectfont\color{popcream}\thepage\,/\,\pageref*{LastPage}};}

% ============================================================
% Titres
% ============================================================
\newcommand{\chaptitre}[2]{%
  \begin{tcolorbox}[enhanced,colback=poporange,colframe=popink,boxrule=2.6pt,arc=11pt,
      drop shadow=popink,left=15pt,right=15pt,top=12pt,bottom=11pt,before skip=3pt,after skip=18pt]
    \poppill{#2}{popink}{popcream}\par\vspace{9pt}
    {\raggedright\hyphenpenalty=10000\exhyphenpenalty=10000\popdisplay\fontsize{22}{24}\selectfont\color{popcream}\MakeUppercase{#1}\par}
  \end{tcolorbox}
}
% Section numérotée : pastille orange avec le numéro + titre Archivo
\newcounter{popsec}
\newcommand{\coursec}[1]{%
  \stepcounter{popsec}\setcounter{popsub}{0}%
  \par\vspace{16pt}\noindent
  \begin{minipage}{\linewidth}
  \tikz[baseline=-0.7ex]\node[draw=popink,line width=1.6pt,fill=poporange,rounded corners=4pt,minimum size=22pt,inner sep=0]{\popdisplay\fontsize{12}{12}\selectfont\color{popcream}\thepopsec};%
  \hspace{8pt}{\popdisplay\fontsize{15}{17}\selectfont\color{popink}\MakeUppercase{#1}}\par
  \vspace{4pt}\noindent{\color{poporange}\rule{36pt}{3.6pt}}
  \end{minipage}\par\nopagebreak\vspace{8pt}
}
\newcounter{popsub}
\newcommand{\courssub}[1]{%
  \stepcounter{popsub}%
  \par\vspace{9pt}\noindent{\popxbold\fontsize{11.5}{13}\selectfont\color{popink}{\color{poporange}\thepopsec.\thepopsub}\hspace{6pt}#1}\par\nopagebreak\vspace{4pt}
}

% ============================================================
% Boîtes pédagogiques — même recette : fond blanc, bordure épaisse colorée,
% ombre dure encre, badge plein. Titre optionnel [#1].
% ============================================================
\tcbset{popbase/.style={breakable,enhanced,colback=white,boxrule=2.3pt,arc=9pt,
  shadow={3pt}{-3pt}{0pt}{popink},left=14pt,right=14pt,top=11pt,bottom=11pt,before skip=11pt,after skip=15pt,
  fontupper=\popsemi\fontsize{10.6}{14.5}\selectfont\color{popink},
  fontlower=\popsemi\fontsize{10.6}{14.5}\selectfont\color{popink}}}
% Coche dessinée (le glyphe ✓ n'existe pas dans les polices du document)
\newcommand{\popcheck}[1]{\tikz[baseline=-0.5ex]\draw[#1,line width=1.6pt,line cap=round,line join=round](-0.13,0.0)--(-0.04,-0.09)--(0.13,0.1);}
\providecommand{\coche}{\popcheck{popgreen}}
\providecommand{\resultat}[1]{\par\smallskip\noindent\fcolorbox{popgreen}{popgreenL}{\parbox{\dimexpr\linewidth-2\fboxsep-2\fboxrule\relax}{\textbf{Résultat.} #1}}\par\smallskip}
\newcommand{\popboxhead}[3]{\poppill{#1}{#2}{white}\ifx\relax#3\relax\else\hspace{6pt}{\popxbold\fontsize{10.6}{12}\selectfont\color{popink}#3}\fi\par\vspace{7pt}}

\newtcolorbox{aretenir}[1][]{popbase,colframe=popgreen,before upper={\popboxhead{À retenir}{popgreen}{#1}}}
% Texte en espagnol (document de lecture, dialogue, poème, extrait) : cadre
% d'étude. Un titre optionnel (source, auteur) s'affiche dans l'en-tête.
\newtcolorbox{textebox}[1][]{popbase,colframe=popblue,before upper={\popboxhead{Texto}{popblue}{#1}\begin{otherlanguage}{spanish}},after upper={\end{otherlanguage}}}
% Marques ✓ / ✗ (forme correcte / fautive) souvent utilisées par les modèles
\providecommand{\cross}{\textcolor{poppink}{\ensuremath{\times}}}
\providecommand{\xmark}{\cross}\providecommand{\crossmark}{\cross}\providecommand{\cmark}{\ensuremath{\checkmark}}
% Ligne de réponse / justification à compléter (inventée par les modèles)
\providecommand{\justification}[1]{\par\noindent\textit{Justification~:} \dotfill\par}
% \textipa (package tipa non chargé) : la police ne couvre pas l'API -> simple passage
\providecommand{\textipa}[1]{#1}
% Soulignement (ulem non chargé) : \uline, \ul
\providecommand{\uline}[1]{\underline{#1}}\providecommand{\ul}[1]{\underline{#1}}
% Ordinaux français écrits en macro par les modèles : 1\ière, 3\ième
\newcommand{\ière}{\textsuperscript{re}}\newcommand{\ième}{\textsuperscript{e}}
% Mot ou phrase en espagnol à l'intérieur d'un paragraphe français
\newcommand{\esp}[1]{\ifmmode\text{\textit{#1}}\else\begingroup\selectlanguage{spanish}\itshape #1\endgroup\fi}
\newtcolorbox{exemplebox}[1][]{popbase,colframe=poporange,before upper={\popboxhead{Exemple}{poporange}{#1}}}
\newtcolorbox{definition}[1][]{popbase,colframe=popblue,before upper={\popboxhead{Définition}{popblue}{#1}}}
\newtcolorbox{propriete}[1][]{popbase,colframe=poppurple,before upper={\popboxhead{Propriété}{poppurple}{#1}}}
\newtcolorbox{methode}[1][]{popbase,colframe=popgold,before upper={\popboxhead{Méthode}{popgold}{#1}}}
\newtcolorbox{attention}[1][]{popbase,colframe=poppink,before upper={\popboxhead{Attention}{poppink}{#1}}}
\newtcolorbox{experience}[1][]{popbase,colframe=popblue,colback=popblueL,before upper={\popboxhead{Expérience}{popblue}{#1}}}
% « Le savais-tu ? » : carte violette pleine (contexte, culture, Cameroun)
\newtcolorbox{savaistu}[1][]{popbase,colback=poppurple,colframe=popink,
  fontupper=\popsemi\fontsize{10.4}{14.2}\selectfont\color{white},
  before upper={\poppill{Le savais-tu\,?}{white}{poppurple}\ifx\relax#1\relax\else\hspace{6pt}{\popxbold\color{white}#1}\fi\par\vspace{7pt}}}
% Exercice résolu : énoncé en haut, \tcblower puis la solution sur fond crème
\newcounter{popexo}\newcounter{popexr}
\newtcolorbox{exoresolu}[1][]{popbase,colframe=poporange,colbacklower=popcream,
  segmentation style={popink,line width=1.2pt,dashed},
  before upper={\stepcounter{popexr}\popboxhead{Exercice résolu \thepopexr}{poporange}{#1}},
  before lower={\poppill{Solution}{popink}{popcream}\par\vspace{6pt}}}
% Exercice (non résolu en place) — la correction est dans la partie Corrigés
\newtcolorbox{exercice}[1][]{popbase,colframe=popink,
  before upper={\stepcounter{popexo}\popboxhead{Exercice \thepopexo}{popink}{#1}}}
% Corrigé
\newtcolorbox{corrige}[1][]{popbase,colframe=popgreen,colback=popgreenL,
  before upper={\popboxhead{Corrigé}{popgreen}{#1}}}
% Objectifs / prérequis / bilan
\newtcolorbox{objectifs}{popbase,colframe=popink,colback=poporangeL,
  before upper={\popboxhead{Objectifs de la leçon}{popink}{}}}
\newtcolorbox{prerequis}{popbase,colframe=popink,colback=popblueL,
  before upper={\popboxhead{Prérequis}{popblue}{}}}
\newtcolorbox{bilan}{popbase,colframe=popink,colback=white,boxrule=2.8pt,
  before upper={\popboxhead{Fiche bilan}{poporange}{L'essentiel à connaître pour l'examen}}}
% Figure encadrée avec légende
\newtcolorbox{popfigure}[1][]{enhanced,breakable=false,colback=white,colframe=popline,boxrule=1.4pt,arc=8pt,
  left=10pt,right=10pt,top=10pt,bottom=8pt,before skip=10pt,after skip=12pt,halign=center,
  lower separated=false,halign lower=center,
  fontlower=\popsemi\fontsize{9}{11.5}\selectfont\color{popmuted},
  #1}
\newcommand{\legende}[1]{\par\vspace{6pt}{\popsemi\fontsize{9}{11.5}\selectfont\color{popmuted}#1}}

% ============================================================
% Signature OnBuch+
% ============================================================
\newcommand{\poplogoinline}[1]{{\poplogo\fontsize{#1}{#1}\selectfont\color{popink}OnBuch\color{poporange}+}}
\newcommand{\signatureonbuch}{%
  \begin{tcolorbox}[enhanced,colback=popink,colframe=popink,boxrule=2.2pt,arc=10pt,
      drop shadow=poporange,left=14pt,right=14pt,top=10pt,bottom=10pt,before skip=0pt,after skip=0pt]
    \tikz[baseline=-0.7ex]\node[circle,fill=popgreen,draw=popcream,line width=1.3pt,minimum size=22pt,inner sep=0]{\popcheck{white}};%
    \hspace{9pt}\begin{minipage}[c]{0.62\linewidth}
      {\popxbold\fontsize{12}{13}\selectfont\color{popcream}Signé\ \textbullet\ {\poplogo OnBuch\color{poporange}+}}\par\vspace{2pt}
      {\raggedright\popsemi\fontsize{8}{10.5}\selectfont\color{popline}\MakeUppercase{Équipe pédagogique OnBuch+}\\\MakeUppercase{Conforme au programme officiel MINESEC}\par}
    \end{minipage}\hfill
    {\poplogo\fontsize{22}{22}\selectfont\color{popcream}OnBuch\color{poporange}+}
  \end{tcolorbox}%
}

% ============================================================
% Page de garde
% #1 titre · #2 matière · #3 niveau · #4 module · #5 n° leçon · #6 durée ·
% #7 description / objectifs (texte ou itemize)
% ============================================================
\newcommand{\pagegarde}[7]{%
  \thispagestyle{empty}
  \newgeometry{a4paper,margin=1.7cm,top=1.6cm,bottom=1.4cm}%
  \begin{tikzpicture}[remember picture,overlay]
    \fill[poporange!18] ([xshift=-3.2cm,yshift=-3.6cm]current page.north east) circle (4.6cm);
    \fill[poppurple!14] ([xshift=2.4cm,yshift=7.6cm]current page.south west) circle (3.8cm);
  \end{tikzpicture}%
  {\poplogo\fontsize{30}{30}\selectfont\color{popink}OnBuch\color{poporange}+}\hfill
  \raisebox{0.6ex}{\poppill{Cours signé \textbullet\ Premium}{popink}{popcream}}\par
  \vspace{1pt}\popsquiggle{poppurple}\par
  \vspace{8pt}
  \begin{tcolorbox}[enhanced,colback=poporange,colframe=popink,boxrule=2.8pt,arc=13pt,
      drop shadow=popink,left=18pt,right=18pt,top=12pt,bottom=13pt,after skip=10pt]
    \poppill{#2 \textbullet\ #3}{popink}{popcream}\hspace{5pt}\poppill{#4}{white}{popink}\par\vspace{8pt}
    {\popdisplay\fontsize{14}{14}\selectfont\color{popink}LEÇON #5}\par\vspace{4pt}
    {\raggedright\hyphenpenalty=10000\exhyphenpenalty=10000\popdisplay\fontsize{25}{27}\selectfont\color{popcream}\MakeUppercase{#1}\par}
    \vspace{8pt}
    {\popxbold\fontsize{9.5}{9.5}\selectfont\color{popink}\MakeUppercase{Durée conseillée : #6}}
  \end{tcolorbox}
  \begin{tcolorbox}[enhanced,colback=white,colframe=popink,boxrule=2.2pt,arc=10pt,
      drop shadow=popink,left=16pt,right=16pt,top=10pt,bottom=10pt,after skip=8pt]
    \poppill{Dans cette leçon}{poppurple}{white}\par\vspace{5pt}
    \popsemi\fontsize{9.3}{12.6}\selectfont\color{popink}#7
  \end{tcolorbox}
  \noindent\begin{minipage}[t]{0.487\linewidth}
    \begin{tcolorbox}[enhanced,colback=popgreen,colframe=popink,boxrule=2.2pt,arc=10pt,
        drop shadow=popink,left=11pt,right=11pt,top=10pt,bottom=10pt,height=3.1cm,valign=center]
      \tcbox[colback=white,colframe=popink,boxrule=1.3pt,arc=5pt,boxsep=0pt,nobeforeafter,
        left=3pt,right=3pt,top=3pt,bottom=3pt,valign=center]{\qrcode[height=1.9cm]{https://play.google.com/store/apps/details?id=com.luvvix.onbuch&hl=fr}}%
      \hspace{8pt}\begin{minipage}[c]{0.5\linewidth}
        {\popdisplay\fontsize{11}{12.5}\selectfont\color{white}\MakeUppercase{Télécharge OnBuch}}\par\vspace{4pt}
        {\raggedright\popsemi\fontsize{8.4}{11}\selectfont\color{white}Gratuit \textbullet\ Google Play\\Tuteur IA, annales, concours, résultats\par}
      \end{minipage}
    \end{tcolorbox}
  \end{minipage}\hfill
  \begin{minipage}[t]{0.487\linewidth}
    \begin{tcolorbox}[enhanced,colback=poppurple,colframe=popink,boxrule=2.2pt,arc=10pt,
        drop shadow=popink,left=11pt,right=11pt,top=10pt,bottom=10pt,height=3.1cm,valign=center]
      \tikz[baseline=-0.7ex]\node[circle,fill=white,draw=popink,line width=1.3pt,minimum size=1.6cm,inner sep=0]{\popdisplay\fontsize{24}{24}\selectfont\color{poppurple}+};%
      \hspace{8pt}\begin{minipage}[c]{0.6\linewidth}
        {\popdisplay\fontsize{11}{12.5}\selectfont\color{white}\MakeUppercase{Passe à OnBuch+}}\par\vspace{4pt}
        {\raggedright\popsemi\fontsize{8.4}{11}\selectfont\color{white}Tous les cours signés, Tuteur IA illimité, fiches PDF — onglet OnBuch+ de l'app\par}
      \end{minipage}
    \end{tcolorbox}
  \end{minipage}
  \vspace{8pt}
  \popcontenu
  \vfill
  \signatureonbuch
  \restoregeometry
  \newpage
}

% Rangée « Ce cours contient » (page de garde)
\newcommand{\poptile}[3]{% #1 couleur #2 gros texte #3 légende
  \begin{tcolorbox}[enhanced,colback=#1,colframe=popink,boxrule=2pt,arc=9pt,shadow={2.5pt}{-2.5pt}{0pt}{popink},
      left=6pt,right=6pt,top=9pt,bottom=9pt,halign=center,valign=center,height=1.75cm,width=\linewidth,nobeforeafter]
    {\popdisplay\fontsize{12.5}{13}\selectfont\color{white}\MakeUppercase{#2}}\par\vspace{3pt}
    {\centering\popsemi\fontsize{7.8}{9.5}\selectfont\color{white}#3\par}
  \end{tcolorbox}}
\newcommand{\popcontenu}{%
  \par\noindent{\popxboldls\fontsize{7.5}{7.5}\selectfont\color{popmuted}CE COURS SIGNÉ CONTIENT}\par\vspace{7pt}
  \noindent\begin{minipage}[t]{0.235\linewidth}\poptile{popblue}{Cours}{complet, illustré, pas à pas}\end{minipage}\hfill
  \begin{minipage}[t]{0.235\linewidth}\poptile{poporange}{Méthodes}{et exercices résolus}\end{minipage}\hfill
  \begin{minipage}[t]{0.235\linewidth}\poptile{poppink}{Exercices}{type examen + corrigés}\end{minipage}\hfill
  \begin{minipage}[t]{0.235\linewidth}\poptile{popgreen}{Fiche bilan}{l'essentiel à retenir}\end{minipage}\par}

% Sommaire
\newtcolorbox{sommaire}{enhanced,colback=white,colframe=popink,boxrule=2.2pt,arc=10pt,
  drop shadow=popink,left=18pt,right=18pt,top=13pt,bottom=13pt,before skip=6pt,after skip=20pt,
  before upper={\poppill{Sommaire}{poppurple}{white}\par\vspace{9pt}\popsemi\fontsize{10.6}{14.5}\selectfont\color{popink}}}

% Signature de fin de cours — poussée tout en bas de la dernière page
\newcommand{\findecours}{%
  \par\vfill\nopagebreak
  \begin{center}\popsquiggle{poporange}\end{center}\vspace{4pt}
  \signatureonbuch
}
\AtBeginDocument{\def\llbracket{\mathopen{[\mkern-3.5mu[}}\def\rrbracket{\mathclose{]\mkern-3.5mu]}}} % intervalles d'entiers (glyphe ⟦ absent de la police)
% Glyphes absents de Fira Math : redéfinis à partir de glyphes disponibles
\AtBeginDocument{%
  \def\setminus{\mathbin{\backslash}}\let\smallsetminus\setminus
  \def\vdots{\mathord{\vbox{\baselineskip3.2pt\lineskiplimit0pt\kern4pt\hbox{.}\hbox{.}\hbox{.}}}}%
  \def\ddots{\mathinner{\mkern1mu\raise6.5pt\hbox{.}\mkern2mu\raise3.6pt\hbox{.}\mkern2mu\raise0.7pt\hbox{.}\mkern1mu}}%
  \def\triangle{\mathord{\scalebox{0.9}{$\symup{\Delta}$}}}%
  \def\nabla{\mathord{\raisebox{\depth}{\scalebox{1}[-1]{$\symup{\Delta}$}}}}%
  \def\checkmark{\popcheck{popdark}}%
  \def\star{\mathbin{\ast}}\def\bigstar{\mathord{\ast}}%
  \def\diamond{\mathbin{\scalebox{0.6}{$\Diamond$}}}%
  \def\bigcup{\mathop{\mathchoice{\raisebox{-0.3em}{\scalebox{1.7}{$\cup$}}}{\scalebox{1.25}{$\cup$}}{\cup}{\cup}}\displaylimits}%
  \def\bigcap{\mathop{\mathchoice{\raisebox{-0.3em}{\scalebox{1.7}{$\cap$}}}{\scalebox{1.25}{$\cap$}}{\cap}{\cap}}\displaylimits}%
  \def\lesseqqgtr{\mathrel{\lessgtr}}\def\bigoplus{\mathop{\oplus}\displaylimits}%
}
\providecommand{\ssi}{si et seulement si} % abréviation fréquente
\AtBeginDocument{\providecommand{\wideparen}{\overparen}} % arc de cercle

\begin{document}

\pagegarde{El emprendimiento y los intercambios comerciales}{Espagnol}{1ère A}{Módulo 4 : Vie économique}{EA14}{2 h}{Cette leçon mixte du module « Vie économique » permet de lire et commenter un texte sur l'entrepreneuriat, d'acquérir le lexique du commerce et de l'entreprise, et de maîtriser le futur et le conditionnel pour présenter un projet. Elle s'achève par la conception guidée d'un mini-projet entrepreneurial en espagnol.
\vspace{4pt}
\begin{multicols}{2}\small
\begin{itemize}
\item À la fin de cette leçon, je sais lire et comprendre un texte espagnol sur l'entrepreneuriat et le commerce.
\item À la fin de cette leçon, je sais utiliser le lexique de l'entreprise et des échanges commerciaux (emprendedor, empresa, mercado, exportar, importar, cliente, beneficio, inversión).
\item À la fin de cette leçon, je sais conjuguer et employer le futur simple pour annoncer un projet (crearé, exportaré).
\item À la fin de cette leçon, je sais utiliser la périphrase « voy a + infinitivo » pour exprimer un projet proche.
\item À la fin de cette leçon, je sais employer le conditionnel (me gustaría, podría, querría) pour exprimer un souhait ou un projet hypothétique.
\item À la fin de cette leçon, je sais distinguer futur proche, futur simple et conditionnel selon le degré de réalisation du projet.
\end{itemize}\end{multicols}}

\begin{sommaire}
\begin{multicols}{2}
\begin{enumerate}
\item Texte support : « Un sueño emprendedor en Duala »
\item Lexique thématique : empresa y comercio
\item Grammaire 1 : le futur du projet (futur proche et futur simple)
\item Grammaire 2 : le conditionnel du souhait (me gustaría)
\item Méthode du commentaire et commentaire modèle
\item Production guidée et entraînement à la traduction
\item Activité d'intégration
\item Méthodes et exercices résolus
\item Exercices
\item Corrigés des exercices
\item Fiche bilan
\end{enumerate}
\end{multicols}
\end{sommaire}

\begin{objectifs}
\begin{itemize}
\item À la fin de cette leçon, je sais lire et comprendre un texte espagnol sur l'entrepreneuriat et le commerce.
\item À la fin de cette leçon, je sais utiliser le lexique de l'entreprise et des échanges commerciaux (emprendedor, empresa, mercado, exportar, importar, cliente, beneficio, inversión).
\item À la fin de cette leçon, je sais conjuguer et employer le futur simple pour annoncer un projet (crearé, exportaré).
\item À la fin de cette leçon, je sais utiliser la périphrase « voy a + infinitivo » pour exprimer un projet proche.
\item À la fin de cette leçon, je sais employer le conditionnel (me gustaría, podría, querría) pour exprimer un souhait ou un projet hypothétique.
\item À la fin de cette leçon, je sais distinguer futur proche, futur simple et conditionnel selon le degré de réalisation du projet.
\item À la fin de cette leçon, je sais présenter oralement et par écrit un projet entrepreneurial simple.
\item À la fin de cette leçon, je sais parler des échanges commerciaux du Cameroun avec l'Afrique et le monde hispanophone.
\item À la fin de cette leçon, je sais traduire des phrases de thème et de version portant sur le commerce et les projets.
\end{itemize}
\end{objectifs}

\begin{prerequis}
\begin{itemize}
\item Présent de l'indicatif des verbes réguliers et irréguliers courants (Seconde)
\item Verbes d'opinion et de volonté : querer, poder, gustar (Seconde)
\item Lexique de base de la vie quotidienne et de l'argent (dinero, trabajo, comprar, vender)
\item Structure de la phrase simple et place de l'adjectif (Seconde)
\item Notions de géographie : pays hispanophones et capitales (début d'année)
\end{itemize}
\end{prerequis}

\newpage

\coursec{Texte support : « Un sueño emprendedor en Duala »}

\courssub{Situation d'accroche et problématique}

À Douala, une jeune diplômée transforme le manioc en chips qu'elle exporte vers le Gabon et rêve d'atteindre le marché mexicain. \esp{Voy a crear mi propia empresa}, dit-elle ; \esp{me gustaría exportar a todo el mundo}. Derrière cette phrase simple se cachent deux questions essentielles : comment \cle{présenter un projet entrepreneurial} en espagnol ? Comment \cle{parler des échanges commerciaux} du Cameroun avec l'Afrique et le monde hispanophone ? C'est tout l'enjeu de cette leçon.

Dans cette première étape, nous allons lire ensemble un portrait : celui de Mballa, jeune entrepreneure camerounaise. Nous observerons d'abord le texte, nous en dégagerons le vocabulaire économique, puis nous repérerons les marqueurs grammaticaux du projet (\esp{voy a crear}, \esp{crearé}, \esp{me gustaría}), que nous étudierons en détail dans les sections suivantes.

\courssub{Lecture globale du texte}

\begin{textebox}[Un sueño emprendedor en Duala]
Mballa tiene 22 años y vive en Duala, la capital económica de Camerún. Después de sus estudios, decidió no buscar un empleo, sino crear su propia empresa. Con una pequeña inversión de su familia, compró máquinas para transformar el maníoc en chips. Hoy, su pequeña fábrica emplea a cinco personas y vende sus productos en el mercado local. Gracias a la calidad de sus chips, Mballa ya exporta a Gabón, donde tiene muchos clientes fieles.

« Voy a crear una empresa más grande, dice Mballa con una sonrisa. Crearé diez empleos nuevos y me gustaría exportar a Guinea Ecuatorial y a México. Mis clientes aprecian la calidad, y mi proyecto tendrá muchos beneficios. Con esos beneficios, compraré más materia prima y mejoraré la balanza comercial de mi país. »

Mballa representa a una nueva generación de jóvenes emprendedores africanos: con una idea, una inversión y mucho trabajo, quieren transformar la economía de su continente.
\end{textebox}

\textbf{Traduction du texte.} Mballa a 22 ans et vit à Douala, la capitale économique du Cameroun. Après ses études, elle a décidé de ne pas chercher un emploi, mais de créer sa propre entreprise. Avec un petit investissement de sa famille, elle a acheté des machines pour transformer le manioc en chips. Aujourd'hui, sa petite usine emploie cinq personnes et vend ses produits sur le marché local. Grâce à la qualité de ses chips, Mballa exporte déjà vers le Gabon, où elle a de nombreux clients fidèles. « Je vais créer une entreprise plus grande, dit Mballa avec un sourire. Je créerai dix nouveaux emplois et j'aimerais exporter vers la Guinée équatoriale et le Mexique. Mes clients apprécient la qualité, et mon projet aura beaucoup de bénéfices. Avec ces bénéfices, j'achèterai plus de matière première et j'améliorerai la balance commerciale de mon pays. » Mballa représente une nouvelle génération de jeunes entrepreneurs africains : avec une idée, un investissement et beaucoup de travail, ils veulent transformer l'économie de leur continent.

\begin{definition}[Glossaire : le vocabulaire de l'entreprise et du commerce]
\begin{itemize}
\item \esp{emprendedor / emprendedora} : entrepreneur / entrepreneure (personne qui crée et dirige une entreprise en assumant des risques).
\item \esp{empresa} : entreprise, société.
\item \esp{proyecto} : projet.
\item \esp{mercado} : marché.
\item \esp{exportar} : exporter ; \esp{importar} : importer.
\item \esp{cliente} : client(e).
\item \esp{beneficio} : bénéfice, profit.
\item \esp{inversión} : investissement.
\item \esp{materia prima} : matière première.
\item \esp{balanza comercial} : balance commerciale (différence entre exportations et importations).
\item \esp{empleo} : emploi ; \esp{fábrica} : usine ; \esp{maníoc} : manioc.
\end{itemize}
\end{definition}

\begin{attention}[Faux amis à éviter]
\esp{empresa} ne signifie pas « emprise » mais « entreprise ». De même, \esp{beneficio} signifie « bénéfice, profit » (gain d'argent), et non « bienfait ». Enfin, \esp{exportar} se prononce « eks-por-tar » avec un x qui sonne « ks » : attention à ne pas écrire \esp{esportar}.
\end{attention}

\courssub{Méthode de lecture d'un texte économique}

\begin{methode}[Lire un texte sur l'entrepreneuriat en trois étapes]
\begin{enumerate}
\item \textbf{Identifier le projet} : qui est le personnage ? Que veut-il créer ? (Cherche les verbes au futur et au conditionnel : \esp{voy a crear}, \esp{crearé}, \esp{me gustaría}.)
\item \textbf{Repérer les moyens} : avec quelles ressources agit-il ? (\esp{inversión}, \esp{máquinas}, \esp{materia prima}, employés.)
\item \textbf{Noter les résultats attendus} : quels bénéfices, pour qui, vers quels marchés ? (\esp{beneficios}, \esp{clientes}, \esp{exportar a...}.)
\end{enumerate}
\end{methode}

Appliquons cette méthode au texte :

\begin{center}
\begin{tabularx}{\linewidth}{|X|X|}\hline
\rowcolor{popblueL} \textbf{Question de repérage} & \textbf{Réponse dans le texte} \\ \hline
Qui ? & Mballa, 22 ans, jeune entrepreneure \\ \hline
Quoi ? & Transformation du manioc en chips \\ \hline
Où ? & À Douala ; vente au marché local, exportation vers le Gabon \\ \hline
Avec quels moyens ? & Une petite inversión familiale, des máquinas, cinq employés \\ \hline
Objectifs ? & Crear una empresa más grande, crear diez empleos, exportar a Guinea Ecuatorial y México \\ \hline
Résultats attendus ? & Muchos beneficios, mejorar la balanza comercial del país \\ \hline
\end{tabularx}
\end{center}

\courssub{Carte mentale du lexique}

\begin{popfigure}
\begin{center}
\begin{tikzpicture}[mindmap, grow cyclic, every node/.style={concept, align=center, font=\small},
root concept/.append style={concept color=popblue, font=\small\bfseries},
level 1/.append style={level distance=2.9cm, sibling angle=72},
level 1 concept/.append style={concept color=poporangeL}]
\node[root concept]{El emprendimiento}
child { node[align=center]{idea \\ \emph{idée}} }
child { node[align=center]{inversión \\ \emph{investissement}} }
child { node[align=center]{mercado \\ \emph{marché}} }
child { node[align=center]{clientes \\ \emph{clients}} }
child { node[align=center]{beneficio \\ \emph{bénéfice}} };
\end{tikzpicture}
\end{center}
\legende{Carte mentale : les cinq mots-clés du projet entrepreneurial.}
\end{popfigure}

\courssub{Questions de compréhension}

\begin{exercice}[Compréhension du texte]
\textbf{A. Questions littérales (la réponse est dans le texte)}
\begin{enumerate}
\item ¿Cuántos años tiene Mballa y dónde vive?
\item ¿Qué compró Mballa con la inversión de su familia?
\item ¿A qué país exporta Mballa actualmente?
\end{enumerate}
\textbf{B. Questions d'interprétation (il faut réfléchir)}
\begin{enumerate}
\item ¿Por qué podemos decir que Mballa es una verdadera emprendedora?
\item ¿Por qué la exportación es importante para la balanza comercial de Camerún?
\end{enumerate}
\textbf{C. Expression personnelle}
\begin{enumerate}
\item ¿Te gustaría crear tu propia empresa? ¿Qué producto transformarías o venderías?
\end{enumerate}
\end{exercice}

\begin{corrige}[Questions littérales]
\begin{enumerate}
\item \esp{Mballa tiene 22 años y vive en Duala, la capital económica de Camerún.} (Mballa a 22 ans et vit à Douala.)
\item \esp{Compró máquinas para transformar el maníoc en chips.} (Elle a acheté des machines pour transformer le manioc en chips.)
\item \esp{Actualmente exporta a Gabón.} (Actuellement, elle exporte vers le Gabon.)
\end{enumerate}
Pour les questions d'interprétation : Mballa est une vraie entrepreneure parce qu'elle crée son entreprise, assume des risques (\esp{inversión}) et crée des emplois. L'exportation rapporte des devises au pays : si le Cameroun exporte plus qu'il n'importe, sa \esp{balanza comercial} s'améliore.
\end{corrige}

\courssub{Relevé des marqueurs de projet}

Relisons les paroles de Mballa et soulignons les trois formes verbales qui expriment son projet :

\begin{exemplebox}[Trois façons d'exprimer le projet]
\begin{itemize}
\item \esp{Voy a crear una empresa más grande.} « Je vais créer une entreprise plus grande. » — \cle{futur proche} : \esp{ir a + infinitivo} (projet proche, déjà décidé).
\item \esp{Crearé diez empleos nuevos.} « Je créerai dix nouveaux emplois. » — \cle{futur simple} : \esp{crearé} (promesse, prévision ferme).
\item \esp{Me gustaría exportar a Guinea Ecuatorial y a México.} « J'aimerais exporter vers la Guinée équatoriale et le Mexique. » — \cle{conditionnel de souhait} : \esp{me gustaría + infinitivo} (rêve, désir poli).
\end{itemize}
\end{exemplebox}

\begin{center}
\begin{tabularx}{\linewidth}{|X|X|X|}\hline
\rowcolor{popblueL} \textbf{Français} & \textbf{Español} & \textbf{Remarque} \\ \hline
Je vais créer & Voy a crear & futur proche, projet décidé \\ \hline
Je créerai & Crearé & futur simple, promesse \\ \hline
J'aimerais exporter & Me gustaría exportar & conditionnel, souhait \\ \hline
Mon projet aura des bénéfices & Mi proyecto tendrá beneficios & futur simple, prévision \\ \hline
\end{tabularx}
\end{center}

\begin{exoresolu}[Repérer les marqueurs de projet]
Énoncé : dans la phrase \esp{« Voy a crear una empresa más grande, crearé diez empleos y me gustaría exportar a México »}, identifie les trois marqueurs de projet et précise la nuance de chacun.
\tcblower
Solution : \esp{voy a crear} = futur proche (\esp{ir a + infinitivo}), le projet est déjà décidé et proche ; \esp{crearé} = futur simple de \esp{crear}, une promesse ferme ; \esp{me gustaría exportar} = conditionnel de \esp{gustar}, un souhait encore incertain. L'entrepreneure passe ainsi du projet concret au rêve lointain.
\end{exoresolu}

\begin{savaistu}[La Guinée équatoriale, voisine hispanophone]
La Guinée équatoriale est le seul pays d'Afrique subsaharienne dont l'espagnol est langue officielle. Frontalière du Cameroun, elle est un partenaire commercial naturel : un entrepreneur camerounais qui parle espagnol peut y vendre ses produits sans barrière linguistique. C'est pourquoi Mballa rêve d'exporter d'abord à Malabo et à Bata, avant de viser le grand marché du Mexique.
\end{savaistu}

\courssub{Lecture expressive à voix haute}

\begin{experience}[Lecture expressive par groupes]
Par groupes de quatre, lisez le texte à voix haute : un élève lit le premier paragraphe (le portrait), un deuxième lit les paroles de Mballa avec enthousiasme (comme dans une interview), un troisième lit le dernier paragraphe, et le quatrième vérifie la prononciation des mots-clés : \esp{emprendedora} « em-pren-de-do-ra », \esp{inversión} « in-ber-sion », \esp{beneficios} « bé-né-fi-sios », \esp{Guinea Ecuatorial} « gui-né-a é-koua-to-rial ». Puis chaque groupe choisit sa meilleure « voix de Mballa » pour lire les paroles du projet devant la classe.
\end{experience}

\begin{aretenir}[Ce qu'il faut retenir de cette lecture]
Le texte présente un \cle{projet entrepreneurial} en trois temps : l'idée et les moyens (\esp{inversión}, \esp{máquinas}), l'action présente (\esp{vende}, \esp{exporta}) et le projet futur (\esp{voy a crear}, \esp{crearé}, \esp{me gustaría}). Le vocabulaire essentiel est : \esp{emprendedor}, \esp{empresa}, \esp{proyecto}, \esp{mercado}, \esp{exportar}, \esp{importar}, \esp{cliente}, \esp{beneficio}, \esp{inversión}, \esp{materia prima}, \esp{balanza comercial}. Ces mots et ces trois marqueurs verbaux seront les outils de toute la leçon.
\end{aretenir}

\coursec{Lexique thématique : empresa y comercio}

Dans la section précédente, nous avons lu le texte \esp{Un sueño emprendedor en Duala}, l'histoire d'Amina, une jeune Camerounaise qui crée sa petite entreprise de jus de fruits. Ce texte contenait de nombreux mots du monde de l'économie : \esp{empresa}, \esp{mercado}, \esp{clientes}, \esp{beneficio}... Pour bien comprendre ce texte, le commenter et, surtout, pour savoir présenter ton propre projet entrepreneurial, il faut maintenant organiser et mémoriser ce vocabulaire. C'est l'objectif de cette section : construire, champ par champ, ton lexique de l'entreprise et du commerce.

\courssub{Observation : le lexique en contexte}

Avant d'apprendre des listes de mots, observons-les dans un texte authentique. Lis ce court article de presse économique et souligne mentalement tous les mots qui appartiennent au monde de l'entreprise et du commerce.

\begin{textebox}[La economía camerunesa en movimiento]
Duala, la capital económica del Camerún, es una ciudad llena de energía. Cada mañana, miles de \textbf{trabajadores} van a sus \textbf{empresas} y miles de \textbf{emprendedores} abren sus tiendas en los mercados.

El puerto de Duala es el corazón del comercio del país. De allí salen barcos cargados de petróleo, cacao, café y madera: son las principales \textbf{exportaciones} camerunesas. El Camerún \textbf{exporta} estas materias primas a muchos países del mundo, entre ellos España y China, que son importantes \textbf{socios comerciales}. A cambio, el país \textbf{importa} máquinas, coches, medicinas y productos electrónicos. Estas \textbf{importaciones} llegan también por el puerto y pasan por la \textbf{aduana}.

En el \textbf{mercado} central de la ciudad, los comerciantes compiten para \textbf{ganar clientes}. Ofrecen productos de buena \textbf{calidad} a un \textbf{precio} justo. La \textbf{competencia} es fuerte, pero es buena para los \textbf{consumidores}, porque los precios bajan.

Muchos jóvenes cameruneses sueñan con \textbf{crear una empresa}. Para \textbf{lanzar un proyecto}, necesitan una \textbf{inversión} inicial. Algunos piden un \textbf{préstamo} al banco; otros buscan un \textbf{socio} que ponga dinero en el negocio. Si el proyecto funciona, la empresa \textbf{obtiene beneficios}; si no funciona, hay \textbf{pérdidas} y a veces \textbf{deudas}.

Amina, la joven de nuestro texto, es un buen ejemplo: ella \textbf{invirtió} sus ahorros en una pequeña fábrica de zumos, conquistó el mercado de su barrio y hoy su empresa da \textbf{empleo} a cinco personas. Su historia demuestra que el \textbf{emprendimiento} puede transformar la vida de una persona y de todo un país.
\end{textebox}

\textbf{Glossaire :}
\begin{itemize}
\item \esp{los ahorros} : les économies ; \esp{la fábrica de zumos} : l'usine de jus ; \esp{el barrio} : le quartier
\item \esp{las materias primas} : les matières premières ; \esp{la madera} : le bois ; \esp{el barco} : le bateau
\item \esp{poner dinero} : investir de l'argent ; \esp{funcionar} : marcher, réussir ; \esp{los comerciantes} : les commerçants
\item \esp{cargado de} : chargé de ; \esp{a cambio} : en échange ; \esp{compiten} : ils rivalisent (verbe \esp{competir}, irrégulier : e $\rightarrow$ i)
\end{itemize}

\textbf{Questions de compréhension :}
\begin{enumerate}
\item ¿Qué productos exporta el Camerún?
\item ¿Qué productos importa el país?
\item ¿Por qué la competencia es buena para los consumidores?
\item ¿Qué necesita un joven para lanzar un proyecto?
\item ¿En qué invirtió Amina sus ahorros?
\end{enumerate}

\courssub{Champ 1 : la empresa y el emprendedor}

Ce premier champ lexical désigne les personnes et les structures qui font l'économie.

\begin{definition}[Los actores de la economía]
\begin{itemize}
\item \esp{el/la emprendedor(a)} : l'entrepreneur, la personne qui crée une activité. \esp{Amina es una emprendedora valiente.} « Amina est une entrepreneuse courageuse. »
\item \esp{el empresario / la empresaria} : le chef d'entreprise, la cheffe d'entreprise. \esp{El empresario dirige tres fábricas.} « Le chef d'entreprise dirige trois usines. »
\item \esp{la empresa} : l'entreprise. \esp{Su empresa tiene veinte trabajadores.} « Son entreprise compte vingt employés. »
\item \esp{el proyecto} : le projet. \esp{Mi proyecto es vender frutas en el mercado.} « Mon projet est de vendre des fruits au marché. »
\item \esp{el socio / la socia} : l'associé, la partenaire. \esp{Busco un socio para mi negocio.} « Je cherche un associé pour mon affaire. »
\item \esp{el empleo} : l'emploi. \esp{La empresa crea empleo para los jóvenes.} « L'entreprise crée de l'emploi pour les jeunes. »
\item \esp{el/la trabajador(a)} : le travailleur, l'employé. \esp{Los trabajadores llegan a las siete.} « Les employés arrivent à sept heures. »
\end{itemize}
\end{definition}

\begin{attention}[Emprendedor ou empresario ?]
Le français utilise souvent « entrepreneur » pour les deux notions. En espagnol, distingue bien : \esp{el emprendedor} est celui qui \emph{lance} une activité nouvelle (souvent petite et innovante), tandis que \esp{el empresario} est le chef d'une entreprise déjà établie. Amina au début est \esp{emprendedora} ; dix ans plus tard, avec trois usines, elle sera \esp{empresaria}. Remarque aussi la formation des féminins : \esp{emprendedor} $\rightarrow$ \esp{emprendedora}, mais \esp{empresario} $\rightarrow$ \esp{empresaria}.
\end{attention}

\courssub{Champ 2 : el mercado y los clientes}

\begin{definition}[El mercado]
\begin{itemize}
\item \esp{el mercado} : le marché (le lieu, mais aussi le marché économique). \esp{Hay mucha demanda de zumos naturales en el mercado.} « Il y a une forte demande de jus naturels sur le marché. »
\item \esp{el cliente / la clienta} : le client, la cliente. \esp{Los clientes están contentos con el producto.} « Les clients sont contents du produit. »
\item \esp{el consumidor / la consumidora} : le consommateur, la consommatrice. \esp{Los consumidores prefieren la calidad.} « Les consommateurs préfèrent la qualité. »
\item \esp{la competencia} : la concurrence. \esp{La competencia es muy fuerte en Duala.} « La concurrence est très forte à Douala. »
\item \esp{el precio} : le prix. \esp{El precio del cacao sube.} « Le prix du cacao augmente. »
\item \esp{la calidad} : la qualité. \esp{Vendemos calidad a buen precio.} « Nous vendons de la qualité à bon prix. »
\item \esp{la demanda} : la demande ; \esp{la oferta} : l'offre. \esp{La oferta y la demanda deciden los precios.} « L'offre et la demande déterminent les prix. »
\end{itemize}
\end{definition}

\begin{attention}[Faux amis et pièges de genre]
\begin{itemize}
\item \esp{el cliente} a pour féminin \esp{la clienta} (on trouve aussi \esp{la cliente}). Ne traduis jamais « la cliente » par « la clienta » mot à mot sans vérifier : en espagnol les deux formes existent, mais \esp{la clienta} est la plus courante.
\item \esp{la competencia} signifie d'abord « la concurrence » dans un contexte économique : \esp{ganar a la competencia} = « battre la concurrence ». Le mot peut aussi signifier « la compétence » (le savoir-faire) dans d'autres contextes, mais dans cette leçon, retiens : \esp{competencia} = concurrence.
\item \esp{el precio} se prononce « pré-sio » (le \esp{c} devant \esp{i} se prononce comme un « s » en Amérique latine et comme le « z » espagnol en Espagne). Ne l'écris jamais « precio » sans réfléchir à sa prononciation, et ne le confonds pas avec \esp{el premio} (le prix au sens de récompense).
\end{itemize}
\end{attention}

\courssub{Champ 3 : los intercambios internacionales}

\begin{definition}[El comercio internacional]
\begin{itemize}
\item \esp{exportar} : exporter. \esp{El Camerún exporta petróleo y cacao.} « Le Cameroun exporte du pétrole et du cacao. »
\item \esp{importar} : importer. \esp{El país importa máquinas y coches.} « Le pays importe des machines et des voitures. »
\item \esp{la exportación / la importación} : l'exportation / l'importation. \esp{La exportación de cacao es importante.} « L'exportation du cacao est importante. »
\item \esp{la aduana} : la douane. \esp{Las mercancías pasan por la aduana del puerto.} « Les marchandises passent par la douane du port. »
\item \esp{los aranceles} : les droits de douane, les tarifs douaniers. \esp{Los aranceles encarecen los productos importados.} « Les droits de douane renchérissent les produits importés. »
\item \esp{la balanza comercial} : la balance commerciale. \esp{La balanza comercial compara exportaciones e importaciones.} « La balance commerciale compare les exportations et les importations. »
\item \esp{el socio comercial} : le partenaire commercial. \esp{España es un socio comercial del Camerún.} « L'Espagne est un partenaire commercial du Cameroun. »
\end{itemize}
\end{definition}

\begin{exemplebox}[Exemples traduits]
\begin{itemize}
\item \esp{El Camerún exporta petróleo y cacao e importa máquinas.} « Le Cameroun exporte du pétrole et du cacao et importe des machines. » (Remarque : la conjonction \esp{y} devient \esp{e} devant un mot commençant par le son « i » : \esp{máquinas e importaciones}, \esp{padre e hijo}.)
\item \esp{Guinea Ecuatorial exporta petróleo a España.} « La Guinée équatoriale exporte du pétrole vers l'Espagne. »
\item \esp{Las exportaciones de café pasan por el puerto de Duala.} « Les exportations de café passent par le port de Douala. »
\item \esp{China es un socio comercial muy importante para África.} « La Chine est un partenaire commercial très important pour l'Afrique. »
\end{itemize}
\end{exemplebox}

\begin{savaistu}[El comercio del Camerún]
Le Cameroun exporte surtout des matières premières : le pétrole, le cacao (le pays est l'un des grands producteurs mondiaux), le café et le bois. Ses partenaires commerciaux incluent des pays européens comme l'Espagne et la France, ainsi que la Chine. Et savais-tu que la Guinée équatoriale, voisine du Cameroun, est le seul pays d'Afrique où l'espagnol est langue officielle ? C'est une porte naturelle vers le monde hispanophone pour les commerçants camerounais : parler espagnol, c'est pouvoir négocier directement avec des millions de clients et de partenaires en Espagne et en Amérique latine !
\end{savaistu}

\courssub{Champ 4 : el dinero y los resultados}

\begin{definition}[Las finanzas de la empresa]
\begin{itemize}
\item \esp{la inversión} : l'investissement. \esp{Necesito una inversión de quinientos mil francos.} « J'ai besoin d'un investissement de cinq cent mille francs. »
\item \esp{invertir} : investir. \esp{Amina invirtió sus ahorros en la fábrica.} « Amina a investi ses économies dans l'usine. »
\item \esp{el beneficio} : le bénéfice, le profit. \esp{La empresa obtiene beneficios cada mes.} « L'entreprise réalise des bénéfices chaque mois. »
\item \esp{la ganancia} : le gain. \esp{La ganancia del año fue buena.} « Le gain de l'année a été bon. »
\item \esp{la pérdida} : la perte. \esp{El primer año hubo pérdidas.} « La première année, il y a eu des pertes. »
\item \esp{el préstamo} : le prêt, l'emprunt. \esp{Pedí un préstamo al banco.} « J'ai demandé un prêt à la banque. »
\item \esp{la deuda} : la dette. \esp{La empresa tiene deudas con el banco.} « L'entreprise a des dettes envers la banque. »
\end{itemize}
\end{definition}

\begin{attention}[Invertir : un verbe à changement radical]
\esp{Invertir} est un verbe irrégulier à double changement de radical (comme \esp{sentir} ou \esp{preferir}) : e $\rightarrow$ ie au présent (sauf à \esp{nosotros} et \esp{vosotros}), et e $\rightarrow$ i aux 3\up{es} personnes du passé simple.
\begin{center}
\begin{tabular}{|l|l|l|}
\hline
\rowcolor{popblueL} \textbf{Presente} & \textbf{Pretérito} & \textbf{Traduction} \\
\hline
yo \esp{invierto} & yo \esp{invertí} & j'investis / j'ai investi \\
\hline
tú \esp{inviertes} & tú \esp{invertiste} & tu investis / tu as investi \\
\hline
él \esp{invierte} & él \esp{invirtió} & il investit / il a investi \\
\hline
nosotros \esp{invertimos} & nosotros \esp{invertimos} & nous investissons / nous avons investi \\
\hline
vosotros \esp{invertís} & vosotros \esp{invertisteis} & vous investissez / vous avez investi \\
\hline
ellos \esp{invierten} & ellos \esp{invirtieron} & ils investissent / ils ont investi \\
\hline
\end{tabular}
\end{center}
Ne dis jamais « yo inverto » ! Et autre piège : \esp{el éxito} signifie « le succès » (prononcé « èk-si-to »), ne le confonds pas avec le mot anglais \emph{exit} (la sortie se dit \esp{la salida}). \esp{Su empresa tiene mucho éxito.} « Son entreprise a beaucoup de succès. »
\end{attention}

\courssub{Las familias de palabras}

Une excellente stratégie pour enrichir ton vocabulaire est de regrouper les mots par famille : un même radical donne un verbe, un nom d'action et un nom de personne.

\begin{propriete}[Formar familias de palabras]
En espagnol, beaucoup de noms d'action se forment avec le suffixe \esp{-ción} (comme le français « -tion ») et les noms de personne avec \esp{-dor/-dora} :
\begin{center}
\begin{tabularx}{\linewidth}{|X|X|X|}
\hline
\rowcolor{popblueL} \textbf{Verbo} & \textbf{Acción (-ción)} & \textbf{Persona (-dor)} \\
\hline
\esp{exportar} (exporter) & \esp{la exportación} & \esp{el exportador} \\
\hline
\esp{importar} (importer) & \esp{la importación} & \esp{el importador} \\
\hline
\esp{invertir} (investir) & \esp{la inversión} & \esp{el inversor} \\
\hline
\esp{emprender} (entreprendre) & \esp{el emprendimiento} & \esp{el emprendedor} \\
\hline
\esp{trabajar} (travailler) & \esp{el trabajo} & \esp{el trabajador} \\
\hline
\esp{consumir} (consommer) & \esp{el consumo} & \esp{el consumidor} \\
\hline
\end{tabularx}
\end{center}
Attention aux exceptions de formation : \esp{invertir} donne \esp{el inversor} (et non « invertidor »), et \esp{emprender} donne un nom d'action en \esp{-miento} : \esp{el emprendimiento}.
\end{propriete}

\begin{propriete}[Los nombres en -ción son femeninos]
Tous les noms espagnols en \esp{-ción} sont \textbf{féminins} : \esp{la exportación}, \esp{la importación}, \esp{la inversión}, \esp{la producción}. Ils correspondent presque toujours à un mot français en « -tion » : \esp{la exportación} = « l'exportation », \esp{la producción} = « la production » (attention parfois le français change de mot : \esp{la inversión} = « l'investissement »). Leur pluriel est en \esp{-ciones}, \textbf{sans accent écrit} : \esp{las exportaciones}, \esp{las inversiones}.
\end{propriete}

Voici une carte mentale pour visualiser les familles de mots autour du commerce :

\begin{popfigure}
\begin{center}
\begin{tikzpicture}[mindmap, grow cyclic, every node/.style={concept, align=center, font=\small},
root concept/.append style={concept color=popblue, font=\small\bfseries},
level 1/.append style={level distance=3.2cm, sibling angle=90},
level 2/.append style={level distance=2.2cm, sibling angle=45, font=\footnotesize}]
\node[root concept]{COMERCIO}
  child[concept color=poporange] { node {exportar}
    child { node[concept color=poporangeL] {la exportación} }
    child { node[concept color=poporangeL] {el exportador} } }
  child[concept color=popgreen] { node {importar}
    child { node[concept color=popgreenL] {la importación} }
    child { node[concept color=popgreenL] {el importador} } }
  child[concept color=poppurple] { node {invertir}
    child { node[concept color=poppurpleL] {la inversión} }
    child { node[concept color=poppurpleL] {el inversor} } }
  child[concept color=poppink] { node {emprender}
    child { node[concept color=poppinkL] {el emprendi- miento} }
    child { node[concept color=poppinkL] {el emprendedor} } };
\end{tikzpicture}
\end{center}
\legende{Carte mentale : les familles de mots du commerce (verbe $\rightarrow$ action $\rightarrow$ personne)}
\end{popfigure}

\courssub{Expresiones útiles para hablar de un proyecto}

Pour présenter un projet entrepreneurial à l'oral ou à l'écrit, apprends ces expressions figées (verbe + nom) :

\begin{exemplebox}[Las expresiones del emprendedor]
\begin{itemize}
\item \esp{crear una empresa} : créer une entreprise. \esp{Quiero crear una empresa de ropa.} « Je veux créer une entreprise de vêtements. »
\item \esp{lanzar un proyecto} : lancer un projet. \esp{Vamos a lanzar un proyecto agrícola.} « Nous allons lancer un projet agricole. »
\item \esp{abrir un mercado} : ouvrir un marché. \esp{La empresa abrió un mercado en Gabón.} « L'entreprise a ouvert un marché au Gabon. »
\item \esp{ganar clientes} : gagner des clients. \esp{Con buena calidad, ganamos clientes.} « Avec de la bonne qualité, nous gagnons des clients. »
\item \esp{obtener beneficios} : réaliser des bénéfices. \esp{El segundo año obtuvimos beneficios.} « La deuxième année, nous avons réalisé des bénéfices. »
\item \esp{pedir un préstamo} : demander un prêt. \esp{Pedí un préstamo para comprar máquinas.} « J'ai demandé un prêt pour acheter des machines. »
\end{itemize}
\end{exemplebox}

Le schéma suivant résume le cycle économique d'une entreprise qui réussit :

\begin{popfigure}
\begin{center}
\begin{tikzpicture}[node distance=0.5cm and 0.9cm, every node/.style={align=center, font=\small},
caja/.style={draw=popblue, fill=popblueL, rounded corners, thick, minimum width=2.4cm, minimum height=0.9cm},
flecha/.style={-{Stealth[length=3mm]}, very thick, poporange}]
\node[caja] (inv) {inversión};
\node[caja, right=of inv] (prod) {producción};
\node[caja, right=of prod] (merc) {mercado};
\node[caja, below=0.8cm of merc] (cli) {clientes};
\node[caja, left=of cli] (ben) {beneficio};
\node[caja, left=of ben, fill=poporangeL, draw=poporange] (reinv) {reinversión};
\draw[flecha] (inv) -- (prod);
\draw[flecha] (prod) -- (merc);
\draw[flecha] (merc) -- (cli);
\draw[flecha] (cli) -- (ben);
\draw[flecha] (ben) -- (reinv);
\draw[flecha, dashed] (reinv.north) .. controls +(0,0.9) and +(0,-0.9) .. (inv.south);
\end{tikzpicture}
\end{center}
\legende{Le cycle de l'entreprise : l'investissement produit, le marché attire les clients, le bénéfice permet de réinvestir}
\end{popfigure}

\courssub{Tableau récapitulatif et mini-dictionnaire}

Voici le tableau de synthèse de tout le lexique de la leçon, classé par champ :

\begin{center}
\begin{tabularx}{\linewidth}{|l|X|X|}
\hline
\rowcolor{popblueL} \textbf{Campo} & \textbf{Español} & \textbf{Français} \\
\hline
Empresa & \esp{el emprendedor / la empresa / el proyecto / el socio / el empleo / el trabajador} & l'entrepreneur / l'entreprise / le projet / l'associé / l'emploi / le travailleur \\
\hline
Mercado & \esp{el mercado / el cliente / el consumidor / la competencia / el precio / la calidad / la oferta / la demanda} & le marché / le client / le consommateur / la concurrence / le prix / la qualité / l'offre / la demande \\
\hline
Intercambios & \esp{exportar / importar / la exportación / la aduana / los aranceles / la balanza comercial / el socio comercial} & exporter / importer / l'exportation / la douane / les droits de douane / la balance commerciale / le partenaire commercial \\
\hline
Dinero & \esp{la inversión / invertir / el beneficio / la ganancia / la pérdida / el préstamo / la deuda / el éxito} & l'investissement / investir / le bénéfice / le gain / la perte / le prêt / la dette / le succès \\
\hline
\end{tabularx}
\end{center}

\begin{methode}[Comment mémoriser le lexique efficacement]
\begin{enumerate}
\item \textbf{Regroupe par famille de mots} : apprends ensemble \esp{exportar}, \esp{la exportación}, \esp{el exportador}. Un seul effort, trois mots !
\item \textbf{Fabrique une phrase personnelle par champ} : par exemple, pour le champ de l'argent : \esp{Pedí un préstamo, invertí el dinero y obtuve beneficios.} « J'ai demandé un prêt, j'ai investi l'argent et j'ai réalisé des bénéfices. » Une phrase vraie ou rêvée sur TA vie se retient dix fois mieux qu'une liste.
\item \textbf{Utilise les ressemblances avec le français} : les mots en \esp{-ción} (\esp{exportación}, \esp{inversión}) sont presque identiques au français ; méfie-toi seulement des faux amis (\esp{éxito}, \esp{competencia}).
\item \textbf{Teste-toi en cache-cache} : recouvre la colonne française du tableau et retrouve les traductions, puis fais l'inverse.
\end{enumerate}
\end{methode}

\begin{exoresolu}[Complète le mini-dictionnaire bilingue]
Énoncé : complète ce mini-dictionnaire en écrivant l'équivalent français ou espagnol de chaque mot.
\begin{enumerate}
\item \esp{el emprendedor} = \ldots
\item le marché = \ldots
\item \esp{la aduana} = \ldots
\item investir = \ldots
\item \esp{el préstamo} = \ldots
\item la concurrence = \ldots
\item \esp{obtener beneficios} = \ldots
\item le partenaire commercial = \ldots
\end{enumerate}
\tcblower
Solution détaillée :
\begin{enumerate}
\item \esp{el emprendedor} = l'entrepreneur (celui qui lance une activité).
\item le marché = \esp{el mercado}.
\item \esp{la aduana} = la douane (nom féminin : on dit bien \esp{la aduana}, au pluriel \esp{las aduanas}).
\item investir = \esp{invertir} (attention au changement de radical : \esp{invierto}, \esp{invirtió}).
\item \esp{el préstamo} = le prêt, l'emprunt.
\item la concurrence = \esp{la competencia} (faux ami : ce n'est pas « compétence » au sens de capacité dans ce contexte).
\item \esp{obtener beneficios} = réaliser des bénéfices (expression figée à apprendre en bloc).
\item le partenaire commercial = \esp{el socio comercial}.
\end{enumerate}
\end{exoresolu}

\begin{exercice}[À ton tour : empresa y comercio]
\begin{enumerate}
\item Traduis en espagnol : a) Le Cameroun exporte du café et du bois. b) Mon associé a investi ses économies dans le projet. c) La concurrence est forte, mais nos clients aiment la qualité. d) Je vais demander un prêt à la banque pour lancer mon projet.
\item Complète avec le mot de la bonne famille : a) \esp{El Camerún es un gran \ldots\ de cacao} (exportar). b) \esp{La \ldots\ de máquinas aumenta cada año} (importar). c) \esp{Amina es una \ldots\ muy valiente} (emprender).
\item Rédige trois phrases personnelles, une par champ lexical (entreprise, marché, argent), pour présenter ton projet de rêve.
\end{enumerate}
\end{exercice}

\begin{corrige}[Exercice : empresa y comercio]
\begin{enumerate}
\item a) \esp{El Camerún exporta café y madera.} b) \esp{Mi socio invirtió sus ahorros en el proyecto.} (passé simple, 3\up{e} personne : e $\rightarrow$ i) c) \esp{La competencia es fuerte, pero nuestros clientes aman la calidad.} (on peut aussi dire \esp{les gusta la calidad}) d) \esp{Voy a pedir un préstamo al banco para lanzar mi proyecto.} (\esp{para} + infinitif exprime le but.)
\item a) \esp{exportador} (nom de personne en \esp{-dor}) ; b) \esp{importación} (nom d'action en \esp{-ción}, féminin : \esp{la importación}) ; c) \esp{emprendedora} (féminin de \esp{emprendedor}).
\item Exemple de production attendue : \esp{Quiero crear una empresa de zumos naturales. En el mercado de mi barrio, ganaré muchos clientes con buena calidad. Pediré un préstamo pequeño y obtendré beneficios el primer año.}
\end{enumerate}
\end{corrige}

\begin{aretenir}[L'essentiel du lexique]
\begin{itemize}
\item Quatre champs à maîtriser : l'entreprise (\esp{emprendedor, empresa, proyecto, socio}), le marché (\esp{mercado, cliente, competencia, precio, calidad, oferta, demanda}), les échanges (\esp{exportar, importar, aduana, balanza comercial, socio comercial}) et l'argent (\esp{inversión, beneficio, pérdida, préstamo, deuda}).
\item Les noms en \esp{-ción} sont féminins et proches du français en « -tion » ; leur pluriel perd l'accent : \esp{las exportaciones}.
\item \esp{Invertir} est irrégulier : \esp{invierto}, \esp{invirtió}, \esp{invirtieron}.
\item Faux amis : \esp{el éxito} = le succès ; \esp{la competencia} = la concurrence ; \esp{el precio} = le prix (et non \esp{el premio}, la récompense).
\item Expressions clés : \esp{crear una empresa}, \esp{lanzar un proyecto}, \esp{ganar clientes}, \esp{obtener beneficios}, \esp{pedir un préstamo}.
\end{itemize}
\end{aretenir}

Maintenant que ton vocabulaire est solide, tu es prêt à passer à la grammaire : dans la prochaine section, tu apprendras à parler de ton projet au futur (\esp{voy a crear}, \esp{crearé}) et à exprimer tes souhaits au conditionnel (\esp{me gustaría}).

\coursec{Grammaire 1 : le futur du projet (futur proche et futur simple)}

Dans la section précédente, nous avons construit le vocabulaire de l'entreprise et du commerce : \esp{empresa}, \esp{mercado}, \esp{exportar}, \esp{beneficio}\ldots Mais un lexique ne suffit pas : pour \cle{présenter un projet entrepreneurial}, il faut pouvoir parler de l'\textbf{avenir}. Or l'espagnol dispose de deux outils pour cela : le \cle{futur proche} (\esp{voy a crear}) et le \cle{futur simple} (\esp{crearé}). C'est exactement ce que fait Amina, l'héroïne de notre texte support \emph{Un sueño emprendedor en Duala}. Observons-la.

\courssub{Observation : deux façons de parler de l'avenir}

\begin{exemplebox}[Deux phrases du texte support]
\esp{Voy a crear una empresa de zumos naturales.} « Je vais créer une entreprise de jus naturels. »

\esp{Crearé diez empleos para los jóvenes de mi barrio.} « Je créerai dix emplois pour les jeunes de mon quartier. »
\end{exemplebox}

Dans les deux phrases, Amina parle de l'avenir. Pourtant les formes sont différentes :

\begin{center}
\begin{tabularx}{\linewidth}{|X|X|X|}
\hline
\rowcolor{popblueL} Phrase & Structure & Nuance \\
\hline
\esp{Voy a crear} & verbe \esp{ir} au présent + \esp{a} + infinitif & projet déjà décidé, proche, concret \\
\hline
\esp{Crearé} & infinitif \esp{crear} + terminaison \esp{-é} & promesse, objectif plus lointain \\
\hline
\end{tabularx}
\end{center}

La première phrase exprime un projet \textbf{déjà mûri} : Amina a son plan, son local, ses fournisseurs. La deuxième est une \textbf{promesse}, une ambition pour plus tard. Le français fait exactement la même distinction : « je \emph{vais créer} » / « je \emph{créerai} ». Retenons cette intuition ; précisons maintenant chaque forme.

\courssub{Règle 1 : le futur proche — \esp{ir} + \esp{a} + infinitif}

\begin{propriete}[Formation du futur proche (\esp{futuro próximo})]
Le futur proche se forme avec trois éléments inséparables :

\textbf{verbe \esp{ir} conjugué au présent + \esp{a} + verbe à l'infinitif}

\begin{center}
\begin{tabular}{|l|l|l|}
\hline
\rowcolor{popgreenL} Personne & \esp{ir} au présent & Exemple \\
\hline
yo & voy & \esp{Voy a exportar.} \\
\hline
tú & vas & \esp{Vas a exportar.} \\
\hline
él / ella / usted & va & \esp{Va a exportar.} \\
\hline
nosotros/-as & vamos & \esp{Vamos a exportar.} \\
\hline
vosotros/-as & vais & \esp{Vais a exportar.} \\
\hline
ellos / ellas / ustedes & van & \esp{Van a exportar.} \\
\hline
\end{tabular}
\end{center}
\end{propriete}

\begin{propriete}[Emplois du futur proche]
On utilise \esp{ir a + infinitivo} pour :
\begin{itemize}
\item un \cle{projet décidé}, déjà préparé : \esp{Voy a abrir una tienda en el mercado.} « Je vais ouvrir un magasin au marché. »
\item une action \cle{proche dans le temps} : \esp{Mañana vamos a visitar la fábrica.} « Demain nous allons visiter l'usine. »
\item une conséquence \cle{imminente}, visible : \esp{¡Cuidado! ¡El camión va a salir!} « Attention ! Le camion va partir ! »
\end{itemize}
\end{propriete}

\begin{exemplebox}[Exemples traduits]
\esp{Voy a abrir una tienda.} « Je vais ouvrir un magasin. »

\esp{Mi hermano va a invertir en mi proyecto.} « Mon frère va investir dans mon projet. »

\esp{¿Vais a importar máquinas de España?} « Allez-vous importer des machines d'Espagne ? »

\esp{El próximo mes vamos a contratar a dos empleados.} « Le mois prochain, nous allons embaucher deux employés. »
\end{exemplebox}

\begin{attention}[Piège du francophone : la préposition \esp{a}]
En français on dit « je vais créer » sans mot entre les deux verbes. En espagnol, la préposition \esp{a} est \textbf{obligatoire} : on ne dit jamais \esp{voy crear}, mais toujours \esp{voy \textbf{a} crear}. De même, l'infinitif reste \textbf{invariable} : \esp{voy a crear}, \esp{vamos a crear} (jamais \esp{vamos a creamos} !). Enfin, on ne mélange pas les deux futurs : \esp{voy a crearé} est une faute grave.
\end{attention}

\courssub{Règle 2 : le futur simple — l'infinitif entier comme base}

\begin{propriete}[Formation du futur simple (\esp{futuro simple})]
Le futur simple espagnol se forme sur l'\textbf{infinitif ENTIER} (avec son \esp{-r} final) auquel on ajoute les mêmes terminaisons pour \textbf{les trois groupes} (-ar, -er, -ir) :

\begin{center}
\begin{tabular}{|l|l|}
\hline
\rowcolor{poporangeL} Personne & Terminaison \\
\hline
yo & \textbf{-é} \\
\hline
tú & \textbf{-ás} \\
\hline
él / ella / usted & \textbf{-á} \\
\hline
nosotros/-as & \textbf{-emos} \\
\hline
vosotros/-as & \textbf{-éis} \\
\hline
ellos / ellas / ustedes & \textbf{-án} \\
\hline
\end{tabular}
\end{center}

Ces terminaisons viennent historiquement du présent du verbe \esp{haber} (\esp{he, has, ha, hemos, habéis, han}) soudé à l'infinitif : \esp{crear + he} $\rightarrow$ \esp{crearé}. C'est une grande différence avec le français, où le radical change souvent (je \emph{fer}-ai, j'\emph{ir}-ai) : en espagnol, l'infinitif reste presque toujours intact.
\end{propriete}

\begin{center}
\begin{tabular}{|l|l|l|l|}
\hline
\rowcolor{popblueL} Personne & \esp{hablar} (parler) & \esp{comer} (manger) & \esp{vivir} (vivre) \\
\hline
yo & hablaré & comeré & viviré \\
\hline
tú & hablarás & comerás & vivirás \\
\hline
él / ella / usted & hablará & comerá & vivirá \\
\hline
nosotros/-as & hablaremos & comeremos & viviremos \\
\hline
vosotros/-as & hablaréis & comeréis & viviréis \\
\hline
ellos / ellas / ustedes & hablarán & comerán & vivirán \\
\hline
\end{tabular}
\end{center}

\begin{attention}[L'accent, signature du futur]
Toutes les formes du futur portent un \textbf{accent aigu}, sauf \esp{nosotros} : \esp{crear\textbf{é}}, \esp{exportar\textbf{ás}}, \esp{vender\textbf{á}}, \esp{crear\textbf{éis}}, \esp{importar\textbf{án}}, mais \esp{crearemos} (sans accent). Sans accent, \esp{creare} ou \esp{exportara} n'existent pas au futur — et \esp{hablara} se confond avec l'imparfait du subjonctif ! L'accent est donc porteur de sens.
\end{attention}

\courssub{Les radicaux irréguliers du futur}

Une douzaine de verbes modifient leur \textbf{radical} avant d'ajouter les terminaisons (qui, elles, restent \textbf{régulières} et accentuées). Trois familles :

\begin{itemize}
\item \textbf{Perte d'une voyelle} (le \esp{e} de l'infinitif tombe) : \esp{poder} $\rightarrow$ \esp{podr-}, \esp{saber} $\rightarrow$ \esp{sabr-}, \esp{haber} $\rightarrow$ \esp{habr-}, \esp{querer} $\rightarrow$ \esp{querr-}, \esp{caber} $\rightarrow$ \esp{cabr-}.
\item \textbf{Apparition d'un \esp{d}} : \esp{poner} $\rightarrow$ \esp{pondr-}, \esp{salir} $\rightarrow$ \esp{saldr-}, \esp{tener} $\rightarrow$ \esp{tendr-}, \esp{venir} $\rightarrow$ \esp{vendr-}, \esp{valer} $\rightarrow$ \esp{valdr-}.
\item \textbf{Verbes très irréguliers} : \esp{decir} $\rightarrow$ \esp{dir-}, \esp{hacer} $\rightarrow$ \esp{har-}.
\end{itemize}

\begin{center}
\begin{tabularx}{\linewidth}{|X|X|X|}
\hline
\rowcolor{popgreenL} Infinitif & Radical du futur & 1\textsuperscript{re} personne (yo) \\
\hline
decir (dire) & dir- & diré \\
\hline
hacer (faire) & har- & haré \\
\hline
poder (pouvoir) & podr- & podré \\
\hline
poner (mettre) & pondr- & pondré \\
\hline
querer (vouloir) & querr- & querré \\
\hline
saber (savoir) & sabr- & sabré \\
\hline
salir (sortir) & saldr- & saldré \\
\hline
tener (avoir) & tendr- & tendré \\
\hline
venir (venir) & vendr- & vendré \\
\hline
haber (avoir, auxiliaire) & habr- & habré \\
\hline
\end{tabularx}
\end{center}

\begin{popfigure}
\begin{center}
\begin{tikzpicture}[font=\small]
\node[draw=popblue, fill=popblueL, rounded corners, align=center, minimum width=2.6cm] (inf) at (0,0) {\textbf{Infinitif}\\ \esp{tener}};
\node[draw=poporange, fill=poporangeL, rounded corners, align=center, minimum width=2.6cm] (rad) at (4.6,0) {\textbf{Radical}\\ \esp{tendr-}};
\node[draw=popgreen, fill=popgreenL, rounded corners, align=center, minimum width=3.2cm] (fut) at (9.6,0) {\textbf{Futur}\\ \esp{tendré, tendrás, tendrá,}\\ \esp{tendremos, tendréis, tendrán}};
\draw[-{Stealth}, very thick, poporange] (inf) -- (rad) node[midway, above, popdark]{+ \esp{d}};
\draw[-{Stealth}, very thick, popgreen] (rad) -- (fut) node[midway, above, popdark]{+ \esp{-é, -ás\ldots}};
\end{tikzpicture}
\end{center}
\legende{De l'infinitif au futur : le radical irrégulier reçoit les terminaisons régulières.}
\end{popfigure}

\begin{attention}[Pièges des francophones]
\begin{itemize}
\item \esp{querré} s'écrit avec \textbf{deux r} (radical \esp{querr-}) ; \esp{queré} n'existe pas.
\item \esp{habré} (« j'aurai ») et surtout \esp{habrá} (« il y aura ») gardent le \esp{h} muet : jamais \esp{abré}, qui se confondrait avec \esp{abre} (« il ouvre ») !
\item \esp{haré} (je ferai) exige son accent : sans accent, la forme n'existe pas.
\item Le français dit « je \emph{ferai} » avec un radical très modifié ; l'espagnol \esp{haré} est plus simple : un seul radical irrégulier à connaître, \esp{har-}.
\item Attention au faux ami : \esp{éxito} signifie « succès », pas « exit » ; « sortie » se dit \esp{salida}.
\end{itemize}
\end{attention}

\courssub{Emplois du futur simple}

\begin{propriete}[Quand utiliser le futur simple ?]
\begin{enumerate}
\item \textbf{Projet lointain, objectif à long terme} : \esp{El año próximo abriré una tienda.} « L'année prochaine, j'ouvrirai un magasin. »
\item \textbf{Promesse, engagement solennel} : \esp{Crearé diez empleos.} « Je créerai dix emplois. »
\item \textbf{Prédiction, hypothèse sur l'avenir} : \esp{Mi empresa tendrá éxito.} « Mon entreprise aura du succès. » / \esp{El mercado crecerá.} « Le marché grandira. »
\item \textbf{Supposition sur le présent} (emploi plus avancé) : \esp{Serán las cinco.} « Il doit être cinq heures. »
\end{enumerate}
\end{propriete}

\begin{exemplebox}[Exemples traduits]
\esp{Abriré una tienda el año próximo.} « J'ouvrirai un magasin l'année prochaine. »

\esp{El próximo año, exportaremos nuestro cacao a España.} « L'année prochaine, nous exporterons notre cacao vers l'Espagne. »

\esp{Mis clientes vendrán de toda la ciudad.} « Mes clients viendront de toute la ville. »

\esp{¿Tendréis suficiente dinero para la inversión?} « Aurez-vous assez d'argent pour l'investissement ? »

\esp{Habrá más oportunidades para los jóvenes emprendedores.} « Il y aura plus d'opportunités pour les jeunes entrepreneurs. »
\end{exemplebox}

Remarque importante : \esp{habrá} est la forme \textbf{impersonnelle} de \esp{haber} au futur ; elle reste toujours au singulier : \esp{Habrá muchos clientes.} « Il y aura beaucoup de clients. » (jamais \esp{habrán muchos clientes} dans la langue soignée).

\courssub{Frise des temps : du présent au futur}

\begin{popfigure}
\begin{center}
\begin{tikzpicture}[font=\small]
\draw[-{Stealth}, very thick, popdark] (0,0) -- (12.5,0);
\foreach \x in {1,6,11}{
  \fill[popblue] (\x,0) circle (0.09);
}
\node[below, align=center, popdark] at (1,-0.15) {\textbf{Présent}\\ \esp{creo}};
\node[below, align=center, popdark] at (6,-0.15) {\textbf{Futur proche}\\ \esp{voy a crear}};
\node[below, align=center, popdark] at (11,-0.15) {\textbf{Futur simple}\\ \esp{crearé}};
\node[above, align=center, popmuted] at (1,0.15) {maintenant};
\node[above, align=center, popmuted] at (6,0.15) {projet décidé, proche};
\node[above, align=center, popmuted] at (11,0.15) {promesse, lointain};
\draw[-{Stealth}, thick, poporange] (3.2,1.3) -- (10.4,1.3) node[midway, above, poporange]{certitude et distance croissantes};
\end{tikzpicture}
\end{center}
\legende{Plus on s'éloigne du présent, plus on passe du futur proche au futur simple.}
\end{popfigure}

\courssub{Comparaison français / espagnol}

\begin{center}
\begin{tabularx}{\linewidth}{|X|X|X|}
\hline
\rowcolor{popblueL} Français & Español & Remarque \\
\hline
Je vais exporter. & \esp{Voy a exportar.} & Même structure : \esp{ir a} = « aller ». \\
\hline
J'exporterai. & \esp{Exportaré.} & L'espagnol garde l'infinitif entier + \esp{-é}. \\
\hline
Je ferai. & \esp{Haré.} & Radical irrégulier dans les deux langues. \\
\hline
Nous aurons du succès. & \esp{Tendremos éxito.} & \esp{tener} $\rightarrow$ \esp{tendr-} ; pas d'accent à \esp{nosotros}. \\
\hline
Il y aura des clients. & \esp{Habrá clientes.} & \esp{habrá} = « il y aura » (une seule forme). \\
\hline
\end{tabularx}
\end{center}

\courssub{Méthode : quel futur choisir ?}

\begin{methode}[Futur proche ou futur simple ?]
Pour présenter ton projet entrepreneurial, pose-toi deux questions :
\begin{enumerate}
\item \textbf{Le projet est-il déjà décidé, préparé, proche ?} $\rightarrow$ utilise \esp{ir a + infinitivo} : \esp{Voy a crear una empresa.} (j'ai déjà mon plan d'affaires).
\item \textbf{S'agit-il d'une promesse, d'un objectif lointain, d'une prédiction ?} $\rightarrow$ utilise le futur simple : \esp{Crearé diez empleos.} \esp{Mi empresa tendrá éxito.}
\end{enumerate}
Dans une présentation orale, alterne les deux : le futur proche pour les étapes concrètes, le futur simple pour la vision à long terme. C'est exactement la stratégie d'Amina dans notre texte.
\end{methode}

\begin{exoresolu}[Conjugue au futur demandé]
Énoncé — Mets le verbe entre parenthèses à la forme convenable (futur proche ou futur simple selon le sens) :

1. Mañana yo (visitar) el mercado central de Duala. [projet proche]
2. Mi empresa (tener) muchos clientes. [prédiction]
3. Nosotros (exportar) cacao a Guinea Ecuatorial el año próximo. [objectif lointain]
4. ¿Tú (pedir) un préstamo al banco? [projet décidé]

\tcblower
Solution détaillée :

1. Projet proche et décidé $\rightarrow$ futur proche : \esp{Mañana yo \textbf{voy a visitar} el mercado central de Duala.} « Demain je vais visiter le marché central de Douala. »
2. Prédiction $\rightarrow$ futur simple, verbe irrégulier \esp{tener} $\rightarrow$ \esp{tendr-} : \esp{Mi empresa \textbf{tendrá} muchos clientes.} « Mon entreprise aura beaucoup de clients. »
3. Objectif lointain $\rightarrow$ futur simple régulier : \esp{Nosotros \textbf{exportaremos} cacao a Guinea Ecuatorial el año próximo.} « L'année prochaine, nous exporterons du cacao vers la Guinée équatoriale. » Attention : pas d'accent à \esp{nosotros}.
4. Projet décidé $\rightarrow$ futur proche : \esp{¿Tú \textbf{vas a pedir} un préstamo al banco?} « Vas-tu demander un prêt à la banque ? » (Si l'on voulait insister sur l'engagement : \esp{¿Pedirás un préstamo?}, futur simple régulier.) Remarque lexicale : en espagnol, « faire un emprunt » se dit \esp{pedir un préstamo} ; \esp{hacer un préstamo} signifie plutôt « consentir un prêt » (c'est la banque qui le fait !).
\end{exoresolu}

\begin{exercice}[À ton tour]
1. Conjugue au futur simple : \esp{crear} (yo), \esp{vender} (ellos), \esp{salir} (tú), \esp{poder} (nosotros), \esp{decir} (usted).
2. Traduis en espagnol : « Je vais investir dans ce projet. » / « Nous créerons une entreprise. » / « Il y aura un nouveau marché à Yaoundé. »
3. Choisis la bonne forme et justifie : \esp{Voy a abrir / Abriré mi tienda mañana por la mañana.}
\end{exercice}

\begin{corrige}[Exercice « À ton tour »]
1. \esp{crearé} ; \esp{venderán} ; \esp{saldrás} (radical \esp{saldr-}) ; \esp{podremos} (radical \esp{podr-}, sans accent à \esp{nosotros}) ; \esp{dirá} (radical \esp{dir-}).
2. \esp{Voy a invertir en este proyecto.} / \esp{Crearemos una empresa.} / \esp{Habrá un nuevo mercado en Yaoundé.}
3. \esp{Voy a abrir mi tienda mañana por la mañana} : l'action est proche (\esp{mañana}) et décidée, donc futur proche. \esp{Abriré} conviendrait pour une promesse ou une échéance plus lointaine (\esp{el año próximo}).
\end{corrige}

\begin{aretenir}[L'essentiel du futur du projet]
\begin{itemize}
\item \textbf{Futur proche} = \esp{ir} (présent) + \esp{a} + infinitif : projet décidé, proche, concret. \esp{Voy a exportar.}
\item \textbf{Futur simple} = infinitif ENTIER + \esp{-é, -ás, -á, -emos, -éis, -án} (identique pour -ar, -er, -ir) : promesse, objectif lointain, prédiction. \esp{Exportaré.}
\item Radicaux irréguliers à mémoriser : \esp{dir-, har-, podr-, pondr-, querr-, sabr-, saldr-, tendr-, vendr-, habr-}.
\item Accent sur toutes les formes sauf \esp{nosotros} ; jamais de mélange du type \esp{voy a crearé}.
\item \esp{Habrá} = « il y aura », forme impersonnelle toujours au singulier.
\end{itemize}
\end{aretenir}

\coursec{Grammaire 2 : le conditionnel du souhait (me gustaría)}

Dans la section précédente, nous avons appris à présenter un projet avec le futur : \esp{voy a crear una empresa} (je vais créer une entreprise, décision proche) et \esp{crearé una empresa} (je créerai une entreprise, promesse ferme). Mais tout projet commence souvent par un rêve, un souhait encore incertain. Comment dire en espagnol « j'aimerais créer une entreprise », « je voudrais exporter mon cacao » ? C'est exactement le rôle du \cle{conditionnel}, le temps du souhait, de la politesse et de l'hypothèse. Bonne nouvelle : tu connais déjà la moitié du travail, car le conditionnel utilise \textbf{exactement les mêmes radicaux que le futur simple}.

\courssub{Observation : un projet rêvé, pas encore certain}

Lisons d'abord ce court dialogue entre deux élèves de Première à Douala, après la visite du port.

\begin{textebox}[Un sueño en el puerto]
— ¿Viste todos esos contenedores de cacao? ¡Qué impresionante!\\
— Sí. Yo \textbf{voy a estudiar} comercio internacional y \textbf{crearé} mi propia empresa de exportación.\\
— ¡Qué ambicioso! A mí \textbf{me gustaría exportar} cacao a México, pero todavía no tengo el dinero.\\
— Si tuvieras un buen socio, \textbf{podrías} empezar pronto.\\
— Tienes razón. \textbf{Me gustaría} hablar con un banco primero. ¿\textbf{Podrías} ayudarme con el proyecto?\\
— Claro. Y \textbf{deberías} visitar también la Cámara de Comercio de Duala.
\end{textebox}

\textbf{Glossaire :} \esp{contenedor} = conteneur ; \esp{socio} = associé, partenaire ; \esp{empezar} = commencer ; \esp{banco} = banque ; \esp{Cámara de Comercio} = Chambre de commerce.

\textbf{Questions de compréhension :}
\begin{enumerate}
\item ¿Qué quiere estudiar el primer alumno?
\item ¿Por qué el segundo alumno no puede exportar todavía?
\item ¿Qué consejo le da su amigo?
\end{enumerate}

Observons les verbes en gras. Certains annoncent des actions sûres (\esp{voy a estudiar}, \esp{crearé}), d'autres expriment un souhait, une possibilité ou un conseil : \esp{me gustaría exportar}, \esp{podrías empezar}, \esp{deberías visitar}. Remarque la terminaison commune : \textbf{-ía}. C'est la marque du conditionnel.

\courssub{Formation du conditionnel simple}

\begin{propriete}[Règle de formation]
Le conditionnel simple se forme sur \textbf{l'infinitif entier} du verbe (comme le futur simple), auquel on ajoute les terminaisons :
\begin{center}
-ía, -ías, -ía, -íamos, -íais, -ían
\end{center}
Ces terminaisons sont \textbf{identiques pour les trois groupes} (-ar, -er, -ir). Ce sont d'ailleurs les terminaisons de l'imparfait de l'indicatif des verbes en -er/-ir : compare \esp{crearía} (je créerais) et \esp{vivía} (je vivais).
\end{propriete}

\begin{propriete}[Futur et conditionnel : le même radical]
Futur et conditionnel partagent \textbf{exactement le même radical} ; seules les terminaisons changent. Les radicaux irréguliers du futur sont donc aussi ceux du conditionnel :
\begin{center}
\begin{tabularx}{\linewidth}{|X|X|X|X|}
\hline
\rowcolor{popblueL} Infinitivo & Radical & Futuro & Condicional \\
\hline
crear & crear- & crearé & crearía \\
\hline
decir & dir- & diré & diría \\
\hline
hacer & har- & haré & haría \\
\hline
poder & podr- & podré & podría \\
\hline
tener & tendr- & tendré & tendría \\
\hline
querer & querr- & querré & querría \\
\hline
salir & saldr- & saldré & saldría \\
\hline
haber & habr- & habrá & habría \\
\hline
\end{tabularx}
\end{center}
Retiens la formule : radical du futur \textbf{+ -é} = futur ; radical du futur \textbf{+ -ía} = conditionnel. Note l'emploi utile de \esp{habría} : \esp{habría más clientes} = « il y aurait plus de clients ».
\end{propriete}

\begin{center}
\begin{tabular}{|l|l|l|l|}
\hline
\rowcolor{popblueL} Persona & crear & poder & querer \\
\hline
yo & crearía & podría & querría \\
\hline
tú & crearías & podrías & querrías \\
\hline
él / ella / usted & crearía & podría & querría \\
\hline
nosotros & crearíamos & podríamos & querríamos \\
\hline
vosotros & crearíais & podríais & querríais \\
\hline
ellos / ustedes & crearían & podrían & querrían \\
\hline
\end{tabular}
\end{center}

\begin{attention}[Accent obligatoire]
Toutes les formes du conditionnel portent un \textbf{accent écrit sur le i} : \esp{crear\textbf{í}a}, \esp{podr\textbf{í}amos}. Sans accent, le mot est fautif : \esp{podria} et \esp{haria} n'existent pas, et l'oubli de l'accent est une erreur grave à l'examen.
\end{attention}

\courssub{Me gustaría : la structure inversée du verbe gustar}

L'expression la plus utile de cette leçon est \esp{me gustaría + infinitivo} = « j'aimerais + infinitif ». Mais attention à sa logique, très différente du français.

\begin{definition}[Le verbe gustar, un verbe inversé]
Avec \esp{gustar}, le sujet grammatical est \textbf{la chose qui plaît}, et la personne est exprimée par un pronom complément (\esp{me, te, le, nos, os, les}). Littéralement, \esp{me gustaría exportar} signifie « exporter me plairait ». On ne conjugue donc jamais \esp{gustar} à la première personne pour dire « j'aime » ou « j'aimerais ».
\end{definition}

\begin{center}
\begin{tabularx}{\linewidth}{|X|X|X|}
\hline
\rowcolor{popblueL} Français & Español & Remarque \\
\hline
J'aimerais créer une entreprise. & Me gustaría crear una empresa. & sujet = « crear » (infinitif) \\
\hline
Tu aimerais voyager. & Te gustaría viajar. & te = à toi \\
\hline
Nous aimerions exporter. & Nos gustaría exportar. & nos = à nous \\
\hline
J'aimerais un meilleur prix. & Me gustaría un precio mejor. & gustaría s'accorde avec « precio » (singulier) \\
\hline
J'aimerais ces machines. & Me gustarían estas máquinas. & gustarían (pluriel, sujet = « máquinas ») \\
\hline
\end{tabularx}
\end{center}

\begin{attention}[Erreurs fatales des francophones]
\begin{itemize}
\item Ne dis \textbf{jamais} \esp{*yo gustaría} : c'est un calque du français « j'aimerais ». Le sujet n'est pas « yo », c'est la chose. Dis : \esp{me gustaría}.
\item Après \esp{me gustaría}, mets toujours un \textbf{infinitif} quand c'est toi qui fais l'action : \esp{me gustaría exportar}. La construction \esp{me gustaría que + subjonctif imparfait} existe, mais seulement quand c'est une \textbf{autre personne} qui agit : \esp{Me gustaría que me ayudaras.} « J'aimerais que tu m'aides. »
\item Ne confonds pas \esp{me gusta} (ça me plaît, réalité présente) et \esp{me gustaría} (ça me plairait, souhait).
\end{itemize}
\end{attention}

\courssub{Les quatre emplois du conditionnel}

\begin{propriete}[Emplois du conditionnel simple]
\begin{enumerate}
\item \textbf{Le souhait poli} : \esp{me gustaría}, \esp{querría}, \esp{desearía} + infinitif. \esp{Querría abrir una tienda en Yaundé.} « Je voudrais ouvrir une boutique à Yaoundé. »
\item \textbf{Le conseil} : \esp{deberías} + infinitif = « tu devrais ». \esp{Deberías estudiar el mercado antes de invertir.} « Tu devrais étudier le marché avant d'investir. »
\item \textbf{L'hypothèse, la possibilité} : \esp{podría} + infinitif = « je pourrais ». \esp{Podríamos vender el cacao en México.} « Nous pourrions vendre le cacao au Mexique. »
\item \textbf{Le projet conditionnel} (après \esp{si}) : \esp{Si tuviera dinero, crearía una empresa.} « Si j'avais de l'argent, je créerais une entreprise. »
\end{enumerate}
\end{propriete}

\begin{exemplebox}[Exemples traduits]
\esp{Me gustaría crear una empresa de cacao.} « J'aimerais créer une entreprise de cacao. »

\esp{¿Podrías ayudarme con mi proyecto?} « Pourrais-tu m'aider avec mon projet ? »

\esp{Querría importar máquinas de España.} « Je voudrais importer des machines d'Espagne. »

\esp{Deberías hablar con tus clientes primero.} « Tu devrais parler d'abord avec tes clients. »

\esp{En tu lugar, yo invertiría en el mercado local.} « À ta place, j'investirais dans le marché local. »
\end{exemplebox}

\begin{exemplebox}[Complément pour le Bac : la phrase hypothétique]
\esp{Si tuviera más dinero, invertiría en máquinas.} « Si j'avais plus d'argent, j'investirais dans des machines. »

La structure est : \textbf{si + imparfait du subjonctif} (\esp{tuviera}, \esp{pudiera}, \esp{fuera}) \textbf{+ conditionnel}. C'est l'équivalent exact du français « si + imparfait, conditionnel ». Attention : après \esp{si}, on ne met \textbf{jamais} le conditionnel (\esp{*si tendría} est impossible). Cette structure sera étudiée en détail en Terminale ; pour l'instant, apprends-la comme un modèle figé avec \esp{tener} et \esp{poder}.
\end{exemplebox}

\courssub{L'échelle du projet : du rêve à la décision}

Les trois temps que tu connais maintenant expriment trois degrés de réalisation du projet. Observe bien la figure suivante.

\begin{popfigure}
\begin{center}
\begin{tikzpicture}[scale=1]
\draw[-{Stealth}, very thick, popmuted] (0,0) -- (12,0);
\node[below, popmuted] at (6,-0.15) {\small degré de réalisation croissant};
\node[draw, rounded corners, fill=poppinkL, align=center, minimum width=3.2cm, minimum height=1.3cm] at (1.8,1.5) {\esp{me gustaría crear}\\ \small rêvé, incertain};
\node[draw, rounded corners, fill=popgoldL, align=center, minimum width=3.2cm, minimum height=1.3cm] at (6,1.5) {\esp{crearé}\\ \small promis, prévu};
\node[draw, rounded corners, fill=popgreenL, align=center, minimum width=3.2cm, minimum height=1.3cm] at (10.2,1.5) {\esp{voy a crear}\\ \small décidé, imminent};
\fill[popink] (1.8,0) circle (2pt);
\fill[popink] (6,0) circle (2pt);
\fill[popink] (10.2,0) circle (2pt);
\draw[popline] (1.8,0) -- (1.8,0.85);
\draw[popline] (6,0) -- (6,0.85);
\draw[popline] (10.2,0) -- (10.2,0.85);
\end{tikzpicture}
\end{center}
\legende{Du souhait au projet décidé : le conditionnel exprime le rêve, le futur simple la promesse, le futur proche la décision.}
\end{popfigure}

\begin{attention}[Conditionnel ou futur : ne pas confondre]
\esp{crearía} (je créerais, sous condition) n'est pas \esp{crearé} (je créerai, c'est sûr). Une seule lettre change tout le sens : \esp{Exportaré a México} = c'est décidé ; \esp{Exportaría a México} = j'en rêve, mais ce n'est pas certain. À l'examen, vérifie toujours si la phrase exprime une certitude (futur) ou un souhait (conditionnel).
\end{attention}

\courssub{Méthode : la politesse commerciale}

\begin{methode}[Adoucir une demande professionnelle]
Dans le monde des affaires hispanophone, le tutoiement direct avec \esp{quiero} (« je veux ») paraît brutal. Pour toute demande professionnelle (lettre, entretien, négociation), applique ces trois étapes :
\begin{enumerate}
\item Repère le verbe de la demande : \esp{quiero}, \esp{necesito}, \esp{puedes}.
\item Remplace-le par sa forme au conditionnel : \esp{querría}, \esp{me gustaría}, \esp{podría usted}.
\item Ajoute éventuellement une formule de courtoisie : \esp{por favor}, \esp{si es posible}, \esp{le agradecería} (« je vous serais reconnaissant »).
\end{enumerate}
\textbf{Avant :} \esp{Quiero un precio más bajo.} « Je veux un prix plus bas. » (trop direct)

\textbf{Après :} \esp{Me gustaría un precio más bajo, por favor.} « J'aimerais un prix plus bas, s'il vous plaît. » (poli et professionnel)

\textbf{Avant :} \esp{¿Puedes enviarme el catálogo?} \textbf{Après :} \esp{¿Podría usted enviarme el catálogo?} « Pourriez-vous m'envoyer le catalogue ? »
\end{methode}

\courssub{Exercice résolu et entraînement}

\begin{exoresolu}[Mets les verbes au conditionnel]
Énoncé. Conjugue les verbes entre parenthèses au conditionnel simple, puis traduis chaque phrase en français.

1. Yo (crear) \dots\ una empresa de zumos naturales.

2. ¿Tú (poder) \dots\ ayudarme con los clientes?

3. Nosotros (hacer) \dots\ una inversión importante.

4. A mí me (gustar) \dots\ exportar a Guinea Ecuatorial.

5. Ustedes (deber) \dots\ estudiar el mercado primero.
\tcblower
Solution détaillée.

1. \esp{Yo crearía una empresa de zumos naturales.} « Je créerais une entreprise de jus naturels. » Radical régulier \esp{crear-} + \esp{-ía}.

2. \esp{¿Tú podrías ayudarme con los clientes?} « Pourrais-tu m'aider avec les clients ? » Radical irrégulier \esp{podr-} (comme au futur \esp{podré}) + \esp{-ías}.

3. \esp{Nosotros haríamos una inversión importante.} « Nous ferions un investissement important. » Radical irrégulier \esp{har-} + \esp{-íamos}, avec l'accent sur le i.

4. \esp{A mí me gustaría exportar a Guinea Ecuatorial.} « J'aimerais exporter en Guinée équatoriale. » Attention : \esp{gustar} reste à la 3e personne ; c'est \esp{exportar} qui est le sujet. Jamais \esp{*yo gustaría}.

5. \esp{Ustedes deberían estudiar el mercado primero.} « Vous devriez étudier le marché d'abord. » Radical régulier \esp{deber-} + \esp{-ían}.
\end{exoresolu}

\begin{exercice}[À toi de jouer]
1. Traduis en espagnol : a) J'aimerais importer des machines. b) Pourriez-vous me donner le prix ? c) Si j'avais plus de clients, je gagnerais plus d'argent. d) Tu devrais parler avec la banque.

2. Transforme ces demandes trop directes en demandes polies : a) \esp{Quiero hablar con el director.} b) \esp{Dame el catálogo.} c) \esp{¿Puedes bajar el precio?}

3. Choisis entre futur et conditionnel : a) \esp{Mañana (empezaré / empezaría) mi proyecto, ya está decidido.} b) \esp{(Me gustará / Me gustaría) abrir una tienda, pero no tengo dinero.}
\end{exercice}

\begin{corrige}[Exercice : le conditionnel du souhait]
1. a) \esp{Me gustaría importar máquinas.} b) \esp{¿Podría usted darme el precio?} c) \esp{Si tuviera más clientes, ganaría más dinero.} d) \esp{Deberías hablar con el banco.} (attention : \esp{banco} est masculin, faux ami du genre français).

2. a) \esp{Querría hablar con el director.} ou \esp{Me gustaría hablar con el director.} b) \esp{¿Podría usted darme el catálogo, por favor?} c) \esp{¿Podría usted bajar el precio?}

3. a) \esp{empezaré} (décision certaine : futur). b) \esp{Me gustaría} (souhait irréalisé : conditionnel ; et jamais \esp{*me gustará} pour un souhait).
\end{corrige}

\begin{aretenir}[L'essentiel de la section]
\begin{itemize}
\item Le conditionnel se forme sur l'infinitif + \textbf{-ía, -ías, -ía, -íamos, -íais, -ían}, avec les mêmes radicaux irréguliers qu'au futur (\esp{diría}, \esp{haría}, \esp{podría}, \esp{tendría}, \esp{querría}).
\item \esp{Me gustaría + infinitivo} = « j'aimerais » : le sujet est la chose, jamais \esp{*yo gustaría}.
\item Quatre emplois : souhait poli, conseil (\esp{deberías}), hypothèse (\esp{podría}), projet conditionnel (\esp{si tuviera..., crearía...}).
\item Futur = certitude (\esp{crearé}) ; conditionnel = souhait ou condition (\esp{crearía}).
\end{itemize}
\end{aretenir}

\coursec{Méthode du commentaire et commentaire modèle}

Dans la section précédente, nous avons appris à exprimer un souhait avec le conditionnel : \esp{me gustaría exportar mi cacao}. Or, à l'examen, ces formes verbales ne servent pas seulement à parler : elles deviennent des \cle{indices} que le commentaire de texte doit repérer et interpréter. Un futur simple (\esp{crearé}) signale la détermination ; un conditionnel (\esp{me gustaría}) signale un rêve encore incertain. C'est exactement ce lien entre la grammaire et le sens que nous allons maintenant exploiter dans le \cle{commentaire de texte}, exercice central de la classe de Première A et épreuve clé du Baccalauréat.

\courssub{Qu'est-ce qu'un commentaire de texte ?}

\begin{definition}[Le commentaire de texte]
Le commentaire de texte (\esp{comentario de texto}) est un exercice d'\cle{analyse} : il ne s'agit pas de raconter le texte, mais d'expliquer \textbf{comment il est écrit} (procédés : lexique, temps verbaux, chiffres, adjectifs) et \textbf{ce qu'il signifie} (le message de l'auteur, le portrait dressé, le contexte économique). On répond à deux questions : \emph{¿Qué dice el texto?} (que dit le texte ?) et surtout \emph{¿Cómo lo dice y para qué?} (comment et pourquoi le dit-il ?).
\end{definition}

\begin{attention}[Le piège n° 1 : résumer au lieu d'analyser]
L'erreur la plus fréquente des élèves francophones est de \textbf{paraphraser} le texte : « Mballa est jeune, elle vend du cacao, elle veut exporter... ». Cela ne vaut presque rien. Le correcteur attend des phrases d'analyse comme : \esp{El uso del condicional «me gustaría» muestra que el proyecto es todavía un sueño.} « L'emploi du conditionnel "j'aimerais" montre que le projet est encore un rêve. » On part toujours d'un \cle{procédé} (une forme, un mot, un chiffre) pour arriver à une \cle{interprétation} (ce que cela révèle).
\end{attention}

\courssub{Rappel du texte support : « Un sueño emprendedor en Duala »}

Relisons le texte étudié dans les sections précédentes, car le commentaire modèle s'appuiera sur lui.

\begin{textebox}[Un sueño emprendedor en Duala]
Mballa Essono tiene veinticuatro años y vive en Duala, la capital económica de Camerún. Hija de agricultores de la región del Centro, creció entre los cacaotales de su abuelo. Después de sus estudios de comercio en la universidad, decidió no buscar un empleo, sino crear su propia empresa.

En 2023, con un pequeño préstamo de su familia, fundó « Cacao Rey », una empresa que compra cacao a los campesinos de su aldea y lo transforma en chocolate artesanal. « Voy a crear diez empleos este año », afirma con orgullo. Su mercado principal es todavía local : vende sus tabletas en los mercados de Duala y Yaundé.

Pero Mballa sueña más alto. « Me gustaría exportar a México y a España, porque allí el chocolate camerunés es muy apreciado », explica. « Crearé una página web y participaré en ferias internacionales. » Sabe que el camino es difícil : la inversión es grande y los clientes extranjeros son exigentes. Sin embargo, no tiene miedo. « El que no arriesga, no gana », sonríe.

Hoy, « Cacao Rey » emplea a seis personas y su beneficio crece cada mes. Para los jóvenes de su barrio, Mballa es ya un modelo : la prueba de que el emprendimiento puede cambiar una vida... y todo un pueblo.
\end{textebox}

\textbf{Glossaire :} \esp{cacaotal} : plantation de cacao ; \esp{préstamo} : prêt ; \esp{campesino} : paysan ; \esp{tableta} : tablette (de chocolat) ; \esp{feria} : foire, salon ; \esp{exigente} : exigeant ; \esp{arriesgar} : risquer ; \esp{beneficio} : bénéfice ; \esp{barrio} : quartier.

\courssub{La méthode du commentaire en trois parties}

\begin{methode}[Les trois parties du commentaire]
\begin{enumerate}
\item \textbf{Introduction} (3-4 lignes) : présenter le texte (type, thème, personnage principal, situation) et annoncer les deux axes. Formules utiles : \esp{Este texto presenta...} « Ce texte présente... » ; \esp{En este texto, el autor describe...} « Dans ce texte, l'auteur décrit... » ; \esp{Analizaremos primero... y después...} « Nous analyserons d'abord... puis... ».
\item \textbf{Développement} (2 axes) : chaque axe = une idée directrice + des citations espagnoles traduites + l'analyse des procédés.
\item \textbf{Conclusion} (2-3 lignes) : bilan (\esp{En conclusión...} « En conclusion... ») et ouverture vers une question plus large (l'entrepreneuriat des jeunes en Afrique, le commerce Cameroun--monde hispanophone).
\end{enumerate}
\end{methode}

\begin{methode}[Comment citer correctement]
\begin{enumerate}
\item Citer le passage \textbf{en espagnol, entre guillemets}, court (3 à 8 mots).
\item \textbf{Traduire ou expliquer} la citation en français.
\item \textbf{Interpréter} : dire ce que ce procédé révèle (lexique du commerce, temps verbal, chiffre).
\end{enumerate}
Modèle de phrase : \esp{La frase «me gustaría exportar a México»} (« la phrase "j'aimerais exporter vers le Mexique" ») \esp{muestra que el proyecto es todavía un sueño} (« montre que le projet est encore un rêve »).
\end{methode}

\begin{exemplebox}[Citations analysées]
\begin{itemize}
\item \esp{«Voy a crear diez empleos este año»} « Je vais créer dix emplois cette année » : le futur proche et le chiffre précis montrent la \textbf{détermination} et le caractère concret du projet.
\item \esp{«Me gustaría exportar a México y a España»} « J'aimerais exporter vers le Mexique et l'Espagne » : le conditionnel marque un \textbf{souhait}, une perspective encore lointaine.
\item \esp{«su beneficio crece cada mes»} « son bénéfice augmente chaque mois » : le présent d'habitude et le lexique économique prouvent la \textbf{réussite} déjà réelle.
\end{itemize}
\end{exemplebox}

\begin{popfigure}
\begin{tikzpicture}[
node distance=0.55cm,
etape/.style={rectangle, rounded corners, draw=popblue, fill=popblueL, text width=9.5cm, align=center, font=\small, inner sep=6pt},
fleche/.style={-{Stealth[length=3mm]}, thick, poporange}]
\node[etape] (intro) {\textbf{Introduction} : présenter le texte + annoncer les 2 axes};
\node[etape, below=of intro] (axe1) {\textbf{Axe 1} : le portrait de l'entrepreneure (qui ? parcours ? moyens ?) + citations traduites};
\node[etape, below=of axe1] (axe2) {\textbf{Axe 2} : le projet et ses perspectives (futur, conditionnel, marchés visés) + citations traduites};
\node[etape, below=of axe2] (concl) {\textbf{Conclusion} : bilan + ouverture (les jeunes et l'entrepreneuriat en Afrique)};
\draw[fleche] (intro) -- (axe1);
\draw[fleche] (axe1) -- (axe2);
\draw[fleche] (axe2) -- (concl);
\end{tikzpicture}
\legende{Organigramme du commentaire de texte : de l'introduction à l'ouverture.}
\end{popfigure}

\courssub{Les deux axes appliqués au texte de Mballa}

\textbf{Axe 1 : le portrait de l'entrepreneure.} On relève : son identité (\esp{tiene veinticuatro años} « elle a vingt-quatre ans »), son origine modeste (\esp{hija de agricultores} « fille d'agriculteurs »), son parcours (\esp{estudios de comercio} « études de commerce »), ses moyens (\esp{un pequeño préstamo de su familia} « un petit prêt de sa famille »). Le texte construit le portrait d'une jeune femme courageuse qui part de peu.

\textbf{Axe 2 : le projet et ses perspectives.} On relève : les temps verbaux (futur proche \esp{voy a crear}, futur simple \esp{crearé}, conditionnel \esp{me gustaría}), les chiffres (\esp{diez empleos} « dix emplois », \esp{seis personas} « six personnes »), les marchés visés (local : Duala et Yaundé ; international : México, España). Le texte montre un projet en marche, entre réalité et rêve.

\begin{center}
\begin{tabularx}{\linewidth}{|X|X|X|}
\hline
\rowcolor{popblueL} Procédé relevé & Citation (español) & Interprétation (français) \\
\hline
Lexique du commerce & \esp{empresa, mercado, beneficio, inversión} & Champ lexical économique : le texte parle d'argent et d'échanges. \\
\hline
Futur proche & \esp{voy a crear diez empleos} & Projet proche et concret : détermination. \\
\hline
Futur simple & \esp{crearé una página web} & Promesse ferme, vision d'avenir. \\
\hline
Conditionnel & \esp{me gustaría exportar} & Souhait, rêve encore incertain. \\
\hline
Chiffres & \esp{veinticuatro años, seis personas} & Précision réaliste : le succès est mesurable. \\
\hline
Proverbe & \esp{el que no arriesga, no gana} & « Qui ne risque rien ne gagne rien » : philosophie de l'entrepreneur. \\
\hline
\end{tabularx}
\end{center}

\courssub{Commentaire modèle rédigé}

\begin{textebox}[Comentario modelo]
Este texto presenta el retrato de una joven emprendedora camerunesa, Mballa Essono, que ha creado en Duala una empresa de chocolate artesanal. El autor describe su recorrido y su proyecto ; analizaremos primero el retrato de la emprendedora y después las perspectivas de su empresa.

En primer lugar, el texto dibuja el retrato de una mujer joven y valiente. Tiene solo « veinticuatro años » (« vingt-quatre ans ») y es « hija de agricultores » (« fille d'agriculteurs ») : su origen es modesto. Sin embargo, después de sus estudios, « decidió no buscar un empleo, sino crear su propia empresa » (« elle a décidé de ne pas chercher un emploi, mais de créer sa propre entreprise »). Este verbo, « decidió », en pretérito indefinido, muestra una decisión firme y definitiva.

En segundo lugar, el proyecto mezcla realidad y sueño. La realidad : el futuro próximo « voy a crear diez empleos » (« je vais créer dix emplois ») y las cifras precisas (« seis personas », « six personnes ») prueban que la empresa ya funciona. El sueño : el condicional « me gustaría exportar a México y a España » (« j'aimerais exporter vers le Mexique et l'Espagne ») revela una ambición todavía incierta, dirigida hacia el mundo hispanohablante.

En conclusión, el texto celebra el emprendimiento juvenil en Camerún. Mballa es « un modelo » para los jóvenes de su barrio. ¿No es este el futuro económico de toda África?
\end{textebox}

\textbf{Traduction du commentaire modèle :} « Ce texte présente le portrait d'une jeune entrepreneure camerounaise, Mballa Essono, qui a créé à Douala une entreprise de chocolat artisanal. L'auteur décrit son parcours et son projet ; nous analyserons d'abord le portrait de l'entrepreneure, puis les perspectives de son entreprise. — Premièrement, le texte dessine le portrait d'une femme jeune et courageuse. Elle n'a que "vingt-quatre ans" et elle est "fille d'agriculteurs" : son origine est modeste. Cependant, après ses études, elle "a décidé de ne pas chercher un emploi, mais de créer sa propre entreprise". Ce verbe, "decidió", au passé simple, montre une décision ferme et définitive. — Deuxièmement, le projet mêle réalité et rêve. La réalité : le futur proche "je vais créer dix emplois" et les chiffres précis ("six personnes") prouvent que l'entreprise fonctionne déjà. Le rêve : le conditionnel "j'aimerais exporter vers le Mexique et l'Espagne" révèle une ambition encore incertaine, tournée vers le monde hispanophone. — En conclusion, le texte célèbre l'entrepreneuriat des jeunes au Cameroun. Mballa est "un modèle" pour les jeunes de son quartier. N'est-ce pas là l'avenir économique de toute l'Afrique ? »

\textbf{Remarque :} ce commentaire compte environ 14 lignes, cite cinq passages traduits et ne résume jamais le texte : chaque citation sert une analyse. Notez aussi l'emploi des connecteurs \esp{en primer lugar} « premièrement », \esp{en segundo lugar} « deuxièmement », \esp{sin embargo} « cependant », \esp{en conclusión} « en conclusion » : ils structurent la copie et guident le correcteur.

\begin{exoresolu}[Rédiger l'introduction d'un commentaire]
Énoncé : voici un nouveau texte. \esp{« Kwame Bekono, treinta años, ha creado en Yaundé una empresa de transporte en moto, "Bendskin Express". Empezó con dos motos ; hoy tiene quince. "Me gustaría importar motos eléctricas de China", dice. "Voy a contratar a veinte jóvenes el próximo año." »} Rédige en espagnol l'introduction d'un commentaire de ce texte (3-4 lignes : présentation + annonce des deux axes).
\tcblower
Solution détaillée : on présente d'abord le texte (qui ? quoi ? où ?), puis on annonce deux axes. \esp{Este texto presenta el retrato de un joven emprendedor camerunés, Kwame Bekono, que ha creado en Yaundé una empresa de transporte en moto. El autor muestra su éxito y sus proyectos. Analizaremos primero el recorrido del emprendedor, desde dos motos hasta quince, y después sus perspectivas de futuro : la importación y la contratación de jóvenes.} Traduction : « Ce texte présente le portrait d'un jeune entrepreneur camerounais, Kwame Bekono, qui a créé à Yaoundé une entreprise de transport en moto. L'auteur montre son succès et ses projets. Nous analyserons d'abord le parcours de l'entrepreneur, de deux motos à quinze, puis ses perspectives d'avenir : l'importation et l'embauche de jeunes. » Remarque grammaticale : \esp{ha creado} est un \cle{prétérit parfait} (\esp{pretérito perfecto}) : il présente l'action passée comme un résultat visible aujourd'hui — c'est le temps du bilan, idéal pour ouvrir un commentaire.
\end{exoresolu}

\begin{attention}[Erreurs fréquentes des francophones au commentaire]
\begin{itemize}
\item \textbf{Calque de l'ordre des mots} : ne pas écrire \esp{una empresa de chocolate artesanal en Duala creada} ; dire \esp{una empresa creada en Duala} ou \esp{una empresa de chocolate artesanal}.
\item \textbf{Faux amis} : \esp{actualmente} signifie « actuellement, en ce moment » (et non « en fait, effectivement », qui se dit \esp{en realidad}) ; \esp{éxito} signifie « succès » (la « sortie » se dit \esp{salida}) ; \esp{una empresa} est une « entreprise » (et non une « empreinte »).
\item \textbf{Oublier de traduire les citations} : une citation espagnole non traduite et non analysée ne rapporte aucun point.
\item \textbf{Accents} : \esp{análisis, económico, creación, inversión, página} prennent un accent écrit ; ne pas oublier le point d'interrogation ouvrant : \esp{¿No es este el futuro?}
\item \textbf{Anglicisme} : « succès » ne se dit jamais \esp{suceso} (qui signifie « événement, fait divers »).
\end{itemize}
\end{attention}

\begin{exercice}[À ton tour]
Rédige en espagnol, en 10 à 12 lignes, le développement (axe 1 + axe 2) d'un commentaire du texte sur Kwame Bekono donné dans l'exercice résolu. Cite au moins quatre passages entre guillemets, traduis-les en français et analyse les temps verbaux (\esp{ha creado}, \esp{empezó}, \esp{me gustaría}, \esp{voy a contratar}).
\end{exercice}

\begin{corrige}[Exercice]
Éléments de réponse attendus : \textbf{Axe 1} : le parcours de Kwame — \esp{« Empezó con dos motos ; hoy tiene quince »} (« il a commencé avec deux motos ; aujourd'hui il en a quinze ») : le contraste des chiffres et le passé simple \esp{empezó} montrent une réussite rapide ; \esp{« ha creado en Yaundé una empresa »} (« il a créé à Yaoundé une entreprise ») : le prétérit parfait souligne un résultat présent. \textbf{Axe 2} : les perspectives — \esp{« Me gustaría importar motos eléctricas de China »} (« j'aimerais importer des motos électriques de Chine ») : le conditionnel exprime un souhait tourné vers le commerce international ; \esp{« Voy a contratar a veinte jóvenes »} (« je vais embaucher vingt jeunes ») : le futur proche montre un projet concret et utile à la société. Le commentaire doit se terminer par une phrase de bilan, par exemple : \esp{Así, el texto muestra que el emprendimiento crea empleo y esperanza en Camerún.} (« Ainsi, le texte montre que l'entrepreneuriat crée de l'emploi et de l'espoir au Cameroun. »)
\end{corrige}

\begin{aretenir}[L'essentiel de la section]
Un commentaire de texte = introduction (présentation + annonce des axes) + deux axes (idée + citation espagnole traduite + analyse du procédé) + conclusion (bilan + ouverture). On ne résume jamais : on part d'un procédé (temps verbal, mot du lexique commercial, chiffre) pour dégager un sens. Le futur proche signale le concret, le futur simple la promesse, le conditionnel le rêve, le prétérit parfait le bilan présent. Les connecteurs (\esp{en primer lugar}, \esp{sin embargo}, \esp{en conclusión}) structurent la copie.
\end{aretenir}

\coursec{Production guidée et entraînement à la traduction}

Après avoir appris à lire et à commenter un texte sur l'entrepreneuriat, il est temps de passer à l'\cle{expression} : c'est toi maintenant qui vas produire en espagnol. Cette dernière étape de la leçon te prépare directement aux épreuves de fin d'année et, déjà, au Bac : expression écrite guidée, thème, version et expression orale. Nous allons procéder pas à pas, en réutilisant tout ce que nous avons construit dans les sections précédentes : le lexique de l'entreprise (\esp{empresa, mercado, cliente, beneficio, inversión}) et les trois temps du projet (futur proche \esp{voy a + infinitivo}, futur simple \esp{crearé}, conditionnel \esp{me gustaría}).

\courssub{Production écrite guidée : « Mi proyecto empresarial »}

\textbf{Consigne.} Rédige un texte de 80 à 100 mots en espagnol pour présenter ton projet d'entreprise. Tu dois respecter le plan imposé en cinq étapes :

\begin{enumerate}
\item \textbf{Idea} : présente ton idée d'entreprise (présent : \esp{Mi proyecto consiste en...}).
\item \textbf{Inversión} : dis ce qu'il faut pour démarrer (argent, matériel, local).
\item \textbf{Mercado y clientes} : à qui vas-tu vendre ? Au Cameroun ? À l'étranger ?
\item \textbf{Futuro} : annonce tes actions avec \esp{voy a + infinitivo} et le futur simple (\esp{crearé, exportaré, venderé}).
\item \textbf{Sueño} : termine par ton rêve avec le conditionnel (\esp{Me gustaría...}).
\end{enumerate}

\begin{popfigure}
\begin{tikzpicture}[
  pasoS/.style={rectangle, rounded corners, draw=popblue, fill=popblueL, text=popdark, minimum width=3.4cm, minimum height=0.9cm, align=center, font=\small\bfseries},
  arr/.style={-{Stealth[length=2.5mm]}, thick, poporange}
]
\node[pasoS, align=center] (s1) at (0,0){1. Idea\\ \normalfont\itshape Mi proyecto consiste en...};
\node[pasoS, fill=poporangeL, draw=poporange, align=center] (s2) at (4.2,-1.3){2. Inversión\\ \normalfont\itshape Necesito...};
\node[pasoS, fill=popgreenL, draw=popgreen, align=center] (s3) at (8.4,-2.6){3. Mercado\\ \normalfont\itshape Mis clientes son...};
\node[pasoS, fill=poppurpleL, draw=poppurple, align=center] (s4) at (4.2,-3.9){4. Futuro\\ \normalfont\itshape Voy a crear... / Crearé...};
\node[pasoS, fill=poppinkL, draw=poppink, align=center] (s5) at (0,-5.2){5. Sueño\\ \normalfont\itshape Me gustaría...};
\draw[arr] (s1.east) -- (s2.west);
\draw[arr] (s2.east) -- (s3.west);
\draw[arr] (s3.west) -- (s4.east);
\draw[arr] (s4.west) -- (s5.east);
\end{tikzpicture}
\legende{Le plan en escalier de la production écrite : cinq étapes, du présent de l'idée au conditionnel du rêve.}
\end{popfigure}

\begin{methode}[La méthode des cinq verbes]
Pour ne jamais bloquer devant la page blanche, procède ainsi :
\begin{enumerate}
\item Écris d'abord \textbf{cinq verbes conjugués au bon temps}, un par étape du plan : par exemple \esp{consiste} (présent), \esp{necesito} (présent), \esp{tengo} (présent), \esp{voy a crear / exportaré} (futur), \esp{me gustaría} (conditionnel).
\item Construis ensuite \textbf{une phrase autour de chaque verbe}, en ajoutant un complément du lexique de la leçon : \esp{una empresa, el mercado, los clientes, el beneficio}.
\item Relie les phrases avec des connecteurs simples : \esp{primero, además, también, finalmente}.
\item Relis en vérifiant trois points : les \textbf{accents du futur} (\esp{crearé, venderá}), la construction de \esp{gustar} (\esp{me gustaría}), et la majuscule en début de phrase.
\end{enumerate}
\end{methode}

\begin{exemplebox}[Banque de phrases-modèles]
Voici des phrases toutes faites à adapter à ton projet :
\begin{itemize}
\item \esp{Mi proyecto consiste en vender jugo natural de mango.} « Mon projet consiste à vendre du jus naturel de mangue. »
\item \esp{Necesito una inversión pequeña para comprar una máquina.} « J'ai besoin d'un petit investissement pour acheter une machine. »
\item \esp{Mis clientes serán los estudiantes y los trabajadores de Duala.} « Mes clients seront les élèves et les travailleurs de Douala. »
\item \esp{Voy a crear una empresa con dos amigos.} « Je vais créer une entreprise avec deux amis. »
\item \esp{Exportaré mi café a España y a Guinea Ecuatorial.} « J'exporterai mon café en Espagne et en Guinée équatoriale. »
\item \esp{El beneficio será importante para mi familia.} « Le bénéfice sera important pour ma famille. »
\item \esp{Me gustaría tener clientes en todo el mundo hispanohablante.} « J'aimerais avoir des clients dans tout le monde hispanophone. »
\end{itemize}
\end{exemplebox}

\begin{attention}[Les pièges de la production]
\begin{itemize}
\item « J'aimerais » se traduit toujours \esp{me gustaría}, jamais \esp{yo amaría} (calque du français qui ne se dit pas).
\item « Il y aura » se traduit \esp{habrá} (futur de \esp{haber}), jamais \esp{será hay} ni \esp{haberá}.
\item N'oublie pas l'accent du futur à la 1\iere{} et à la 3\ieme{} personne du singulier : \esp{crearé, exportaré, venderá, tendrá}. En revanche, pas d'accent à « nosotros » : \esp{crearemos, venderemos}.
\item Faux amis : \esp{una empresa} = « une entreprise », pas « une empreinte » ; \esp{el beneficio} = « le bénéfice, le profit » ; \esp{el mango} = « la mangue », alors que \esp{la manga} signifie « la manche » (d'un vêtement) !
\item « Jus de fruit » se dit \esp{jugo} en Amérique latine et \esp{zumo} en Espagne : les deux sont corrects.
\end{itemize}
\end{attention}

\courssub{Modèle de production corrigé}

\begin{textebox}[Mi proyecto empresarial — modèle (95 mots)]
Mi proyecto consiste en cultivar y vender café de calidad en la región del Oeste de Camerún. Necesito una inversión pequeña para comprar las semillas y las herramientas. Mis clientes serán primero las tiendas de Bafoussam y de Yaundé, y después los mercados extranjeros. Voy a crear una empresa con mi hermano mayor y exportaré mi café a España y a Guinea Ecuatorial. El beneficio será importante para mi familia y para mi pueblo. Me gustaría tener clientes en todo el mundo hispanohablante y me gustaría también dar trabajo a los jóvenes de mi región. ¡Este es mi sueño!
\end{textebox}

\textbf{Glossaire} : \esp{las semillas} = les graines ; \esp{las herramientas} = les outils ; \esp{las tiendas} = les magasins ; \esp{extranjeros} = étrangers ; \esp{dar trabajo} = donner du travail, employer.

Observe bien ce modèle : il contient \textbf{quatre temps} (présent, futur proche, futur simple, conditionnel), cinq mots du lexique de la leçon, et il respecte exactement le plan en escalier. C'est le niveau visé à l'examen.

\courssub{Thème : du français vers l'espagnol}

Le thème entraîne la traduction précise. Voici cinq phrases sur le commerce camerounais ; traduis-les, puis compare avec la correction.

\begin{exoresolu}[Thème : le commerce camerounais]
Énoncé. Traduis en espagnol :
\begin{enumerate}
\item Je vais créer une entreprise et j'exporterai mon café.
\item Le Cameroun exporte du cacao et importe des machines.
\item Mes clients seront les jeunes de Yaoundé.
\item J'aimerais vendre mes produits en Espagne.
\item Il y aura un grand marché à Douala demain.
\end{enumerate}
\tcblower
Solution détaillée :
\begin{enumerate}
\item \esp{Voy a crear una empresa y exportaré mi café.} — Futur proche (\esp{voy a crear}) puis futur simple avec accent (\esp{exportaré}).
\item \esp{Camerún exporta cacao e importa máquinas.} — Présent de vérité générale ; attention à l'accent de \esp{Camerún} et de \esp{máquinas}.
\item \esp{Mis clientes serán los jóvenes de Yaundé.} — Futur simple de \esp{ser}, 3\ieme{} personne du pluriel : \esp{serán}.
\item \esp{Me gustaría vender mis productos en España.} — « J'aimerais » = \esp{me gustaría} + infinitif ; jamais \esp{amaría}.
\item \esp{Mañana habrá un gran mercado en Duala.} — « Il y aura » = \esp{habrá}, futur irrégulier de \esp{haber} ; \esp{gran} est la forme apocopée de \esp{grande} devant un nom masculin singulier.
\end{enumerate}
\end{exoresolu}

\begin{exercice}[À ton tour]
Traduis en espagnol :
\begin{enumerate}
\item Mon frère va importer des téléphones de Chine.
\item Nous vendrons nos mangues au marché de Douala.
\item J'aimerais avoir une petite entreprise de jus naturel.
\item L'entreprise aura beaucoup de clientes l'année prochaine.
\item Le bénéfice sera important pour ma famille.
\end{enumerate}
\end{exercice}

\begin{corrige}[Exercice : thème]
\begin{enumerate}
\item \esp{Mi hermano va a importar teléfonos de China.}
\item \esp{Venderemos nuestros mangos en el mercado de Duala.} (\esp{venderemos} : futur de \esp{vender}, 1\iere{} personne du pluriel, sans accent à « nosotros » ; « mangue » = \esp{el mango}, attention au faux ami \esp{la manga} = « la manche ».)
\item \esp{Me gustaría tener una pequeña empresa de jugo natural.}
\item \esp{La empresa tendrá muchas clientas el año próximo.} (\esp{tendrá} : futur irrégulier de \esp{tener}.)
\item \esp{El beneficio será importante para mi familia.}
\end{enumerate}
\end{corrige}

\courssub{Version : de l'espagnol vers le français}

Pour la version, lis d'abord tout le paragraphe pour comprendre le sens global, repère les temps, puis traduis phrase par phrase sans calquer la construction espagnole.

\begin{textebox}[Las exportaciones de España]
España importa cacao de África y exporta productos industriales a todo el mundo. El país también vende aceite de oliva, vino y frutas a muchos países europeos. El año próximo, las empresas españolas exportarán más productos a América Latina y a Guinea Ecuatorial. Muchos empresarios jóvenes quieren abrir nuevos mercados en África. Les gustaría tener más clientes en Camerún y en otros países africanos.
\end{textebox}

\textbf{Glossaire} : \esp{el aceite de oliva} = l'huile d'olive ; \esp{el vino} = le vin ; \esp{abrir nuevos mercados} = ouvrir de nouveaux marchés ; \esp{los empresarios} = les chefs d'entreprise.

\begin{exemplebox}[Version modèle]
\esp{España importa cacao de África y exporta productos industriales.} se traduit : « L'Espagne importe du cacao d'Afrique et exporte des produits industriels. »

Traduction complète du paragraphe : « L'Espagne importe du cacao d'Afrique et exporte des produits industriels dans le monde entier. Le pays vend aussi de l'huile d'olive, du vin et des fruits à de nombreux pays européens. L'année prochaine, les entreprises espagnoles exporteront davantage de produits vers l'Amérique latine et la Guinée équatoriale. Beaucoup de jeunes chefs d'entreprise veulent ouvrir de nouveaux marchés en Afrique. Ils aimeraient avoir plus de clients au Cameroun et dans d'autres pays africains. »
\end{exemplebox}

Remarque la traduction de \esp{les gustaría} : « ils aimeraient ». Le pronom \esp{les} espagnol devient le sujet « ils » en français, car \esp{gustar} fonctionne à l'envers du verbe français « aimer » : littéralement, « avoir plus de clients leur plairait ».

\courssub{Correction collective : la grille d'erreurs types}

En classe, nous corrigeons ensemble quelques productions d'élèves (anonymisées) avec cette grille. Apprends à t'auto-corriger avec les mêmes critères :

\begin{center}
\begin{tabularx}{\linewidth}{|X|X|X|}
\hline
\rowcolor{popblueL} \textbf{Erreur type} & \textbf{Correction} & \textbf{Règle} \\
\hline
\esp{yo amaría tener...} & \esp{me gustaría tener...} & « j'aimerais » = \esp{me gustaría} \\
\hline
\esp{creare / exportare} & \esp{crearé / exportaré} & accent obligatoire au futur (yo, él) \\
\hline
\esp{será hay} & \esp{habrá} & « il y aura » = futur de \esp{haber} \\
\hline
\esp{voy a creo} & \esp{voy a crear} & après \esp{voy a}, toujours l'infinitif \\
\hline
\esp{me gustaria} (sans accent) & \esp{me gustaría} & accent sur le \emph{í} du conditionnel \\
\hline
\esp{una empreza} & \esp{una empresa} & orthographe : \emph{s}, pas \emph{z} \\
\hline
\esp{exportaré mi café en España} & \esp{exportaré mi café a España} & direction : \esp{exportar a + pays} \\
\hline
\esp{nuestras mangas} & \esp{nuestros mangos} & « mangue » = \esp{el mango} (masculin) \\
\hline
\end{tabularx}
\end{center}

\begin{attention}[Bilan des pièges]
Trois fautes reviennent chez presque tous les francophones : (1) l'oubli des \textbf{accents du futur} (\esp{crearé}, pas \esp{creare}) ; (2) le calque \esp{amaría} pour « j'aimerais » ; (3) l'oubli de l'accent du conditionnel \esp{gustaría}. Vérifie ces trois points avant de rendre ta copie : c'est souvent là que se jouent deux ou trois points à l'examen.
\end{attention}

\courssub{Présentation orale du projet en une minute}

\begin{experience}[Mi proyecto en un minuto]
À tour de rôle, présente ton projet devant la classe en une minute, sans lire intégralement ton texte.
\begin{enumerate}
\item Prépare une petite fiche avec seulement cinq mots-clés : \esp{idea, inversión, clientes, futuro, sueño}.
\item Commence par saluer : \esp{Buenos días, voy a presentar mi proyecto empresarial.}
\item Suis le plan en escalier, en regardant la classe.
\item Termine par : \esp{Este es mi sueño. ¿Hay preguntas?} « C'est mon rêve. Y a-t-il des questions ? »
\item La classe pose deux questions ; tu réponds avec des phrases simples au présent ou au futur.
\end{enumerate}
Critères d'évaluation : prononciation claire (le \emph{j} se prononce « r » guttural, le \emph{c} devant \emph{e, i} se prononce « s » en Amérique latine et « th » en Espagne), au moins un futur proche, un futur simple et un conditionnel, et le respect du temps d'une minute.
\end{experience}

\begin{savaistu}[Le secret des copies d'excellence au Bac]
Au baccalauréat camerounais, l'expression écrite sur la vie économique exige, pour viser l'excellence, au moins \textbf{deux temps de projet} : le futur proche (\esp{voy a crear}) et le conditionnel (\esp{me gustaría}). Les correcteurs repèrent immédiatement un candidat qui sait passer du présent de l'idée au futur de l'action, puis au conditionnel du rêve. Et sache que le Cameroun entretient des échanges commerciaux réels avec le monde hispanophone : le cacao camerounais part notamment vers l'Espagne, tandis que la Guinée équatoriale, notre voisine hispanophone, est un partenaire commercial naturel de la région. Parler de ces échanges dans ta copie montre ta culture économique — les examinateurs adorent !
\end{savaistu}

\begin{exercice}[Devoir à la maison]
Rédige la version définitive de \esp{Mi proyecto empresarial} (80 à 100 mots), en corrigeant ton brouillon avec la grille d'erreurs. Souligne en rouge tes deux temps de projet et encadre cinq mots du lexique de la leçon. Tu présenteras ton texte oralement à la séance suivante.
\end{exercice}

\newpage

\coursec{Activité d'intégration}

\courssub{Objectifs de l'activité}

À la fin de cette activité d'intégration, l'élève sera capable de :

\begin{itemize}
\item \cle{comprendre} un appel à projets entrepreneurial rédigé en espagnol et en dégager les informations essentielles (qui, quoi, combien, où) ;
\item \cle{mobiliser} le lexique de l'entreprise et du commerce étudié dans la leçon : \esp{emprendedor}, \esp{empresa}, \esp{proyecto}, \esp{mercado}, \esp{cliente}, \esp{inversión}, \esp{beneficio}, \esp{exportar}, \esp{importar} ;
\item \cle{concevoir} un mini-projet entrepreneurial réaliste, lié à une richesse du Cameroun (cacao, café, manioc, tourisme, artisanat) ;
\item \cle{employer correctement} le futur proche (\esp{voy a + infinitivo}), le futur simple (\esp{crearé}, \esp{venderemos}) et le conditionnel de politesse et de souhait (\esp{me gustaría}, \esp{necesitaríamos}) ;
\item \cle{présenter à l'oral}, en deux minutes, un projet devant un public jouant le rôle d'investisseurs, et répondre à des questions simples ;
\item \cle{justifier un choix} à l'oral lors du vote final (\esp{Me gustaría invertir en... porque...}).
\end{itemize}

\courssub{Situation}

La mairie de Douala, en partenariat avec l'ambassade d'Espagne et le lycée, organise la première \esp{Feria de Jóvenes Emprendedores} (Foire des jeunes entrepreneurs). Les élèves de Première A sont invités à présenter, en espagnol, un projet d'entreprise fondé sur une richesse du Cameroun. Un jury d'« investisseurs » — les autres groupes de la classe — écoutera chaque présentation, posera des questions et votera pour le meilleur projet. Le groupe gagnant verra son affiche exposée au CDI du lycée.

Voici l'appel à projets publié dans le journal du lycée :

\begin{textebox}[Convocatoria : Feria de Jóvenes Emprendedores de Duala]
\textbf{CONVOCATORIA}\par
El Ayuntamiento de Duala y el Liceo Bilingüe organizan la primera Feria de Jóvenes Emprendedores. El tema de este año es : « Mi proyecto para el Camerún ».\par
Cada grupo de tres o cuatro estudiantes debe presentar un proyecto de empresa basado en una riqueza del país : el cacao, el café, la mandioca, el turismo o la artesanía.\par
La presentación debe explicar :\par
-- la idea del proyecto y el nombre de la empresa ;\par
-- la inversión necesaria (¿cuánto dinero necesitan?) ;\par
-- el mercado y los clientes (¿a quiénes van a vender?) ;\par
-- los beneficios esperados y los países de exportación.\par
Cada grupo tendrá dos minutos para hablar. Después, los « inversores » harán preguntas. Al final, todos votarán por el mejor proyecto.\par
¡Anímense, jóvenes emprendedores! El futuro del Camerún está en sus manos.
\end{textebox}

\textbf{Traduction.} \emph{Appel à candidatures : la mairie de Douala et le Lycée Bilingue organisent la première Foire des jeunes entrepreneurs. Le thème de cette année est : « Mon projet pour le Cameroun ». Chaque groupe de trois ou quatre élèves doit présenter un projet d'entreprise fondé sur une richesse du pays : le cacao, le café, le manioc, le tourisme ou l'artisanat. La présentation doit expliquer : l'idée du projet et le nom de l'entreprise ; l'investissement nécessaire (de combien d'argent ont-ils besoin ?) ; le marché et les clients (à qui vont-ils vendre ?) ; les bénéfices espérés et les pays d'exportation. Chaque groupe disposera de deux minutes pour parler. Ensuite, les « investisseurs » poseront des questions. À la fin, tous voteront pour le meilleur projet. Lancez-vous, jeunes entrepreneurs ! L'avenir du Cameroun est entre vos mains.}

\begin{definition}[Glossaire utile]
\esp{la convocatoria} : l'appel à candidatures ; \esp{el Ayuntamiento} : la mairie ; \esp{la riqueza} : la richesse ; \esp{la mandioca} : le manioc ; \esp{la artesanía} : l'artisanat ; \esp{los beneficios esperados} : les bénéfices espérés ; \esp{anímense} : lancez-vous (impératif de \esp{animarse}) ; \esp{está en sus manos} : est entre vos mains.
\end{definition}

\courssub{Consignes}

\textbf{Tâche 1 — Comprendre l'appel à projets (10 min).} Lis la \esp{convocatoria} ci-dessus en silence, puis réponds par écrit, en espagnol et en phrases courtes :
\begin{enumerate}
\item ¿Quiénes organizan la feria ?
\item ¿Cuál es el tema de la feria ?
\item ¿Qué riquezas del Camerún se pueden utilizar ?
\item ¿Cuánto tiempo tiene cada grupo para presentar su proyecto ?
\end{enumerate}

\textbf{Tâche 2 — Choisir une idée et compléter la fiche-projet (20 min).} Par groupes de 3 ou 4, choisissez une richesse du Cameroun et complétez la fiche suivante. Rédigez des phrases complètes en respectant les quotas grammaticaux : \textbf{au moins deux phrases au futur proche} (\esp{vamos a + infinitivo}), \textbf{deux au futur simple} (\esp{crearemos}, \esp{venderemos}, \esp{exportaremos}) et \textbf{deux au conditionnel} (\esp{nos gustaría}, \esp{necesitaríamos}).

\begin{center}
\begin{tabularx}{\linewidth}{|l|X|}
\hline
\rowcolor{popblueL} \textbf{Rubrique} & \textbf{Contenu à rédiger} \\
\hline
\esp{Nombre de la empresa} & Un nom court et original \\
\hline
\esp{La idea} & 2 phrases (futur proche) \\
\hline
\esp{La inversión} & Combien d'argent ? Pour quoi ? (conditionnel) \\
\hline
\esp{El mercado y los clientes} & À qui vendre ? Où ? (futur simple) \\
\hline
\esp{Los beneficios y la exportación} & Quels gains ? Quels pays ? (futur simple + conditionnel) \\
\hline
\end{tabularx}
\end{center}

\textbf{Tâche 3 — Créer l'affiche (15 min).} Sur une feuille, réalisez une affiche avec : le nom de l'entreprise, un slogan en espagnol (ex. \esp{¡El cacao de Duala para el mundo!}), un dessin simple et les cinq rubriques de la fiche-projet en mots-clés.

\textbf{Tâche 4 — La feria : présentations orales (20 min).} Chaque groupe présente son projet en \textbf{deux minutes}, en espagnol, sans lire intégralement sa fiche. Les autres groupes jouent le rôle d'\esp{inversores} et posent au moins une question chacun, par exemple : \esp{¿Cuánto dinero necesitan?}, \esp{¿A qué países exportarán?}, \esp{¿Quiénes serán sus clientes?}, \esp{¿Por qué es buena su idea?}

\textbf{Tâche 5 — Le vote final (10 min).} Chaque élève vote pour le meilleur projet (autre que le sien) et justifie oralement son choix avec la structure : \esp{Me gustaría invertir en... porque...} (J'aimerais investir dans... parce que...).

\courssub{Ressources à mobiliser}

\begin{aretenir}[Grammaire : les trois formes du projet]
\begin{itemize}
\item \textbf{Futur proche} (projet proche, déjà décidé) : \esp{ir a + infinitivo}. \esp{Vamos a crear una empresa.} « Nous allons créer une entreprise. »
\item \textbf{Futur simple} (prévision, promesse) : infinitif + \esp{-é, -ás, -á, -emos, -éis, -án}. \esp{Venderemos nuestro cacao en España.} « Nous vendrons notre cacao en Espagne. »
\item \textbf{Conditionnel} (souhait, politesse) : infinitif + \esp{-ía, -ías, -ía, -íamos, -íais, -ían}. \esp{Nos gustaría exportar a México.} « Nous aimerions exporter vers le Mexique. »
\end{itemize}
\end{aretenir}

\begin{propriete}[Rappel : futur simple et conditionnel de \esp{crear} et \esp{necesitar}]
\begin{center}
\begin{tabular}{|l|l|l|}
\hline
\rowcolor{popblueL} \textbf{Personne} & \textbf{Futur : \esp{crear}} & \textbf{Conditionnel : \esp{necesitar}} \\
\hline
yo & crearé & necesitaría \\
\hline
tú & crearás & necesitarías \\
\hline
él / ella / usted & creará & necesitaría \\
\hline
nosotros/as & crearemos & necesitaríamos \\
\hline
vosotros/as & crearéis & necesitaríais \\
\hline
ellos / ellas / ustedes & crearán & necesitarían \\
\hline
\end{tabular}
\end{center}
Les deux temps se construisent sur le même radical (l'infinitif) ; seules les terminaisons changent. Les radicaux irréguliers sont identiques aux deux temps : \esp{tener → tendr-}, \esp{poder → podr-}, \esp{hacer → har-}, \esp{haber → habr-}, \esp{querer → querr-}, \esp{saber → sabr-}, \esp{salir → saldr-}, \esp{venir → vendr-}, \esp{poner → pondr-}, \esp{decir → dir-}.
\end{propriete}

\begin{attention}[Pièges fréquents]
\begin{itemize}
\item Ne dis pas \esp{yo voy a crearé} : après \esp{ir a}, toujours l'\textbf{infinitif} (\esp{voy a crear}).
\item \esp{Me gustaría} est suivi de l'\textbf{infinitif} : \esp{me gustaría invertir}, jamais \esp{me gustaría invierto}.
\item Attention aux futurs irréguliers : \esp{tener → tendré}, \esp{poder → podré}, \esp{hacer → haré}, \esp{haber → habrá} (\esp{habrá muchos clientes} : il y aura beaucoup de clients).
\item \esp{el beneficio} se prononce « bé-né-fi-si-o » ; ne confonds pas \esp{empresa} (entreprise) et \esp{embarazada} (enceinte) !
\item « Pour » devant un infinitif de but se traduit par \esp{para} : \esp{para comprar las máquinas} (pour acheter les machines), jamais \esp{por}.
\end{itemize}
\end{attention}

\begin{definition}[Lexique à réutiliser obligatoirement]
\esp{emprendedor/a} (entrepreneur), \esp{la empresa} (l'entreprise), \esp{el proyecto} (le projet), \esp{el mercado} (le marché), \esp{el/la cliente/a} (le/la client(e)), \esp{la inversión} (l'investissement), \esp{el beneficio} (le bénéfice), \esp{exportar} (exporter), \esp{importar} (importer), \esp{vender} (vendre), \esp{los productos locales} (les produits locaux).
\end{definition}

\courssub{Proposition de production et corrigé commenté}

\begin{exoresolu}[Fiche-projet modèle : « Cacao Duala »]
Rédige la fiche-projet complète du groupe « Cacao Duala » en respectant les quotas grammaticaux (2 futurs proches, 2 futurs simples, 2 conditionnels), puis prépare la phrase de vote d'un investisseur.
\tcblower
\textbf{Fiche-projet :}\par
\esp{Nuestra empresa se llama « Cacao Duala ». Vamos a comprar el cacao de los campesinos de la región del Litoral y vamos a transformarlo en chocolate artesanal. Necesitaríamos una inversión de dos millones de francos CFA para comprar las máquinas. Venderemos nuestro chocolate en los mercados de Duala y Yaundé, y nuestros clientes serán las familias y las tiendas de la ciudad. En el futuro, exportaremos a España y a Guinea Ecuatorial. Nos gustaría ser la primera empresa de chocolate camerunés famosa en el mundo. ¡Habrá beneficios para todos!}\par
\textbf{Phrase de vote :} \esp{Me gustaría invertir en « Cacao Duala » porque el proyecto es realista y el chocolate camerunés es delicioso.}\par
\textbf{Commentaire des choix de langue :}
\begin{itemize}
\item \esp{Vamos a comprar... y vamos a transformarlo...} : deux \textbf{futurs proches} corrects — \esp{ir a} + infinitif ; note le pronom \esp{-lo} soudé à l'infinitif (\esp{transformarlo} : le transformer, c'est-à-dire le cacao).
\item \esp{Necesitaríamos... / Nos gustaría ser...} : deux \textbf{conditionnels} qui expriment le souhait et la demande polie, suivis de l'infinitif.
\item \esp{Venderemos... / exportaremos...} : deux \textbf{futurs simples} réguliers (\esp{-emos} à la première personne du pluriel) pour annoncer des prévisions.
\item \esp{Habrá beneficios} : futur irrégulier de \esp{haber} (il y aura), très utile dans un projet.
\item \esp{para comprar las máquinas} : \esp{para} + infinitif exprime le but (« pour acheter »).
\item Lexique de la leçon bien mobilisé : \esp{empresa}, \esp{inversión}, \esp{mercados}, \esp{clientes}, \esp{exportaremos}, \esp{beneficios}.
\item Le vote utilise la structure imposée : \esp{Me gustaría invertir en... porque...} + adjectifs de justification (\esp{realista}, \esp{delicioso}).
\end{itemize}
\end{exoresolu}

\begin{corrige}[Tâche 1]
\begin{enumerate}
\item \esp{El Ayuntamiento de Duala y el Liceo Bilingüe organizan la feria.} (La mairie de Douala et le Lycée Bilingue organisent la foire.)
\item \esp{El tema es « Mi proyecto para el Camerún ».} (Le thème est « Mon projet pour le Cameroun ».)
\item \esp{Se pueden utilizar el cacao, el café, la mandioca, el turismo o la artesanía.} (On peut utiliser le cacao, le café, le manioc, le tourisme ou l'artisanat.)
\item \esp{Cada grupo tiene dos minutos para presentar su proyecto.} (Chaque groupe a deux minutes pour présenter son projet.)
\end{enumerate}
\end{corrige}

\courssub{Bilan de l'activité}

Cette \esp{Feria de Jóvenes Emprendedores} a permis de réinvestir, dans une situation de communication authentique et motivante, l'ensemble des acquis de la leçon :

\begin{itemize}
\item \textbf{Compréhension de l'écrit} : lire et exploiter un appel à projets en espagnol ;
\item \textbf{Lexique} : réemployer dans un contexte réel les mots-clés de l'entreprise et du commerce (\esp{empresa}, \esp{inversión}, \esp{mercado}, \esp{clientes}, \esp{beneficios}, \esp{exportar}) ;
\item \textbf{Grammaire} : articuler les trois formes verbales du projet — futur proche pour les actions décidées, futur simple pour les prévisions, conditionnel pour les souhaits et les demandes polies ;
\item \textbf{Expression orale} : présenter en deux minutes, répondre aux questions des investisseurs, justifier un vote ;
\item \textbf{Compétence de vie} : concevoir un projet entrepreneurial réaliste valorisant les richesses du Cameroun et ouvert sur le monde hispanophone (Espagne, Guinée équatoriale, Amérique latine).
\end{itemize}

\begin{experience}[Pour aller plus loin]
En devoir à la maison, chaque élève rédige une courte lettre (8 à 10 lignes) à un investisseur espagnol imaginaire pour lui présenter le projet de son groupe, en commençant par \esp{Estimado señor :} et en utilisant au moins une fois \esp{me gustaría}, une fois le futur proche et une fois le futur simple. Les meilleures lettres seront lues en classe lors de la séance suivante.
\end{experience}

\newpage

\coursec{Méthodes et exercices résolus}

\coursec{Leçon EA14 : El emprendimiento y los intercambios comerciales}

\begin{textebox}[Texte support : Un sueño emprendedor en Duala]
Amina Njoya tiene veinticuatro años y vive en Duala, la capital económica de Camerún. Después de obtener su BTS en agroalimentaria en el Instituto Superior del Sahel, decidió no buscar un empleo asalariado. Su sueño era crear su propia empresa de jugos de frutas cien por cien naturales, sin conservantes ni azúcar añadido. Bautizó su marca « Vitalidad de Camerún ».

Al principio, las dificultades eran numerosas: ni local, ni máquinas, ni capital. Amina presentó su plan de negocio a un banco local y obtuvo un microcrédito de quinientos mil francos CFA. Compró una pasteurizadora de segunda mano y alquiló un pequeño taller en Makepe.

Hoy, dos años después, « Vitalidad de Camerún » produce doscientas botellas al día. Amina exporta sus jugos hacia Guinea Ecuatorial y Gabón, donde la demanda de productos naturales cameruneses es fuerte. Su empresa da empleo ahora a diez jóvenes de su barrio, entre ellos tres madres solteras.

« El emprendimiento no es fácil », dice Amina, « pero cuando veo las sonrisas de mis empleados y el orgullo de mis clientes, vale la pena. El año que viene, espero abrir una fábrica más grande y exportar hacia Europa. »
\end{textebox}

\begin{definition}[Lexique thématique : La empresa y el comercio]
\begin{center}
\begin{tabularx}{\linewidth}{|X|X|X|X|}
\hline
\rowcolor{popblueL}
\textbf{Español} & \textbf{Français} & \textbf{Prononciation} & \textbf{Remarque} \\
\hline
\esp{el emprendedor / la emprendedora} & l'entrepreneur / l'entrepreneuse & « em-pre-en-de-dor / -do-ra » & Accent sur l'antépénultième. \\
\hline
\esp{la empresa} & l'entreprise & « em-pre-sa » & Faux ami : \textbf{pas} « emprise ». \\
\hline
\esp{el proyecto} & le projet & « pro-yek-to » & \emph{j} aspiré (son [x]). \\
\hline
\esp{el mercado} & le marché & « mer-ka-do » & \esp{mercado internacional / nacional}. \\
\hline
\esp{el cliente} & le client & « kli-en-te » & Masculin invariable au pluriel : \esp{los clientes}. \\
\hline
\esp{el beneficio} & le bénéfice, le profit & « be-ne-fi-sio » & Faux ami : \textbf{pas} « bénéfice » au sens religieux. \\
\hline
\esp{la inversión} & l'investissement & « in-ver-sion » & Accent sur le \emph{ó} (mot aigu en -\emph{n}). \\
\hline
\esp{exportar} & exporter & « ex-por-tar » & \esp{exportar algo a un país}. \\
\hline
\esp{importar} & importer & « im-por-tar » & \esp{importar algo de un país}. \\
\hline
\esp{el socio comercial} & le partenaire commercial & « so-syo co-mer-sial & \emph{ci} = [si]. \\
\hline
\esp{la balanza comercial} & la balance commerciale & « ba-lan-tha co-mer-sial » & \emph{z} = [th] (Espagne) / [s] (Amérique). \\
\hline
\esp{el préstamo} & le prêt & « pres-ta-mo » & \esp{pedir un préstamo}. \\
\hline
\esp{el empleo} & l'emploi & « em-ple-o » & \esp{crear empleo, dar empleo a}. \\
\hline
\esp{la materia prima} & la matière première & « ma-te-ria pri-ma » & Pluriel : \esp{las materias primas}. \\
\hline
\end{tabularx}
\end{center}
\end{definition}

\begin{propriete}[Grammaire 1 : Le futur du projet — Futur proche et Futur simple]
\textbf{A. Futur proche : \esp{ir a + infinitivo}} \\
Utilisé pour une action \cle{déjà décidée}, \cle{prochaine} ou \cle{très probable}. Structure : \esp{ir} (présent indicatif) + \esp{a} + \emph{infinitif}.
\begin{center}
\begin{tabular}{|l|l|l|}
\hline
\rowcolor{popblueL}
\textbf{Personne} & \textbf{Ir (présent)} & \textbf{Exemple : \esp{crear}} \\
\hline
\esp{yo} & \esp{voy} & \esp{voy a crear} \\
\hline
\esp{tú} & \esp{vas} & \esp{vas a crear} \\
\hline
\esp{él / ella / usted} & \esp{va} & \esp{va a crear} \\
\hline
\esp{nosotros / -as} & \esp{vamos} & \esp{vamos a crear} \\
\hline
\esp{vosotros / -as} & \esp{vais} & \esp{vais a crear} \\
\hline
\esp{ellos / ellas / ustedes} & \esp{van} & \esp{van a crear} \\
\hline
\end{tabular}
\end{center}
\esp{¡Ojo!} Le \esp{a} est \textbf{obligatoire} : \esp{Voy a crear} (jamais \esp{*voy crear}). \\
\emph{Exemple :} \esp{Maana voy a presentar mi proyecto al banco.} (Demain je vais présenter mon projet à la banque.)

\textbf{B. Futur simple : \emph{infinitif} + désinences} \\
Utilisé pour une \cle{prévision}, une \cle{promesse}, un \cle{projet moins immédiat} ou une \cle{hypothèse sur le futur}. Désinences \textbf{identiques pour les 3 groupes} (-\emph{é, -ás, -á, -emos, -éis, -án}) — \textbf{toujours accentuées}.
\begin{center}
\begin{tabular}{|l|l|l|l|}
\hline
\rowcolor{popblueL}
\textbf{Personne} & \textbf{\esp{hablar}} & \textbf{\esp{comer}} & \textbf{\esp{vivir}} \\
\hline
\esp{yo} & \esp{hablaré} & \esp{comeré} & \esp{viviré} \\
\hline
\esp{tú} & \esp{hablarás} & \esp{comerás} & \esp{vivirás} \\
\hline
\esp{él/ella/usted} & \esp{hablará} & \esp{comerá} & \esp{vivirá} \\
\hline
\esp{nosotros/-as} & \esp{hablaremos} & \esp{comeremos} & \esp{viviremos} \\
\hline
\esp{vosotros/-as} & \esp{hablaréis} & \esp{comeréis} & \esp{viviréis} \\
\hline
\esp{ellos/ellas/ustedes} & \esp{hablarán} & \esp{comerán} & \esp{vivirán} \\
\hline
\end{tabular}
\end{center}

\textbf{C. Verbes irréguliers au futur simple (radical modifié + mêmes désinences)} \\
\begin{center}
\begin{tabularx}{\linewidth}{|X|X|X|}
\hline
\rowcolor{popblueL}
\textbf{Infinitif} & \textbf{Radical futur} & \textbf{1ère pers. sg (\esp{yo})} \\
\hline
\esp{tener} & \esp{tendr-} & \esp{tendré} \\
\hline
\esp{poder} & \esp{podr-} & \esp{podré} \\
\hline
\esp{querer} & \esp{querr-} & \esp{querré} \\
\hline
\esp{saber} & \esp{sabr-} & \esp{sabré} \\
\hline
\esp{poner} & \esp{pondr-} & \esp{pondré} \\
\hline
\esp{salir} & \esp{saldr-} & \esp{saldré} \\
\hline
\esp{venir} & \esp{vendr-} & \esp{vendré} \\
\hline
\esp{hacer} & \esp{har-} & \esp{haré} \\
\hline
\esp{decir} & \esp{dir-} & \esp{diré} \\
\hline
\esp{haber} & \esp{habr-} & \esp{habrá} (impersonnel) \\
\hline
\end{tabularx}
\end{center}
\emph{Exemples :} \esp{Tendré éxito.} (J'aurai du succès.) — \esp{¿Podrás ayudarme?} (Pourras-tu m'aider ?) — \esp{No habrá problema.} (Il n'y aura pas de problème.)
\end{propriete}

\begin{propriete}[Grammaire 2 : Le conditionnel du souhait — \esp{me gustaría + infinitivo}]
\textbf{Formation} : Mêmes radicaux irréguliers que le futur simple + désinences \textbf{-\emph{ía, -ías, -ía, -íamos, -íais, -ían}} (toujours accentuées sur le \emph{i}).
\begin{center}
\begin{tabular}{|l|l|}
\hline
\rowcolor{popblueL}
\textbf{Personne} & \textbf{\esp{gustar} (conditionnel)} \\
\hline
\esp{yo} & \esp{me gustaría} \\
\hline
\esp{tú} & \esp{te gustaría} \\
\hline
\esp{él/ella/usted} & \esp{le gustaría} \\
\hline
\esp{nosotros/-as} & \esp{nos gustaría} \\
\hline
\esp{vosotros/-as} & \esp{os gustaría} \\
\hline
\esp{ellos/ellas/ustedes} & \esp{les gustaría} \\
\hline
\end{tabular}
\end{center}

\textbf{Emploi principal (niveau A2/B1)} : Exprimer un \cle{souhait}, un \cle{désir poli}, un \cle{rêve} ou une \cle{politesse}.
\begin{itemize}
\item \textbf{Sujet identique} (celui qui souhaite = celui qui agit) $\to$ \esp{me gustaría + infinitif}.
\esp{Me gustaría crear una empresa.} (J'aimerais créer une entreprise.)
\item \textbf{Sujet différent} $\to$ \esp{me gustaría que + subjonctif} (niveau B1/B2, simple reconnaissance ici).
\esp{Me gustaría que mi hermano me ayudara.}
\end{itemize}

\textbf{Autres verbes fréquents au conditionnel de politesse/souhait} :
\esp{querría} (je voudrais, de \esp{querer}), \esp{podría} (je pourrais, de \esp{poder}), \esp{debería} (je devrais, de \esp{deber}).
\emph{Attention :} \esp{Me gustaría} s'emploie \textbf{sans subjonctif} quand le sujet ne change pas. Erreur fréquente : \esp{*Me gustaría que creo} $\to$ \textbf{Correct :} \esp{Me gustaría crear}.
\end{propriete}

\begin{aretenir}[Résumé : Choisir son futur pour un projet]
\begin{center}
\begin{tabularx}{\linewidth}{|X|X|X|}
\hline
\rowcolor{popblueL}
\textbf{Intention} & \textbf{Structure} & \textbf{Exemple} \\
\hline
Projet décidé, étape prochaine, plan concret & \esp{ir a + infinitivo} (Futur proche) & \esp{Voy a buscar un local mañana.} \\
\hline
Souhait, rêve, politesse, conditionnel & \esp{me gustaría + infinitivo} & \esp{Me gustaría exportar a España.} \\
\hline
Promesse, prévision, engagement formel & \esp{Futur simple} & \esp{Crearé empleo para diez jóvenes.} \\
\hline
\end{tabularx}
\end{center}
\end{aretenir}

\courssub{Savoir-faire 1 : Lire et comprendre un texte sur l'entrepreneuriat}

\begin{methode}[Comprendre un texte économique en 5 étapes]
\begin{enumerate}
\item \textbf{Première lecture globale} : lis le texte une fois sans dictionnaire. Repère le \cle{thème} (emprendimiento), le \cle{personnage} (\esp{¿quién?}), le \cle{lieu} (\esp{¿dónde?}) et le \cle{problème/but} (\esp{¿qué dificultad / qué objetivo?}).
\item \textbf{Repérage du lexique clé} : souligne les mots du champ lexical de l'économie : \esp{emprendedor, empresa, proyecto, mercado, cliente, beneficio, inversión, exportar, importar, préstamo, empleo}. Ils portent le sens du texte.
\item \textbf{Lecture analytique} : pour chaque paragraphe, formule une idée principale en une phrase simple en espagnol : \esp{Amina quiere crear una empresa} / \esp{Necesita un préstamo} / \esp{Exporta a Guinea Ecuatorial}.
\item \textbf{Réponse aux questions} : cite le texte en espagnol (\esp{El texto dice que...}) puis reformule avec tes mots. Ne recopie jamais un paragraphe entier.
\item \textbf{Vérification} : ta réponse doit être en espagnol correct, avec verbe conjugué, accords (genre/nombre) et accents.
\end{enumerate}
\textbf{Structures à retenir} : \esp{El texto trata de...} (le texte parle de...) ; \esp{Según el texto...} (selon le texte) ; \esp{Esto significa que...} (cela signifie que...) ; \esp{La idea principal es...} (l'idée principale est...).
\end{methode}

\begin{exoresolu}[Compréhension de l'écrit : « Un sueño emprendedor en Duala »]
\textbf{Énoncé.} Lis le texte support (ci-dessus) et réponds en espagnol aux questions suivantes.
\begin{enumerate}
\item ¿Qué decidió hacer Amina después de sus estudios?
\item ¿Cómo consiguió el dinero para su empresa?
\item ¿A qué países exporta sus productos actualmente?
\item ¿Cuál es el beneficio social de su proyecto?
\item ¿Qué proyecto tiene para el futuro?
\end{enumerate}
\tcblower
\textbf{Raisonnement.} Chaque question commence par un mot interrogatif qui indique l'information cherchée : \esp{qué} (action/objet), \esp{cómo} (moyen/manière), \esp{a qué países} (lieu/direction), \esp{cuál} (choix/identification), \esp{qué proyecto} (contenu futur). On repère dans le texte le paragraphe correspondant, on identifie le temps verbal (passé pour les faits accomplis, présent pour la situation actuelle, futur pour les projets) et on reformule.

\textbf{Réponses rédigées.}
\begin{enumerate}
\item \esp{Después de sus estudios, Amina decidió crear su propia empresa de jugos de frutas naturales.} — Structure \esp{decidir + infinitivo} (sans préposition). \esp{Decidió} = passé simple (indéfini) 3e pers. sg.
\item \esp{Consiguió el dinero gracias a un microcrédito (un préstamo) de un banco local.} — \esp{Conseguir} (e$\to$i) : \esp{consiguió}. \esp{Gracias a} = « grâce à ».
\item \esp{Actualmente exporta sus jugos a Guinea Ecuatorial y a Gabón.} — Présent de l'indicatif \esp{exporta} (vérité actuelle). \esp{A} marque la direction.
\item \esp{El beneficio social es que su empresa da empleo a diez jóvenes de su barrio, incluidas tres madres solteras.} — \esp{Dar empleo a} = « donner de l'emploi à ». Accord : \esp{incluidas} (féminin pluriel accordé avec \esp{madres}).
\item \esp{Para el futuro, espera abrir una fábrica más grande y exportar hacia Europa.} — \esp{Espera} (présent) + infinitif = intention proche ; \esp{abrir, exportar} infinitifs.
\end{enumerate}
\textbf{Conclusion.} Une bonne réponse de compréhension est courte, exacte, reprend le temps du texte et soigne les accords.
\end{exoresolu}

\courssub{Savoir-faire 2 : Présenter un projet entrepreneurial}

\begin{methode}[Présenter un projet en 5 étapes]
\begin{enumerate}
\item \textbf{Se présenter} : \esp{Me llamo... y soy estudiante de... en...}
\item \textbf{Annoncer le projet au futur proche} : \esp{Voy a crear una empresa de...} (\cle{ir a + infinitivo} = projet proche, déjà décidé, planifié).
\item \textbf{Expliquer la motivation au conditionnel} : \esp{Me gustaría ayudar a los jóvenes de mi barrio.} (\cle{me gustaría + infinitivo} = souhait poli, rêve).
\item \textbf{Décrire les étapes concrètes au futur simple} : \esp{Primero buscaré un local, después compraré las máquinas y contrataré a cinco personas.} (Futur simple = actions prévues, promesses, engagement).
\item \textbf{Conclure sur les bénéfices (économiques/sociaux)} : \esp{Así, mi empresa creará empleo y venderá productos de calidad.}
\end{enumerate}
\textbf{Connecteurs logiques à utiliser} : \esp{Para empezar / En primer lugar} (pour commencer), \esp{luego / después} (ensuite), \esp{finalmente / al final} (finalement), \esp{por eso / así que} (donc), \esp{además} (de plus).
\end{methode}

\begin{exoresolu}[Expression écrite : Présente ton projet pour le concours « Jóvenes Emprendedores »]
\textbf{Énoncé.} Tu es élève en Première A à Yaoundé. Rédige un texte de 8 à 10 lignes en espagnol pour présenter ton projet entrepreneurial. Utilise \textbf{au moins une fois} : \esp{voy a + infinitivo}, \esp{me gustaría + infinitivo} et \textbf{deux verbes au futur simple} (réguliers ou irréguliers).
\tcblower
\textbf{Raisonnement.} On suit le plan de la méthode : présentation (présent), annonce (futur proche), motivation (conditionnel), étapes (futur simple), conclusion (futur simple). On vérifie chaque forme verbale : \esp{voy a crear} (ir 1ère pers. + a + inf.) ; \esp{me gustaría} (conditionnel invariable à cette personne + inf.) ; futurs simples : \esp{necesitaré, compraré, contrataré, crearé, ganaré, tendré}.

\textbf{Production modèle.}
\esp{Me llamo Nadège y soy estudiante de Primera A en Yaoundé. Voy a crear una pequeña empresa de pasteles y jugos naturales caseros. Me gustaría vender productos sanos y ricos a los estudiantes de mi barrio. Para empezar, necesitaré un pequeño local cerca del mercado central. Después compraré un horno y una nevera grande. También contrataré a dos jóvenes para repartir los pedidos en moto. Con este proyecto, crearé empleo y ganaré un beneficio para ayudar a mi familia. Estoy segura de que mi empresa tendrá éxito.}

\textbf{Justifications grammaticales.}
\begin{itemize}
\item \esp{Voy a crear} : Futur proche, projet décidé. \textbf{Jamais} \esp{voy crear} (sans \esp{a}).
\item \esp{Me gustaría vender} : Conditionnel de politesse/souhait + infinitif (sujet identique $\to$ pas de subjonctif).
\item \esp{necesitaré, compraré, contrataré, crearé, ganaré} : Futur simple régulier (infinitif + \emph{é, ás, á, emos, éis, án}). Accent obligatoire sur la dernière syllabe (sauf nous/vous).
\item \esp{tendrá éxito} : Futur irrégulier de \esp{tener} (radical \esp{tendr-}) + \emph{á}. \esp{Éxito} prend l'accent (mot aigu en voyelle).
\item \esp{Estoy segura} : \esp{estar} + adjectif d'état (accord féminin). \esp{Ser} serait pour une qualité permanente.
\end{itemize}
\textbf{Conclusion.} Un bon texte de projet alterne les trois modalités (proche, souhait, prévision) et reste simple : mieux vaut des phrases courtes et correctes que des phrases longues fautives.
\end{exoresolu}

\courssub{Savoir-faire 3 : Parler des échanges commerciaux}

\begin{methode}[Décrire des échanges commerciaux en 4 étapes]
\begin{enumerate}
\item \textbf{Nommer les partenaires} : \esp{Camerún comercia con España, Guinea Ecuatorial y otros países hispanohablantes.}
\item \textbf{Dire ce qu'on exporte} : \esp{Camerún exporta cacao, café, madera y petróleo al mercado internacional.} (Verbe \cle{exportar} + produit + \esp{a/al} pays ou marché).
\item \textbf{Dire ce qu'on importe} : \esp{Importa máquinas, vehículos, medicamentos y tecnología de España.} (Verbe \cle{importar} + produit + \esp{de} pays d'origine).
\item \textbf{Évaluer l'échange / Perspectives} : \esp{Estos intercambios son importantes para la economía nacional.} / \esp{En el futuro, Camerún exportará más productos transformados.}
\end{enumerate}
\textbf{Structures à retenir} : \esp{el mercado nacional / internacional} ; \esp{los productos de exportación / de importación} ; \esp{un socio comercial} ; \esp{la balanza comercial} (excédentaire / déficitaire) ; \esp{a cambio} (en échange) ; \esp{sobre todo} (surtout).
\end{methode}

\begin{exoresolu}[Expression orale préparée : Mini-exposé sur les échanges Cameroun / Monde hispanophone]
\textbf{Énoncé.} Prépare un exposé de 6 à 8 phrases sur : \esp{Los intercambios comerciales entre Camerún y el mundo hispanohablante}. Tu dois utiliser : \esp{exportar}, \esp{importar}, \esp{vender}, \esp{mercado}.
\tcblower
\textbf{Raisonnement.} On construit l'exposé : partenaires $\to$ exportations $\to$ importations $\to$ spécificité Guinée équatoriale $\to$ bilan $\to$ perspective future. Temps : présent (faits), futur (perspective). Attention aux articles : \esp{el cacao, el café, la madera, el petróleo, la fruta}.

\textbf{Exposé modèle.}
\esp{Camerún mantiene intercambios comerciales con varios países hispanohablantes, como España y Guinea Ecuatorial. Nuestro país exporta cacao, café, madera y petróleo al mercado internacional. España compra sobre todo cacao camerunés. A cambio, Camerún importa máquinas, vehículos y medicamentos de España. Con Guinea Ecuatorial, el comercio es más fácil porque es un país vecino: los comerciantes venden frutas, verduras y pescado en los mercados de la frontera. Estos intercambios crean empleo y traen divisas al país. En el futuro, Camerún exportará más productos transformados, como chocolate, y no solamente materias primas.}

\textbf{Justifications.} \esp{A cambio} = « en échange » (locution figée) ; \esp{sobre todo} = « surtout » (faux ami : ne pas traduire « au-dessus de tout ») ; \esp{materias primas} = « matières premières » (pluriel) ; \esp{exportará} : futur simple pour perspective ; \esp{divisas} = « devises ».
\textbf{Conclusion.} Un exposé commercial efficace utilise le présent pour les faits constants, le futur pour les perspectives, et un lexique précis : \esp{exportar / importar / socio comercial / mercado / balanza comercial}.
\end{exoresolu}

\begin{savaistu}[Le Cameroun et la Guinée équatoriale : une zone d'échange privilégiée]
La Guinée équatoriale (\esp{Guinea Ecuatorial}) est le seul pays africain dont la langue officielle est l'espagnol. Elle partage plus de 180 km de frontière avec le Cameroun (région du Sud). Les échanges y sont intenses : le Cameroun y exporte des produits agricoles (bananes, plantains, cacao, bois) et y importe du pétrole et du gaz. La ville de \esp{Kyé-Ossi} (Cameroun) et \esp{Ebebiyín} (Guinée équatoriale) forment un axe commercial transfrontalier majeur. La monnaie commune (FCFA) facilite les transactions. Pour un jeune camerounais, parler espagnol est un atout professionnel direct pour travailler dans le commerce, la logistique ou la douane à cette frontière.
\end{savaistu}

\courssub{Savoir-faire 4 : Traduire (thème) avec futur et conditionnel}

\begin{methode}[Réussir le thème en 5 étapes]
\begin{enumerate}
\item \textbf{Identifier le temps français} : « je créerai » = futur simple ; « je vais créer » = futur proche (\esp{ir a + inf.}) ; « j'aimerais » = conditionnel (\esp{me gustaría}).
\item \textbf{Choisir le bon auxiliaire/structure} : \esp{voy a crear} (projet décidé) vs \esp{crearé} (promesse/prévision).
\item \textbf{Traduire le vocabulaire exactement} : « entreprise » = \esp{empresa} ; « bénéfice » = \esp{beneficio} ; « client » = \esp{cliente} ; « investissement » = \esp{inversión} ; « transporter » = \esp{transportar}, nom : \esp{el transporte}.
\item \textbf{Accorder} : articles et adjectifs avec le nom (\esp{una empresa pequeña}, \esp{los clientes satisfechos}, \esp{el mercado internacional}).
\item \textbf{Relire : accents et ponctuation} : accents sur futurs (\esp{crearé, venderá}), conditionnels (\esp{gustaría, podría}) ; \esp{¿ ?} \esp{¡ !} si question/exclamation.
\end{enumerate}
\end{methode}

\begin{exoresolu}[Thème : phrases à traduire en espagnol]
\textbf{Énoncé.} Traduis en espagnol :
\begin{enumerate}
\item Je vais créer une entreprise de transport à Douala.
\item J'aimerais exporter le cacao camerounais vers l'Espagne.
\item Mon frère cherchera des clients sur le marché international.
\item Nous aurons besoin d'un investissement important.
\item Pourront-ils vendre leurs produits au Gabon ?
\end{enumerate}
\tcblower
\textbf{Solutions et justifications.}
\begin{enumerate}
\item \esp{Voy a crear una empresa de transporte en Duala.} — « Je vais créer » : futur proche \esp{voy a + infinitivo}. « Transport » = \esp{transporte} (masculin, \emph{e} final). \esp{En Duala} (dans Douala).
\item \esp{Me gustaría exportar el cacao camerunés a España.} — « J'aimerais » : \esp{me gustaría} + infinitif (même sujet). « Vers l'Espagne » = \esp{a España} (direction). \esp{Camerunés} : accent sur \emph{é} (mot aigu en -\emph{s}).
\item \esp{Mi hermano buscará clientes en el mercado internacional.} — Futur simple de \esp{buscar} : \esp{buscará} (régulier, accent sur \emph{á}). « Sur le marché » = \esp{en el mercado}.
\item \esp{Necesitaremos una inversión importante.} — Futur simple 1ère pers. pl. : \esp{necesitaremos}. « Avoir besoin de » = \esp{necesitar} (éviter le calque \esp{*tendremos necesidad de}). \esp{Inversión} : accent sur \emph{ó} (mot aigu en -\emph{n}).
\item \esp{¿Podrán vender sus productos en Gabón?} — Futur simple irrégulier de \esp{poder} : \esp{podrán} (radical \esp{podr-}). Ponctuation \esp{¿ ?}.
\end{enumerate}
\textbf{Conclusion.} Au thème, chaque point vient de trois choses : le bon temps, le bon mot, le bon accent. Relis toujours ces trois éléments.
\end{exoresolu}

\courssub{Savoir-faire 5 : Traduire (version) un texte commercial}

\begin{methode}[Réussir la version en 4 étapes]
\begin{enumerate}
\item \textbf{Traduire phrase par phrase}, en respectant l'ordre des mots espagnol (adjectif souvent après le nom : \esp{productos naturales} = « produits naturels »).
\item \textbf{Repérer les temps} : Passé simple (\esp{creó, vendió, trabajó}) = passé composé ou passé simple français ; Imparfait (\esp{tenía, quería}) = imparfait français ; Présent (\esp{exportan, piensan}) = présent français.
\item \textbf{Attention aux faux amis} : \esp{empresa} = « entreprise » ; \esp{beneficio} = « bénéfice, profit » ; \esp{actualmente} = « actuellement » (pas « actuellement » = « vraiment ») ; \esp{realizar} = « accomplir, réaliser » (pas « se rendre compte ») ; \esp{éxito} = « succès » (pas « sortie »).
\item \textbf{Relire en français} : la phrase doit être naturelle, complète et respecter la concordance des temps.
\end{enumerate}
\end{methode}

\begin{exoresolu}[Version : un paragraphe commercial]
\textbf{Énoncé.} Traduis en français :
\esp{« El año pasado, un grupo de jóvenes de Bafoussam creó una cooperativa para vender café. Al principio tenían pocos clientes, pero trabajaron mucho y mejoraron la calidad de su producto. Actualmente exportan su café a España y piensan abrir una tienda en Internet. »}
\tcblower
\textbf{Raisonnement.} \esp{El año pasado} = « l'année dernière » (repère passé) ; \esp{creó} (passé simple 3e sg) = « a créé » ; \esp{tenían} (imparfait 3e pl) = « avaient » (description état passé) ; \esp{trabajaron, mejoraron} (passé simple 3e pl) = « ont travaillé, ont amélioré » (actions ponctuelles) ; \esp{actualmente} = « actuellement » (faux ami : \textbf{pas} « effectivement ») ; \esp{piensan abrir} = « ils pensent ouvrir / ont l'intention d'ouvrir » (\esp{pensar + infinitivo} = intention).

\textbf{Traduction.}
« L'année dernière, un groupe de jeunes de Bafoussam a créé une coopérative pour vendre du café. Au début, ils avaient peu de clients, mais ils ont beaucoup travaillé et ont amélioré la qualité de leur produit. Actuellement, ils exportent leur café vers l'Espagne et pensent ouvrir une boutique en ligne. »

\textbf{Justifications.} \esp{Al principio} = « au début / au commencement » ; \esp{pocos clientes} = « peu de clients » ; \esp{mejorar la calidad} = « améliorer la qualité » ; \esp{tienda} = « boutique, magasin » (faux ami : \textbf{pas} « tente » = \esp{tienda de campaña}) ; \esp{en Internet} = « en ligne / sur Internet ».
\textbf{Conclusion.} En version, la fidélité aux temps (passé simple vs imparfait) et la vigilance sur les faux amis (\esp{actualmente, tienda, realizar, éxito}) font la différence.
\end{exoresolu}

\courssub{Savoir-faire 6 : Commenter un texte sur l'entrepreneuriat}

\begin{methode}[Commentaire de texte en 4 étapes]
\begin{enumerate}
\item \textbf{Introduction} (2-3 lignes) : Présente le texte (thème, source si connue, personnage principal, lieu). \esp{Este texto trata de una joven emprendedora de Duala que crea su empresa de jugos.}
\item \textbf{Étude des idées} (2 idées principales, 4-5 lignes chacune) : Dégage 2 ou 3 idées forces (ex: la naissance du projet, les difficultés, la réussite, l'impact social). Pour chaque idée : \textbf{cite le texte en espagnol} (entre guillemets) $\to$ \textbf{explique} avec tes mots.
\item \textbf{Étude de la langue} (2-3 lignes) : Relève le \cle{lexique économique} (\esp{empresa, mercado, exportar, beneficio, préstamo, empleo}) et les \cle{temps verbaux} utilisés (passé pour le récit, présent pour la situation, futur pour les projets).
\item \textbf{Conclusion / Avis personnel} (2-3 lignes) : Donne ton avis argumenté. \esp{En mi opinión, el emprendimiento es vital para África porque crea empleo y frena la emigración.}
\end{enumerate}
\textbf{Structures à retenir} : \esp{El autor muestra que...} ; \esp{Esto demuestra que...} ; \esp{Para concluir...} ; \esp{En mi opinión...} ; \esp{Es importante destacar que...}.
\end{methode}

\begin{exoresolu}[Commentaire guidé]
\textbf{Énoncé.} Texte : \esp{« Carlos, un joven de Malabo, quería ser empresario. Con poco dinero, abrió un taller de reparación de teléfonos. Hoy tiene tres tiendas y da trabajo a ocho personas. En el futuro, exportará sus servicios a otros países de la región. »}

Rédige un commentaire de 10 à 12 lignes en suivant le plan : introduction, deux idées principales, étude de la langue, avis personnel.
\tcblower
\textbf{Raisonnement.} Idée 1 : Volonté et départ modeste (\esp{quería ser empresario, con poco dinero, abrió un taller}). Idée 2 : Réussite actuelle et expansion future (\esp{hoy tiene tres tiendas, da trabajo, exportará}). Langue : Imparfait (\esp{quería}), passé simple (\esp{abrió}), présent (\esp{tiene, da}), futur (\esp{exportará}). Lexique : \esp{empresario, taller, tiendas, trabajo, exportar}.

\textbf{Commentaire modèle.}
\esp{Este texto trata de Carlos, un joven emprendedor de Malabo, en Guinea Ecuatorial, que crea su propia empresa. La primera idea es la voluntad del protagonista: el texto dice que « quería ser empresario » y que empezó « con poco dinero ». Esto demuestra que no hace falta ser rico para emprender; hace falta una idea y esfuerzo. La segunda idea es el éxito del proyecto: hoy Carlos « tiene tres tiendas y da trabajo a ocho personas ». Su empresa crea empleo para su comunidad. Además, el futuro « exportará » muestra que el proyecto todavía puede crecer. En mi opinión, este texto es una lección para los jóvenes africanos: el emprendimiento permite realizar los sueños y ayudar al país. Para concluir, pienso que todos podemos crear algo si trabajamos con paciencia.}

\textbf{Justifications.} \esp{No hace falta} = « il n'est pas nécessaire / il ne faut pas » (impersonnel + infinitif) ; \esp{esfuerzo} = « effort » ; \esp{crecer} = « grandir, se développer » ; \esp{realizar los sueños} = « réaliser ses rêves » (faux ami : \esp{realizar} = accomplir, \textbf{pas} « se rendre compte » = \esp{darse cuenta}) ; \esp{si trabajamos} = « si nous travaillons » (si + présent $\to$ futur/présent dans la principale, condition réelle).
\textbf{Conclusion.} Un bon commentaire cite le texte, explique les idées, analyse la langue et finit par un avis personnel argumenté, le tout en espagnol simple et correct.
\end{exoresolu}

\begin{experience}[Jeu de rôle : « Pitch ton projet » (3 minutes par binôme)]
\textbf{Consigne.} En binôme. L'élève A est le jeune entrepreneur ; l'élève B est un banquier / investisseur.
\begin{enumerate}
\item \textbf{Entrepreneur (A)} : Se présente, présente son projet (\esp{voy a crear...}), sa motivation (\esp{me gustaría...}), ses besoins (\esp{necesitaré...}), l'impact social (\esp{crearé empleo...}).
\item \textbf{Banquier (B)} : Pose 3 questions : \esp{¿Cuánto dinero necesita?}, \esp{¿A qué mercado se dirige?}, \esp{¿Cuándo empezará a ganar beneficios?}.
\item \textbf{Entrepreneur (A)} : Répond en utilisant le futur simple et le conditionnel (\esp{Esperaré..., Podría..., Tendré...}).
\item \textbf{Échange} : Le banquier accepte (\esp{Le doy el préstamo}) ou refuse (\esp{No es viable}) avec justification.
\end{enumerate}
\textbf{Objectif :} Réactiver \esp{ir a + inf.}, \esp{me gustaría}, futur simple, lexique \esp{préstamo, mercado, beneficio, inversión, cliente}.
\end{experience}

\begin{attention}[Les 5 erreurs qui coûtent des points au Bac]
\begin{enumerate}
\item \textbf{Oublier les accents des futurs et conditionnels} : écris \esp{crearé, venderá, gustaría, podría, tendrá}. Sans accent, \esp{cree} (présent subjonctif/indicatif) ou \esp{vendera} (forme inexistante) sont fautifs : l'accent marque le temps et la personne.
\item \textbf{Calquer « je vais créer » par \esp{voy crear}} : le futur proche exige \esp{ir a + infinitivo} : \esp{voy a crear}. Le \esp{a} est \textbf{obligatoire}.
\item \textbf{Mettre le subjonctif après \esp{me gustaría}} quand le sujet ne change pas : dis \esp{Me gustaría crear una empresa} (infinitif), jamais \esp{*me gustaría que cree}.
\item \textbf{Les faux amis économiques} : \esp{empresa} = « entreprise » (pas « emprise ») ; \esp{beneficio} = « bénéfice, profit » ; \esp{actualmente} = « actuellement » (pas « effectivement ») ; \esp{realizar} = « accomplir, concrétiser » (pas « se rendre compte ») ; \esp{éxito} = « succès » (pas « sortie ») ; \esp{tienda} = « boutique » (pas « tente » sauf \esp{tienda de campaña}).
\item \textbf{Les accords avec le vocabulaire commercial} : \esp{una empresa pequeña} (féminin sg), \esp{los clientes satisfechos} (masculin pl), \esp{el mercado internacional} (masculin sg), \esp{unas inversiones importantes} (féminin pl). Vérifie chaque article et chaque adjectif avant de rendre ta copie.
\end{enumerate}
\end{attention}

\begin{exercice}[Entraînement personnel : Thème, Version, Expression]
\textbf{Exercice 1 (Thème).} Traduis en espagnol :
\begin{enumerate}
\item Nous allons ouvrir une boutique de vêtements à Yaoundé.
\item Ils aimeraient importer des machines d'Espagne.
\item Mon père aura besoin d'un prêt pour son projet.
\item Vous (vous autres) vendrez des fruits sur le marché local.
\item Je créerai une application pour les agriculteurs.
\end{enumerate}

\textbf{Exercice 2 (Version).} Traduis en français :
\esp{« Mi hermana mayor, Elena, es una emprendedora nata. Hace tres años, montó una pequeña fábrica de jabón artesanal en su pueblo, cerca de Bafía. Al principio no tenía clientes, pero ahora exporta sus jabones a Francia y a Bélgica. El año que viene, quiere contratar a más mujeres del pueblo para aumentar la producción. »}

\textbf{Exercice 3 (Expression écrite).} Sujet : \esp{« Mi proyecto para mejorar mi barrio »}. Rédige 10-12 lignes. Contraintes : utilise \esp{voy a + inf.} (1x), \esp{me gustaría + inf.} (1x), futur simple (3 verbes minimum dont un irrégulier : \esp{tendré, podré, haré, sabré, pondré, saldré, vendré, diré, querré}). Lexique imposé : \esp{empresa, mercado, beneficio, empleo, inversión}.
\end{exercice}

\begin{corrige}[Exercice d'entraînement]
\textbf{Exercice 1 (Thème).}
\begin{enumerate}
\item \esp{Vamos a abrir una tienda de ropa en Yaundé.} (\esp{abrir} régulier, \esp{ropa} = vêtements, \esp{Yaundé} graphie espagnole).
\item \esp{Les gustaría importar máquinas de España.} (\esp{les gustaría} 3e pl, \esp{máquinas} fém. pl, \esp{de España} origine).
\item \esp{Mi padre necesitará un préstamo para su proyecto.} (\esp{necesitará} futur simple, \esp{préstamo} accent).
\item \esp{Venderéis frutas en el mercado local.} (\esp{venderéis} futur 2e pl - \emph{vostros} forme scolaire).
\item \esp{Crearé una aplicación para los agricultores.} (\esp{crearé} futur 1e sg, \esp{agricultores} masc. pl).
\end{enumerate}

\textbf{Exercice 2 (Version).}
« Ma grande sœur, Elena, est une entrepreneuse née. Il y a trois ans, elle a monté une petite fabrique de savon artisanal dans son village, près de Bafia. Au début, elle n'avait pas de clients, mais maintenant elle exporte ses savons en France et en Belgique. L'année prochaine, elle veut embaucher plus de femmes du village pour augmenter la production. »
\emph{Notes :} \esp{emprendedora nata} = « entrepreneuse née » ; \esp{montó} (passé simple \esp{montar}) = « a monté / a créé » ; \esp{jabón artesanal} = « savon artisanal » ; \esp{aumentar la producción} = « augmenter la production ».

\textbf{Exercice 3 (Expression écrite - Proposition de correction).}
\esp{Me llamo Paul y vivo en un barrio popular de Duala. Voy a crear una empresa de reciclaje de plástico para limpiar las calles. Me gustaría transformar los residuos en objetos útiles como macetas o bancos. Para empezar, necesitaré un terreno pequeño y una máquina trituradora. Después, contrataré a cinco jóvenes sin empleo para recoger el plástico. También haré campañas de sensibilización en las escuelas. Con este proyecto, tendré un beneficio económico y crearé empleo. Además, el barrio estará más limpio. ¡Estoy seguro de que podré hacerlo!}
\emph{Vérification :} \esp{voy a crear} (futur proche), \esp{me gustaría transformar} (conditionnel), \esp{necesitaré, contrataré, haré, tendré, crearé, podré} (7 futurs simples dont 4 irréguliers : \esp{haré, tendré, podré, crearé} - \esp{crearé} régulier). Lexique tout présent.
\end{corrige}

\begin{popfigure}
\begin{tikzpicture}[
  mindmap,
  grow cyclic,
  every node/.style=concept,
  concept color=popblue,
  level 1/.append style={level distance=3.5cm, sibling angle=90},
  level 2/.append style={level distance=2.5cm, sibling angle=45},
  text=white,
  font=\small
]
\node {El emprendimiento}
  child[concept color=popgreen] { node[align=center]{La idea\\y el proyecto}
    child { node {Innovación} }
    child { node {Necesidad del mercado} }
  }
  child[concept color=poporange] { node[align=center]{El capital\\y la inversión}
    child { node {Préstamo bancario} }
    child { node {Microcrédito} }
    child { node {Ahorro propio} }
  }
  child[concept color=poppurple] { node[align=center]{El mercado\\y los clientes}
    child { node {Mercado local} }
    child { node {Exportación} }
    child { node {Comercio transfronterizo} }
  }
  child[concept color=poppink] { node[align=center]{El impacto\\social}
    child { node {Crear empleo} }
    child { node {Beneficio comunitario} }
    child { node {Desarrollo local} }
  };
\end{tikzpicture}
\legende{Carte mentale : Les 4 piliers de l'entrepreneuriat (vocabulaire clé de la leçon).}
\end{popfigure}

\newpage

\coursec{Exercices}

\courssub{Je vérifie mes connaissances}

\begin{exercice}[QCM : le vocabulaire de l'entreprise]
Entoure la bonne réponse.
\begin{enumerate}
\item Una persona que crea una empresa es :
\begin{itemize}
\item un cliente \item un emprendedor \item un mercado
\end{itemize}
\item \esp{Exportar} significa :
\begin{itemize}
\item vender productos al extranjero \item comprar productos en el extranjero \item fabricar productos
\end{itemize}
\item El dinero que se gana con un negocio es :
\begin{itemize}
\item la inversión \item el beneficio \item el préstamo
\end{itemize}
\item \esp{Me gustaría crear una empresa} expresa :
\begin{itemize}
\item un hecho pasado \item una certeza absoluta \item un deseo
\end{itemize}
\end{enumerate}
\end{exercice}

\begin{exercice}[Vrai ou faux ? Justifie ta réponse]
Réponds par V (vrai) ou F (faux) et justifie en une phrase en français.
\begin{enumerate}
\item \esp{Voy a abrir una tienda mañana} exprime un futur proche.
\item \esp{Crearé} est la première personne du singulier du conditionnel de \esp{crear}.
\item Le Cameroun exporte du cacao vers l'Espagne.
\item \esp{Me gustaría} se traduit par « je voudrais / j'aimerais ».
\item \esp{Importar} signifie « vendre à l'étranger ».
\end{enumerate}
\end{exercice}

\begin{exercice}[Trouve le mot]
Complète chaque phrase avec le mot espagnol qui convient, choisi dans la liste : \esp{mercado -- cliente -- proyecto -- inversión -- empresa}.
\begin{enumerate}
\item Mi \ldots\ldots\ldots es vender jugo natural en el barrio.
\item Para empezar, necesito una \ldots\ldots\ldots de 100\,000 francos.
\item Cada \ldots\ldots\ldots compra dos botellas por día.
\item Los sábados, vendo mis productos en el \ldots\ldots\ldots de Mokolo.
\item Quiero crear una pequeña \ldots\ldots\ldots con mis hermanos.
\end{enumerate}
\end{exercice}

\begin{exercice}[Conjugaison à compléter]
Conjugue les verbes entre parenthèses au futur simple ou au conditionnel selon le sens.
\begin{enumerate}
\item El año próximo, yo \ldots\ldots\ldots (crear) mi empresa.
\item ¿Le \ldots\ldots\ldots (gustar) a usted invertir en mi proyecto?
\item Nosotros \ldots\ldots\ldots (exportar) cacao a España el próximo año.
\item Ella \ldots\ldots\ldots (querer) abrir una tienda en Duala, pero no tiene dinero.
\item Yo \ldots\ldots\ldots (poder) ayudarte, pero no tengo tiempo.
\end{enumerate}
\end{exercice}

\courssub{Je m'entraîne}

\begin{exercice}[Traduis en espagnol]
\begin{enumerate}
\item Je vais ouvrir un restaurant à Yaoundé.
\item J'aimerais exporter mon café vers l'Espagne.
\item Nous créerons une entreprise l'année prochaine.
\item Les clients achètent beaucoup de cacao au marché.
\item Elle voudrait trouver des investisseurs au Mexique.
\end{enumerate}
\end{exercice}

\begin{exercice}[Transforme les phrases]
Mets chaque phrase au futur proche (\esp{ir a + infinitivo}), puis au conditionnel avec \esp{me gustaría}.
\begin{enumerate}
\item Creo una empresa de transporte.
\item Exportamos frutas a Guinea Ecuatorial.
\item Abro una tienda en el mercado central.
\item Busco clientes en España.
\end{enumerate}
Exemple : \esp{Compro una moto.} $\rightarrow$ \esp{Voy a comprar una moto.} / \esp{Me gustaría comprar una moto.}
\end{exercice}

\begin{textebox}[Mi plan de negocio]
Soy Amina, una joven \ldots\ldots\ldots de Bafoussam. Mi \ldots\ldots\ldots es producir miel natural. Con una pequeña \ldots\ldots\ldots de mi familia, compré diez colmenas. Ahora tengo muchos \ldots\ldots\ldots en la ciudad. El próximo año, quiero \ldots\ldots\ldots mi miel a Gabón. Si todo va bien, el \ldots\ldots\ldots será grande y podré emplear a cinco jóvenes del barrio.
\end{textebox}

\begin{exercice}[Texte à trous]
Complète le texte « Mi plan de negocio » ci-dessus avec les mots suivants : \esp{beneficio -- exportar -- emprendedora -- inversión -- clientes -- proyecto}.
\end{exercice}

\begin{exercice}[Appariements et classement]
\textbf{A.} Relie chaque mot espagnol à sa traduction française.

\begin{center}
\begin{tabularx}{\linewidth}{|X|X|}
\hline
\rowcolor{popblueL} \textbf{Español} & \textbf{Français} \\
\hline
1. el emprendedor & a. le bénéfice \\
\hline
2. el mercado & b. exporter \\
\hline
3. exportar & c. l'entrepreneur \\
\hline
4. el beneficio & d. le marché \\
\hline
5. el préstamo & e. le prêt \\
\hline
\end{tabularx}
\end{center}

\textbf{B.} Classe les formes verbales suivantes dans deux colonnes, « Futur simple » ou « Conditionnel », puis traduis chaque forme en français : \esp{venderé -- compraría -- iremos -- haríamos -- abrirán -- podría -- seré -- tendrían -- crearé -- gustaría -- importaremos -- querría -- diré -- pondría -- habrá}.
\end{exercice}

\courssub{Je me prépare au Bac}

\begin{textebox}[El cacao camerunés conquista España]
En la región del Centro de Camerún, un grupo de cuarenta agricultores ha creado una cooperativa de cacao. Antes, cada agricultor vendía su cosecha solo, a un precio muy bajo. Ahora, juntos, producen más y negocian mejores precios.

El año pasado, la cooperativa exportó sus primeras toneladas de cacao a España. « Fue un momento histórico para nosotros », explica Marta Etoundi, la presidenta. « El próximo año exportaremos el doble y contrataremos a diez jóvenes más. »

Para lograrlo, la cooperativa recibió un préstamo del banco y una formación sobre la calidad del cacao. Los agricultores aprendieron a secar bien los granos y a guardarlos en buenas condiciones. Gracias a este esfuerzo, su cacao tiene ahora una reputación excelente en el mercado europeo.

Marta tiene muchos sueños: « Me gustaría visitar España para conocer a nuestros clientes. También crearíamos una pequeña fábrica de chocolate aquí, en Yaundé. Así, el beneficio quedaría en nuestro país y daría trabajo a nuestra gente. »

El proyecto de la cooperativa demuestra que, con organización y esfuerzo, los pequeños productores pueden participar en los intercambios comerciales internacionales.
\end{textebox}

\begin{exercice}[Type Bac — Compréhension écrite]
Lis le texte « El cacao camerunés conquista España » ci-dessus, puis réponds aux questions.

\textbf{Questions.}
\begin{enumerate}
\item Vrai ou faux ? Justifie avec une phrase du texte.
\begin{enumerate}
\item Los agricultores de la cooperativa trabajan solos.
\item La cooperativa ya exportó cacao a España.
\item Marta Etoundi quiere construir una fábrica en España.
\end{enumerate}
\item Réponds en espagnol avec des phrases complètes.
\begin{enumerate}
\item ¿Qué recibió la cooperativa para mejorar su producción?
\item ¿Cuáles son los dos sueños de Marta?
\end{enumerate}
\item Trouve dans le texte l'équivalent des mots suivants : \emph{la récolte ; le prêt ; le bénéfice ; les échanges commerciaux}.
\item Relève dans le texte deux verbes au futur simple et deux verbes au conditionnel, puis traduis-les.
\end{enumerate}
\end{exercice}

\begin{exercice}[Type Bac — Grammaire, thème et version]
\textbf{A. Grammaire.} Complète les phrases en conjuguant le verbe entre parenthèses au futur simple (projet certain) ou au conditionnel (souhait, politesse).
\begin{enumerate}
\item Mañana nosotros \ldots\ldots\ldots (visitar) el mercado de Duala.
\item ¿\ldots\ldots\ldots (poder, usted) darme un préstamo, por favor?
\item El mes próximo, ellos \ldots\ldots\ldots (importar) arroz.
\item A mí \ldots\ldots\ldots (gustar) crear una empresa de zumos.
\item Yo \ldots\ldots\ldots (ser) empresaria dentro de cinco años.
\item Nosotros \ldots\ldots\ldots (querer) exportar nuestro café, pero es difícil.
\item ¿Cuánto \ldots\ldots\ldots (costar) este proyecto? — Un millón de francos.
\item Mi hermano \ldots\ldots\ldots (abrir) una tienda el año próximo.
\item ¿\ldots\ldots\ldots (venir, tú) a la reunión de la cooperativa mañana?
\item Con más dinero, \ldots\ldots\ldots (contratar, yo) a cinco empleados.
\end{enumerate}

\textbf{B. Thème.} Traduis en espagnol.
\begin{enumerate}
\item L'année prochaine, nous exporterons notre café et j'aimerais trouver des clients au Mexique.
\item Je vais créer une petite entreprise avec mes amis.
\item Voudriez-vous investir dans mon projet ?
\item Les entrepreneurs camerounais vendent leurs produits au marché.
\item Avec ce prêt, elle achètera dix machines et embauchera des jeunes.
\end{enumerate}

\textbf{C. Version.} Traduis en français le texte « Comercio entre vecinos » ci-dessous.
\end{exercice}

\begin{textebox}[Comercio entre vecinos]
Guinea Ecuatorial y Camerún son países vecinos y mantienen intercambios comerciales importantes. Los comerciantes cameruneses venden frutas, verduras y cacao en los mercados de Malabo y Bata. A cambio, Camerún importa pescado y otros productos. Muchos jóvenes emprendedores de los dos países sueñan con crear empresas juntos. « Me gustaría abrir una tienda en Duala », dice María, una joven de Bata. El próximo año, los dos gobiernos firmarán nuevos acuerdos para facilitar este comercio.
\end{textebox}

\begin{exercice}[Type Bac — Expression écrite : Presenta tu proyecto empresarial]
\textbf{Sujet.} Tu participes à un concours de jeunes entrepreneurs organisé par ton lycée. Rédige en espagnol la présentation de ton projet d'entreprise (entre 90 et 110 mots). Tu indiqueras :
\begin{itemize}
\item ton nom, ta ville et l'idée de ton projet ;
\item ce dont tu as besoin pour commencer (\esp{inversión}, \esp{préstamo}, local\ldots) ;
\item tes projets certains pour l'avenir (futur simple ou \esp{voy a + infinitivo}) ;
\item tes rêves et souhaits (conditionnel : \esp{me gustaría\ldots}).
\end{itemize}

\textbf{Grille d'évaluation.}

\begin{center}
\begin{tabularx}{\linewidth}{|X|l|}
\hline
\rowcolor{popblueL} \textbf{Critère} & \textbf{Points} \\
\hline
Lexique de l'entreprise et du commerce (au moins 6 mots du champ lexical) & 4 pts \\
\hline
Emploi correct du futur simple et du conditionnel (au moins 4 verbes) & 6 pts \\
\hline
Cohérence, organisation et respect du nombre de mots & 4 pts \\
\hline
\end{tabularx}
\end{center}
\end{exercice}



\newpage

\coursec{Corrigés des exercices}

\courssub{Je vérifie mes connaissances}

\begin{aretenir}[Rappel : Les temps du projet et du souhait]
Pour parler d'un projet entrepreneurial, on utilise trois structures principales :
\begin{enumerate}
\item \textbf{Futur proche (\esp{ir a + infinitif})} : projet déjà décidé, action imminente. \esp{Voy a crear mi empresa.}
\item \textbf{Futur simple} : projet certain, prévision, promesse, souvent avec indicateur de temps (\esp{el año próximo, mañana}). \esp{Crearé una empresa.}
\item \textbf{Conditionnel (\esp{me gustaría, querría, podría})} : souhait, rêve, hypothèse, politesse. \esp{Me gustaría exportar a España.}
\end{enumerate}
\emph{Attention :} les radicaux irréguliers du futur simple (\esp{dir-, har-, podr-, querr-, tendr-, pondr-, habr-, sabr-, saldr-, vendr-}) sont \textbf{identiques} au conditionnel.
\end{aretenir}

\begin{corrige}[Exercice 1 — QCM : Le vocabulaire de l'entreprise]
\begin{enumerate}
\item \textbf{un emprendedor} : c'est la personne qui crée une entreprise (le \esp{cliente} est celui qui achète, le \esp{mercado} est le lieu ou l'espace d'échange).
\item \textbf{vender productos al extranjero} : \esp{exportar} = vendre à l'étranger ; « acheter à l'étranger » se dit \esp{importar}.
\item \textbf{el beneficio} : c'est le gain, le bénéfice ; \esp{la inversión} est l'argent investi, \esp{el préstamo} est le prêt.
\item \textbf{un deseo} : le conditionnel \esp{me gustaría} exprime un souhait (« j'aimerais »), pas un fait réel ni une certitude.
\end{enumerate}
\end{corrige}

\begin{corrige}[Exercice 2 — Vrai ou faux ? Justifie ta réponse]
\begin{enumerate}
\item \textbf{Vrai} : la structure \esp{ir a + infinitivo} (\esp{voy a abrir}) exprime un projet proche, « je vais ouvrir ».
\item \textbf{Faux} : \esp{crearé} est la première personne du singulier du \emph{futur simple} (« je créerai ») ; le conditionnel serait \esp{crearía} (« je créerais »).
\item \textbf{Vrai} : le cacao est l'un des principaux produits que le Cameroun exporte, notamment vers l'Europe et l'Espagne.
\item \textbf{Vrai} : \esp{me gustaría} est le conditionnel de \esp{gustar} et se traduit par « j'aimerais / je voudrais ».
\item \textbf{Faux} : \esp{importar} signifie « acheter à l'étranger » ; « vendre à l'étranger » se dit \esp{exportar}.
\end{enumerate}
\end{corrige}

\begin{corrige}[Exercice 3 — Trouve le mot]
\begin{enumerate}
\item \esp{Mi \textbf{proyecto} es vender jugo natural en el barrio.} « Mon projet est de vendre du jus naturel dans le quartier. »
\item \esp{Para empezar, necesito una \textbf{inversión} de 100\,000 francos.} « Pour commencer, j'ai besoin d'un investissement de 100\,000 francs. »
\item \esp{Cada \textbf{cliente} compra dos botellas por día.} « Chaque client achète deux bouteilles par jour. »
\item \esp{Los sábados, vendo mis productos en el \textbf{mercado} de Mokolo.} « Le samedi, je vends mes produits au marché de Mokolo. »
\item \esp{Quiero crear una pequeña \textbf{empresa} con mis hermanos.} « Je veux créer une petite entreprise avec mes frères et sœurs. »
\end{enumerate}
\end{corrige}

\begin{aretenir}[Formation du Futur Simple et du Conditionnel]
\begin{center}
\begin{tabularx}{\linewidth}{|X|X|}
\hline
\rowcolor{popblueL} \textbf{Futur Simple} & \textbf{Conditionnel} \\
\hline
Radical (souvent l'infinitif) + \esp{-é, -ás, -á, -emos, -éis, -án} & Même radical + \esp{-ía, -ías, -ía, -íamos, -íais, -ían} \\
\hline
\esp{crear-é, crear-ás, crear-á...} & \esp{crear-ía, crear-ías, crear-ía...} \\
\hline
\multicolumn{2}{|l|}{\textbf{Radicaux irréguliers communs aux deux temps :}} \\
\hline
\esp{dir-} (\esp{decir}) & \esp{har-} (\esp{hacer}) \\
\esp{podr-} (\esp{poder}) & \esp{querr-} (\esp{querer}) \\
\esp{tendr-} (\esp{tener}) & \esp{pondr-} (\esp{poner}) \\
\esp{habr-} (\esp{hay}) & \esp{sabr-} (\esp{saber}) \\
\esp{saldr-} (\esp{salir}) & \esp{vendr-} (\esp{venir}) \\
\hline
\end{tabularx}
\end{center}
\end{aretenir}

\begin{corrige}[Exercice 4 — Conjugaison à compléter]
\begin{enumerate}
\item \esp{El año próximo, yo \textbf{crearé} mi empresa.} Futur simple : projet certain, annoncé avec une date précise.
\item \esp{¿Le \textbf{gustaría} a usted invertir en mi proyecto?} Conditionnel de politesse : « Aimeriez-vous investir dans mon projet ? » (\esp{Le} = leísmo courtois pour \esp{usted}).
\item \esp{Nosotros \textbf{exportaremos} cacao a España el próximo año.} Futur simple : « Nous exporterons ».
\item \esp{Ella \textbf{querría} abrir una tienda en Duala, pero no tiene dinero.} Conditionnel de \esp{querer} (radical irrégulier \esp{querr-}) : souhait contrarié.
\item \esp{Yo \textbf{podría} ayudarte, pero no tengo tiempo.} Conditionnel de \esp{poder} (radical \esp{podr-}) : « Je pourrais t'aider ».
\end{enumerate}
\end{corrige}

\begin{popfigure}
\begin{tikzpicture}[
  mindmap,
  concept color=popblue,
  level 1 concept/.append style={level distance=3.5cm, sibling angle=90, concept color=popblue!70},
  level 2 concept/.append style={level distance=2.8cm, sibling angle=45, concept color=poporange!70},
  text=white,
  font=\small\bfseries,
  every node/.style={align=center}
]
\node[concept, align=center]{Empresa\\y Comercio}
  child[concept color=popgreen] { node[concept] {Personas}
    child { node[concept] {emprendedor/a} }
    child { node[concept] {cliente} }
    child { node[concept] {inversor} }
  }
  child[concept color=poppurple] { node[concept] {Lugares}
    child { node[concept] {mercado} }
    child { node[concept] {empresa} }
    child { node[concept] {tienda} }
  }
  child[concept color=poppink] { node[concept] {Acciones}
    child { node[concept] {exportar} }
    child { node[concept] {importar} }
    child { node[concept] {invertir} }
    child { node[concept] {crear} }
  }
  child[concept color=popgold] { node[concept] {Resultados}
    child { node[concept] {beneficio} }
    child { node[concept] {préstamo} }
    child { node[concept] {inversión} }
  };
\end{tikzpicture}
\legende{Carte mentale : Lexique de l'entrepreneuriat et du commerce (Módulo 4).}
\end{popfigure}

\courssub{Je m'entraîne}

\begin{corrige}[Exercice 5 — Traduis en espagnol]
\begin{enumerate}
\item \esp{Voy a abrir un restaurante en Yaundé.} Futur proche : \esp{ir a + infinitivo} ; « à Yaoundé » se traduit par \esp{en Yaundé} (adaptation phonétique espagnole).
\item \esp{Me gustaría exportar mi café a España.} Conditionnel du souhait ; attention : « vers l'Espagne » se dit simplement \esp{a España} (verbe de mouvement/direction).
\item \esp{Crearemos una empresa el año próximo.} Futur simple de \esp{crear}, première personne du pluriel.
\item \esp{Los clientes compran mucho cacao en el mercado.} « Au marché » = \esp{en el mercado} (lieu, pas \esp{a el mercado}).
\item \esp{Ella querría encontrar inversores en México.} Conditionnel de \esp{querer} : \esp{querría} ; pas d'article devant \esp{inversores} (« des investisseurs » indéfini).
\end{enumerate}
\end{corrige}

\begin{corrige}[Exercice 6 — Transforme les phrases (Futur proche $\leftrightarrow$ Conditionnel)]
\begin{enumerate}
\item \esp{Voy a crear una empresa de transporte.} / \esp{Me gustaría crear una empresa de transporte.}
\item \esp{Vamos a exportar frutas a Guinea Ecuatorial.} / \esp{Nos gustaría exportar frutas a Guinea Ecuatorial.} \emph{Rappel :} avec \esp{nosotros}, le pronom IO est \esp{nos} ; le verbe \esp{gustar} reste au singulier car le sujet est l'infinitif \esp{exportar}.
\item \esp{Voy a abrir una tienda en el mercado central.} / \esp{Me gustaría abrir una tienda en el mercado central.}
\item \esp{Voy a buscar clientes en España.} / \esp{Me gustaría buscar clientes en España.}
\end{enumerate}
\end{corrige}

\begin{corrige}[Exercice 7 — Texte à trous]
\esp{Soy Amina, una joven \textbf{emprendedora} de Bafoussam. Mi \textbf{proyecto} es producir miel natural. Con una pequeña \textbf{inversión} de mi familia, compré diez colmenas. Ahora tengo muchos \textbf{clientes} en la ciudad. El próximo año, quiero \textbf{exportar} mi miel a Gabón. Si todo va bien, el \textbf{beneficio} será grande y podré emplear a cinco jóvenes del barrio.}

Traduction : « Je suis Amina, une jeune entrepreneure de Bafoussam. Mon projet est de produire du miel naturel
\end{corrige}


\newpage

\coursec{Fiche bilan}

\begin{popfigure}
\begin{tikzpicture}[mindmap, grow cyclic, every node/.style={concept, align=center, font=\small},
root concept/.append style={concept color=popblue, font=\bfseries},
level 1/.append style={level distance=3.6cm, sibling angle=51.5},
level 2/.append style={level distance=2.2cm, font=\footnotesize}]
\node[root concept]{El emprendimiento y los intercambios comerciales}
child[concept color=poporange]{node{Lexique de l'entreprise}
child{node{empresa, proyecto, mercado}}
child{node{cliente, beneficio, inversión}}}
child[concept color=popgreen]{node{Lexique du commerce}
child{node{exportar / importar}}
child{node{cacao, café, petróleo}}}
child[concept color=poppurple]{node{Futur proche}
child{node{ir a + infinitivo}}
child{node{voy a crear}}}
child[concept color=poppink]{node{Futur simple}
child{node{infinitivo + é, ás, á...}}
child{node{crearé, será, tendré}}}
child[concept color=popgold]{node{Conditionnel du souhait}
child{node{me gustaría + inf.}}
child{node{podría, querría}}}
child[concept color=popblue!60]{node{Cameroun et monde hispanophone}
child{node{Guinée équatoriale}}
child{node{Espagne, Mexique}}}
child[concept color=popgreen!70]{node{Méthode}
child{node{commentaire de texte}}
child{node{présenter un projet}}};
\end{tikzpicture}
\legende{Carte mentale de la leçon : entrepreneuriat et échanges commerciaux}
\end{popfigure}

\begin{aretenir}[L'essentiel de la leçon]
\textbf{1. Lexique clé à connaître par cœur}

\begin{center}
\begin{tabularx}{\linewidth}{|X|X|X|X|}
\hline
\rowcolor{popblueL} \textbf{Español} & \textbf{Français} & \textbf{Español} & \textbf{Français} \\
\hline
el/la emprendedor(a) & l'entrepreneur / l'entrepreneure & el mercado & le marché \\
\hline
la empresa & l'entreprise & exportar & exporter \\
\hline
el proyecto & le projet & importar & importer \\
\hline
el cliente / la clienta & le client / la cliente & el beneficio & le bénéfice, le profit \\
\hline
la inversión & l'investissement & los intercambios comerciales & les échanges commerciaux \\
\hline
crear una empresa & créer une entreprise & vender / comprar & vendre / acheter \\
\hline
\end{tabularx}
\end{center}

\textbf{2. Grammaire à connaître par cœur}

\begin{center}
\begin{tabularx}{\linewidth}{|X|X|X|}
\hline
\rowcolor{popgreenL} \textbf{Structure} & \textbf{Formation} & \textbf{Exemple} \\
\hline
Futur proche & \esp{ir} (présent) + \esp{a} + infinitif & \esp{Voy a crear mi empresa.} « Je vais créer mon entreprise. » \\
\hline
Futur simple & infinitif + \esp{-é, -ás, -á, -emos, -éis, -án} & \esp{Crearé una tienda en Duala.} « Je créerai une boutique à Douala. » \\
\hline
Futurs irréguliers & \esp{tener $\to$ tendré} ; \esp{hacer $\to$ haré} ; \esp{poder $\to$ podré} ; \esp{ser $\to$ será} & \esp{El proyecto será un éxito.} « Le projet sera un succès. » \\
\hline
Conditionnel du souhait & infinitif + \esp{-ía, -ías, -ía, -íamos, -íais, -ían} & \esp{Me gustaría exportar cacao.} « J'aimerais exporter du cacao. » \\
\hline
\end{tabularx}
\end{center}

Rappel : le verbe \esp{ir} au présent se conjugue \esp{voy, vas, va, vamos, vais, van}. Le futur proche exprime un projet concret et proche ; le futur simple exprime une prédiction, une promesse ou un projet plus lointain ; le conditionnel exprime le souhait, le rêve ou la politesse.

\textbf{3. Méthodes express}
\begin{itemize}
\item \textbf{Lire un texte économique} : repérer le projet (qui ? quoi ? où ?), souligner le lexique du commerce, relever les temps utilisés pour le projet (futur proche, futur simple, conditionnel).
\item \textbf{Présenter un projet} en trois temps : \esp{Primero, voy a...} (projet proche) ; \esp{Después, crearé...} (projet lointain) ; \esp{Me gustaría...} (souhait, rêve).
\item \textbf{Parler des échanges commerciaux} : nommer les produits (\esp{el cacao, el café, el petróleo}), les partenaires (\esp{España, Guinea Ecuatorial, México}) et utiliser \esp{exportar / importar}.
\end{itemize}

\textbf{4. Pièges à éviter}
\begin{itemize}
\item \esp{Beneficio} = « bénéfice » (commerce) ; ne pas confondre avec \esp{beneficioso} = « bénéfique ».
\item \esp{Empresa} = « entreprise » ; « entreprendre » se dit \esp{emprender} (faux ami partiel).
\item Ne jamais dire \esp{me gustaría de crear} : après \esp{me gustaría}, l'infinitif directement, sans préposition.
\item Attention aux accents obligatoires du futur et du conditionnel : \esp{crearé, será, tendré, gustaría, podría}. Sans accent, le mot change de sens ou de temps (\esp{creare} n'existe pas ; \esp{gustaria} est fautif).
\item « À Douala » se dit \esp{en Duala} : la préposition de lieu est \esp{en}, jamais \esp{a} seule devant une ville.
\end{itemize}
\end{aretenir}

\begin{aretenir}[Checklist avant l'examen]
\begin{itemize}
\item[$\square$] Je connais le lexique : \esp{emprendedor, empresa, proyecto, mercado, cliente, beneficio, inversión}.
\item[$\square$] Je sais conjuguer \esp{ir} au présent (\esp{voy, vas, va, vamos, vais, van}) et former le futur proche : \esp{voy a exportar}.
\item[$\square$] Je connais les terminaisons du futur simple (-é, -ás, -á, -emos, -éis, -án).
\item[$\square$] Je connais les futurs irréguliers : \esp{tendré, haré, podré, será, diré, saldré}.
\item[$\square$] Je sais exprimer un souhait avec \esp{me gustaría + infinitivo}, sans préposition.
\item[$\square$] Je sais distinguer futur proche (projet concret) et futur simple (prédiction, promesse).
\item[$\square$] Je sais nommer les produits d'exportation du Cameroun : \esp{cacao, café, petróleo, madera}.
\item[$\square$] Je sais citer les partenaires hispanophones du Cameroun : \esp{Guinea Ecuatorial, España, México}.
\item[$\square$] Je sais présenter un projet entrepreneurial en trois étapes (proche, lointain, souhait).
\item[$\square$] Je sais construire un commentaire de texte : introduction, repérage du lexique, analyse des temps, conclusion.
\item[$\square$] Je n'oublie pas les accents : \esp{crearé, gustaría, inversión, éxito}.
\item[$\square$] Je sais traduire « exporter / importer » sans calquer le français : \esp{exportar / importar}.
\end{itemize}
\end{aretenir}

\findecours

\end{document}
